 % 本文件由 scripts/assemble.py 自动装配，勿手改（改 parts/ 后重跑）
% 第9章 自旋液体与量子序的平均场理论

% 第9章 MEAN-FIELD THEORY OF SPIN LIQUIDS AND QUANTUM ORDER

\chapter{自旋液体与量子序的平均场理论}

\begin{keypoints}
\item 用从属玻色子方法（或投影构造）建立自旋液体态的平均场理论。
\item 规范相互作用的重要性，以及如何应用（含规范相互作用的）平均场理论来研究自旋液体的物理性质。
\item 自旋液体的表征，以及拓扑序与量子序的概念。
\end{keypoints}

在本章中，我们将为自旋液体态建立一个平均场理论。这里所谓"自旋液体态"，是指具有自旋旋转对称性、且每个元胞含\emph{奇数}个电子的绝缘体。通常，每个元胞含奇数个电子的态具有半填充能带，是导体。因此，自旋液体——如果它们存在的话——是非常奇特的态。自旋液体如此奇异，以致许多人根本不相信它们能够存在。事实上，即使到了现在，我们也还不知道有哪个自旋哈密顿量能被证明可以可靠地给出自旋液体基态。

对于那些相信自旋液体存在的人，他们必须接受或相信自旋液体的下列奇特（或者说迷人）的性质。（i）自旋液体中的激发总是携带分数量子数，例如电中性的自旋 $1/2$，这从电子的任何集合出发都不可能得到；有时激发甚至可以携带分数统计。（ii）不同的自旋液体无法靠它们的对称性质来区分；自旋液体是量子有序态的例子。（iii）有能隙自旋液体的基态总是具有拓扑简并，这种简并不能与任何对称性相联系。（iv）自旋液体总是包含某种规范涨落。

在本章中，我们将假定自旋液体确实存在，并为自旋液体建立平均场理论。平均场理论使我们能够理解自旋液体的某些物理性质。在计入重要的平均场涨落之后，我们将证明，某些平均场态对这些涨落是稳定的，因而代表真实的自旋液体态。由于自旋液体是具有非平凡拓扑序/量子序的典型态，我们将以自旋液体为载体来建立量子序的理论。

\section{量子自旋液体态的投影构造}

在这一节中，我们将使用从属玻色子方法（slave-boson approach，或称投影构造（projective construction））（Baskaran et al., 1987; Affleck et al., 1988; Baskaran and Anderson, 1988; Dagotto et al., 1988; Wen and Lee, 1996; Senthil and Fisher, 2000）来构造二维自旋液体。Baskaran 和 Anderson（1988）在从属玻色子方法中发现的规范结构，对我们理解强关联自旋液体起着关键作用。除从属玻色子方法之外，还可以用另一类投影构造——从属费米子方法或 Schwinger 玻色子方法——来构造自旋液体态（Arovas and Auerbach, 1988; Read and Sachdev, 1991; Sachdev and Park, 2002）。由于从属费米子方法只能给出具有有限能隙的自旋液体，这里我们将集中于从属玻色子方法。

\subsection{自旋液体态的平均场理论}

\begin{keypoints}
\item 自旋液体的平均场理论可以由投影构造得到。
\item 平均场理论以及 $\pi$ 通量相的平均场拟设（ansatz）。
\item $a_0(\pmb{i})$ 的规范涨落实施约束条件。
\end{keypoints}

让我们考虑二维正方格子上的 Hubbard 模型（式 (5.5.1)）。在半填充时，Hubbard 模型化为如下的海森堡模型：
\begin{equation*}
H=\sum_{\langle ij\rangle}J_{ij}\,\pmb{S}_i\cdot\pmb{S}_j
\tag{9.1.1}
\end{equation*}

我们想指出，上述自旋 $1/2$ 模型也可以看作一个硬核玻色子（hard-core boson）模型，其中 $|\downarrow\rangle$ 对应空格点，$|\uparrow\rangle$ 对应被一个玻色子占据的格点。上述哈密顿量很难求解，于是我们用平均场近似来理解它的物理性质。在平均场近似中，我们把 $\pmb{S}_i$ 之一换成其量子平均值 $\langle\pmb{S}_i\rangle$，得到如下平均场哈密顿量：
\begin{equation*}
H_{\text{mean}}=\sum_{\langle ij\rangle}J_{ij}\left(\langle\pmb{S}_i\rangle\cdot\pmb{S}_j+\pmb{S}_i\cdot\langle\pmb{S}_j\rangle-\langle\pmb{S}_i\rangle\cdot\langle\pmb{S}_j\rangle\right)
\end{equation*}

平均场哈密顿量 $H_{\text{mean}}$ 容易求解，我们可以得到 $H_{\text{mean}}$ 的基态，即 $|\Phi_{\text{mean}}\rangle$。我们唯一需要做的，是仔细选取 $\langle\pmb{S}_i\rangle$ 的值，使它们满足所谓的自洽方程
\begin{equation*}
\langle\pmb{S}_i\rangle=\langle\Phi_{\text{mean}}|\pmb{S}_i|\Phi_{\text{mean}}\rangle
\end{equation*}

这就是自旋系统的标准平均场方法。

上述标准平均场方法的问题在于，它只能用于研究有序的自旋态，因为我们从一开始就假定了 $\langle\pmb{S}_i\rangle\neq0$。长期以来，自旋液体的平均场理论似乎是不可能建立的。1987 年，我们终于找到了一个奇特的技巧——从属玻色子方法\footnote{从属玻色子（slave boson）和从属费米子（slave fermion）是两个非常奇怪且令人困惑的术语。从属玻色子方法里没有玻色子，从属费米子方法里没有费米子。这就是我更愿意用"投影构造"来称呼这两种方法的原因。}——来做这件事（Baskaran et al., 1987）。为了得到自旋液体的平均场基态，我们引入自旋子（spinon）算符 $f_{i\alpha}$，$\alpha=1,2$，它们是自旋 $1/2$ 的电中性算符。自旋算符 $\pmb{S}_i$ 表示为
\begin{equation*}
\pmb{S}_i=\frac{1}{2}f^{\dagger}_{i\alpha}\sigma_{\alpha\beta}f_{i\beta}
\tag{9.1.2}
\end{equation*}

用自旋子算符表示，哈密顿量 (9.1.1) 可以改写为
\begin{equation*}
H=\sum_{\langle ij\rangle}-\frac{1}{2}J_{ij}f^{\dagger}_{i\alpha}f_{j\alpha}f^{\dagger}_{j\beta}f_{i\beta}+\sum_{\langle ij\rangle}J_{ij}\left(\frac{1}{2}n_i-\frac{1}{4}n_in_j\right)
\tag{9.1.3}
\end{equation*}

这里我们用了 $\sigma_{\alpha\beta}\cdot\sigma_{\alpha'\beta'}=2\delta_{\alpha\beta'}\delta_{\alpha'\beta}-\delta_{\alpha\beta}\delta_{\alpha'\beta'}$，而 $n_i$ 是格点 $i$ 上的费米子数。式 (9.1.3) 中的第二项是常数，在下面的讨论中将略去。注意，式 (9.1.3) 的希尔伯特空间每个格点有四个态，比每格点只有两个态的式 (9.1.1) 的希尔伯特空间要大。式 (9.1.1) 与式 (9.1.3) 之间的等价只在每个格点恰好有一个费米子的子空间中才成立。因此，要用式 (9.1.3) 来描写自旋态，我们需要施加约束（Baskaran et al., 1987; Baskaran and Anderson, 1988）
\begin{equation*}
f^{\dagger}_{i\alpha}f_{i\alpha}=1,\qquad f_{i\alpha}f_{i\beta}\epsilon_{\alpha\beta}=0
\tag{9.1.4}
\end{equation*}

第二个约束实际上是第一个约束的推论。

零阶平均场基态可以通过如下近似得到。首先，我们用约束 (9.1.4) 的基态平均值
\begin{equation*}
\langle f^{\dagger}_{i\alpha}f_{i\alpha}\rangle=1
\tag{9.1.5}
\end{equation*}

来代替约束本身。这一约束可以通过在哈密顿量中引入一个依赖格点且不依赖时间的拉格朗日乘子 $a_0(\pmb{i})(f^{\dagger}_{i\alpha}f_{i\alpha}-1)$ 来实施。其次，我们把算符 $f^{\dagger}_{i\alpha}f_{j\alpha}$ 换成其基态期望值 $\chi_{ij}$，同样忽略它们的涨落。这样，我们便得到零阶平均场哈密顿量
\begin{equation*}
H_{\text{mean}}=\sum_{\langle ij\rangle}-\frac{1}{2}J_{ij}\left[\left(f^{\dagger}_{i\alpha}f_{j\alpha}\chi_{ji}+\text{h.c}\right)-|\chi_{ij}|^{2}\right]+\sum_{i}a_0(\pmb{i})\left(f^{\dagger}_{i\alpha}f_{i\alpha}-1\right)
\tag{9.1.6}
\end{equation*}

式 (9.1.6) 中的 $\chi_{ij}$ 必须满足自洽条件
\begin{equation*}
\chi_{ij}=\langle f^{\dagger}_{i\alpha}f_{j\alpha}\rangle
\tag{9.1.7}
\end{equation*}

而依赖格点的化学势 $a_0(\pmb{i})$ 则选得使平均场基态满足式 (9.1.5)。为方便起见，我们将把 $\chi_{ij}$ 称为自旋液体态的"拟设"。

设 $\bar{\chi}_{ij}$ 是式 (9.1.7) 的一个解，$\bar{a}(\pmb{i})$ 是式 (9.1.5) 的一个解。这样的拟设对应于平均场基态。由于 $\chi_{ij}$ 和 $a_0$ 在自旋旋转变换下不变，平均场基态在自旋旋转变换下不变。如果我们忽略 $\bar{\chi}_{ij}$ 和 $\bar{a}(\pmb{i})$ 的涨落，则平均场态周围的激发由下式描写：
\begin{equation*}
H_{\text{mean}}=\sum_{\langle ij\rangle}-\frac{J_{ij}}{2}\left[\left(f^{\dagger}_{i\alpha}f_{j\alpha}\bar{\chi}_{ji}+\text{h.c}\right)-|\bar{\chi}_{ij}|^{2}\right]+\sum_{i}\bar{a}_0(\pmb{i})\left(f^{\dagger}_{i\alpha}f_{i\alpha}-1\right)
\end{equation*}

我们看到，零阶平均场理论中的激发是由 $f_{i\alpha}$ 描写的自由自旋子。自旋子是自旋 $1/2$ 的中性费米子。令人惊叹的是，费米子激发竟然能从一个纯玻色的模型——式 (9.1.1)——中演生出来。

现在的问题是：我们是否应该相信这个平均场结果？我们是否应该相信具有自旋 $1/2$ 中性费米子激发的自旋液体能够存在？检验的办法之一是把平均场拟设周围的涨落包括进来，看看涨落是否会改变平均场结果。因此，下面我们将考虑涨落的效应。

首先，我们想指出，如果当初把 $a_0$ 的涨落（即时间依赖性）包括进来，约束 (9.1.5) 就会变回原来的约束 (9.1.4)。为看清这一点，让我们考虑 $H_{\text{mean}}$ 的路径积分表述：
\begin{gather*}
Z=\int\mathcal{D}f\,\mathcal{D}[a_0(\pmb{i})]\,\mathcal{D}\chi_{ij}\,\mathrm{e}^{\mathrm{i}\int\mathrm{d}t\,(\mathcal{L}-\sum_{i}a_0(\pmb{i},t)(f^{\dagger}_if_i-1))}\\
\mathcal{L}=\sum_{i}f^{\dagger}_i\mathrm{i}\partial_tf_i-\sum_{\langle ij\rangle}-\frac{1}{2}J_{ij}\left[\left(f^{\dagger}_{i\alpha}f_{j\alpha}\chi_{ji}+\text{h.c}\right)-|\chi_{ij}|^{2}\right]
\tag{9.1.8}
\end{gather*}

我们看到，对依赖时间的 $a_0(\pmb{i},t)$ 作积分将产生一个约束
\begin{equation*}
\prod_{i,t}\delta\left(f^{\dagger}_i(t)f_i(t)-1\right)
\end{equation*}

它在每个格点 $i$、每个时刻 $t$ 都得到实施。我们还注意到，对 $\chi_{ij}(t)$ 作积分将复原原来的哈密顿量 (9.1.3)。因此，式 (9.1.8) 是自旋模型 (9.1.3) 的一个严格表示。$\chi_{ij}$ 和 $a_0(\pmb{i})$ 中的涨落描写平均场基态之上的集体激发。

这里 $\chi_{ij}$ 有两种涨落：振幅涨落和相位涨落。振幅涨落具有有限能隙，在我们的讨论中并不关键。所以这里我们只考虑平均场拟设 $\bar{\chi}_{ij}$ 周围的相位涨落 $a_{ij}$：
\begin{equation*}
\chi_{ij}=\bar{\chi}_{ij}\mathrm{e}^{-\mathrm{i}a_{ij}}
\tag{9.1.9}
\end{equation*}

把上述相位涨落和 $a_0$ 的涨落包括进来后，平均场哈密顿量变为
\begin{equation*}
H=\sum_{\langle ij\rangle}-J_{ij}\left(f^{\dagger}_{i\alpha}f_{j\alpha}\bar{\chi}_{ji}\mathrm{e}^{-\mathrm{i}a_{ji}}+\text{h.c}\right)-\sum_{i}a_0(\pmb{i})\left(f^{\dagger}_{i\alpha}f_{i\alpha}-1\right)
\tag{9.1.10}
\end{equation*}

我们想把式 (9.1.10) 称为一阶平均场哈密顿量。由式 (9.1.10) 可见，$a_0$ 和 $a_{ij}$ 所描写的涨落正是一个 $U(1)$ 格点规范场（见式 (6.4.9)）（Baskaran and Anderson, 1988; Lee and Nagaosa, 1992）。该哈密顿量在规范变换
\begin{equation*}
a_{ij}\rightarrow a_{ij}+\theta_i-\theta_j,\qquad f_i\rightarrow f_i\,\mathrm{e}^{\mathrm{i}\theta_i}.
\tag{9.1.11}
\end{equation*}

之下不变。

因此，在一阶平均场理论中，激发由耦合到 $U(1)$ 规范场的自旋子描写（而非零阶平均场理论中的自由自旋子）。

在通常的平均场理论中，我们对相互作用的哈密顿量作近似以简化问题，希尔伯特空间不因平均场近似而改变。投影构造（或从属玻色子方法）在这一点上非常不同：不仅哈密顿量改变了，希尔伯特空间也改变了。这使得由投影构造得到的平均场结果不仅定量上不正确，而且定性上也不正确。例如，平均场基态 $|\Psi_{\text{mean}}\rangle$（式 (9.1.6) 中 $H_{\text{mean}}$ 的基态）甚至不是一个合法的自旋波函数，因为有些格点上可能没有费米子或有两个费米子。然而，不要放弃。在通常的平均场理论中，我们可以通过把平均场态周围的涨落包括进来，在定量上改进近似。在上面我们看到，在投影构造中，我们可以通过把平均场态周围的规范涨落（$a_0,a_{ij}$）包括进来，在定性上改进近似并复原原来的希尔伯特空间。因此，要想从投影构造得到哪怕是定性的结果，至少把平均场态周围的规范涨落包括进来是至关重要的。换句话说，我们在自旋液体的投影构造中把一个自旋切成了两半；把这两半重新粘合起来，以得到正确的物理结果，是十分重要的。

按照零阶平均场理论，自旋液体中的低能激发是自由自旋子。由上面的讨论我们看到，这样的结果是不正确的。按照一阶平均场理论，自旋子通过规范涨落相互作用。规范相互作用可能会剧烈地改变自旋子的性质。我们稍后将回到这个问题。

\subsection*{问题 9.1.1}

证明式 (9.1.3)。

\subsection{信还是不信}

\begin{keypoints}
\item 一阶平均场理论的解禁闭相导致新一类物质态——量子有序态。
\item 向物理自旋波函数的投影，以及 $U(1)$ 规范结构的含义。
\item 作为纠缠之涨落的演生规范玻色子与费米子。
\end{keypoints}

按照一阶平均场理论，平均场基态周围的涨落由规范场和费米子场描写。别忘了，我们原来的模型只是一个相互作用的纯玻色自旋模型。一个纯玻色的模型怎么会有一个由规范场和费米子场描写的有效理论呢？这简直难以置信。让我们检查一下自己是怎么走到这一步的。我们先把玻色的自旋算符拆成两个费米子算符（自旋子算符）的乘积，然后引入一个规范场把自旋子重新粘合成玻色的自旋。从这个角度看，一阶平均场理论好像只是一个假的理论：规范玻色子和费米子似乎都是假的，到头来我们拥有的只不过是玻色自旋涨落。

然而，我们不应过快地抛弃一阶平均场理论。如果规范场处于禁闭相（见 6.4.3 节），它其实能够复现上述玻色自旋涨落的图像。在禁闭相中，自旋子之间通过线性势相互作用，永远不能在低能下作为准粒子出现；规范玻色子在禁闭相中有很大的能隙，不出现在低能谱中。唯一的低能激发是自旋子对，它们对应于玻色自旋涨落。所以一阶平均场理论也许没有用，但它并没有错。它能够给出与常识一致的图像（尽管绕了很远的弯路）。

另一方面，当规范场处于解禁闭相时，一阶平均场理论也能给出有悖常识的图像。在这种情况下，自旋子和规范玻色子将作为良定义的准粒子出现。问题是：我们相信解禁闭相的图像吗？我们相信从纯玻色模型中能够演生出规范玻色子和费米子吗？显然，上面概述的投影构造太过形式化，很难让大多数人相信如此激进的结果。

我不得不说，把自旋切成两半再粘回去这套做法让很多人望而却步。很难看出这种形式操作能带来什么新的物理洞见和结果。如果这条路径真的给出了什么新结果，许多人也会把它们归咎于这一奇怪构造的人为假象，而不是自旋液体的真实物理性质。的确，投影构造非常形式化，但它也是一件超越时代的礼物。我们现在知道，投影构造真正所做的，是产生一个弦网凝聚态。弦网凝聚（string-net condensation）给出一种新型的关联态，它具有非平凡的量子序（见第 10 章）。所以，请相信我：只要恰当地计入规范涨落的效应，投影构造给出的惊人结果是可以信赖的。

这里让我们暂且不用弦网凝聚的图像，试着理解投影构造背后的物理图像。首先，我们需要理解平均场拟设 $\chi_{ij}$ 是如何与每个格点恰好有一个费米子的物理自旋波函数相联系的。我们知道，平均场波态 $|\Psi^{(\chi_{ij})}_{\text{mean}}\rangle$（$H_{\text{mean}}$ 的基态）不是自旋系统的合法波函数，因为它未必每格点有一个费米子。为了把它联系到物理自旋波函数，我们需要把 $a_0$ 的涨落包括进来，以实施每格点一个费米子的约束。有了这一理解，我们就可以把平均场态投影到每格点一个费米子的子空间上，得到自旋系统的一个合法波函数 $\Psi_{\text{spin}}(\{\alpha_i\})$：
\begin{equation*}
\Psi^{(\chi_{ij})}_{\text{spin}}(\{\alpha_i\})=\langle 0_f|\prod_{i}f_{i\alpha_i}|\Psi^{(\chi_{ij})}_{\text{mean}}\rangle.
\tag{9.1.12}
\end{equation*}

其中 $|0_f\rangle$ 是没有 f 费米子的态，即 $f_{i\alpha}|0_f\rangle=0$。式 (9.1.12) 把平均场拟设（连同它的涨落）与物理自旋波函数联系起来，使我们能够理解平均场拟设和平均场涨落的物理意义。

例如，投影 (9.1.12) 赋予规范变换 (9.1.11) 以物理意义。由规范变换
\begin{equation*}
\tilde{\chi}_{ij}=\mathrm{e}^{\mathrm{i}\theta_i}\chi_{ij}\mathrm{e}^{-\mathrm{i}\theta_j}
\tag{9.1.13}
\end{equation*}

联系的两个平均场拟设 $\chi_{ij}$ 与 $\tilde{\chi}_{ij}$ 给出同一个投影后的自旋态
\begin{equation*}
\Psi^{(\tilde{\chi}_{ij})}_{\text{spin}}(\{\alpha_i\})=\mathrm{e}^{\mathrm{i}\sum_i\theta_i}\Psi^{(\chi_{ij})}_{\text{spin}}(\{\alpha_i\})
\tag{9.1.14}
\end{equation*}

对于 $\chi_{ij}$ 的不同取值，式 (9.1.6) 中 $H_{\text{mean}}$ 的基态对应于不同的平均场波函数 $|\Psi^{(\chi_{ij})}_{\text{mean}}\rangle$；投影之后，它们给出不同的物理自旋波函数 $\Psi^{(\tilde{\chi}_{ij})}_{\text{spin}}(\{\alpha_i\})$。因此，我们可以把 $\chi_{ij}$ 看作标记不同物理自旋态的标签。式 (9.1.14) 告诉我们，这个标签不是一对一的，而是多对一的。这一性质对我们理解 $\chi_{ij}$ 涨落的反常动力学性质十分重要。用多个标签标记同一个物理态，也使我们的理论按照 6.1.1 节中的定义成为一门规范理论。

让我们考虑多对一性质，或者说 $\chi_{ij}$ 的规范结构，如何影响它的动力学性质。如果 $\chi_{ij}$ 是物理态的一对一标签，那它就会像玻色超流体中的凝聚玻色子振幅 $\langle\phi(\pmb{x},t)\rangle$，或 SDW 态中的凝聚自旋矩 $\langle\pmb{S}_i(t)\rangle$ 一样，而 $\chi_{ij}$ 的涨落将对应于一种类似自旋波模式的玻色型模式。\footnote{更确切地说，自旋波模式对应于标量玻色子。局域序参量的涨落总是给出标量玻色子。}然而，$\chi_{ij}$ 的行为并不像局域序参量 $\langle\phi(\pmb{x},t)\rangle$ 和 $\langle\pmb{S}_i(t)\rangle$ 那样无冗余地标记物理态。作为多对一的标签，$\chi_{ij}$ 的某些涨落不改变物理态，因而是非物理的；这些涨落被称为纯规范涨落（pure gauge fluctuations）。对一个一般的涨落 $\delta\chi_{ij}$，其中一部分是物理的，另一部分是非物理的。$\chi_{ij}$ 的有效理论必须是规范不变的。这一点剧烈地改变了涨落的动力学性质。正是这一性质使 $\chi_{ij}$ 的涨落表现得像规范玻色子，与声模式和自旋波模式大不相同。

我们已经论证过，$\chi_{ij}$ 的相位涨落对应于规范涨落。投影构造 (9.1.12) 使我们能够得到与规范涨落 $a_{ij}$ 相应的物理自旋波函数：
\begin{equation*}
\Psi^{(a_{ij})}_{\text{spin}}=\langle 0|\prod_{i}f_{i\alpha_i}|\Psi^{(\bar{\chi}_{ij}\mathrm{e}^{\mathrm{i}a_{ij}})}_{\text{mean}}\rangle.
\end{equation*}

类似地，投影构造也使我们能够得到与一对自旋子激发相应的物理自旋波函数。我们从带有一对粒子–空穴激发的平均场基态出发，经过投影 (9.1.12)，便得到含有一对自旋子的物理自旋波函数：
\begin{equation*}
\Psi^{\text{spinon}}_{\text{spin}}(\pmb{i}_1,\alpha_1;\pmb{i}_2,\alpha_2)=\langle 0|\Big(\prod_{i}f_{i\alpha_i}\Big)f^{\dagger}_{\pmb{i}_1\lambda_1}f_{\pmb{i}_2\lambda_2}|\Psi^{(\bar{\chi}_{ij})}_{\text{mean}}\rangle.
\end{equation*}

如果你对这里给出的物理图像不满意，仍然不相信投影构造能够产生演生的规范玻色子和费米子，那么你可以直接跳到第 10 章。在第 10 章中，我们构造了几个能够用投影构造精确（或准精确）求解的自旋模型。这些模型具有演生的 $Z_2/U(1)$ 规范结构和费米子。这些精确可解模型揭示了演生规范玻色子和费米子的弦网起源。在本章其余部分，我们假定投影构造是有意义的，并研究它的推论。

\begin{figure}[H]
\centering
\includegraphics[width=0.72\textwidth]{fig_9.1.pdf}
\caption*{图 9.1\quad 二聚体态由自旋单态对构成。把一个二聚体变成三重态便产生一个自旋 1 激发。两个分离开的自旋 $1/2$ 激发由一条由错位二聚体构成的弦相连。}
\end{figure}

\subsection{二聚体态}

\begin{keypoints}
\item 二聚体态的一阶平均场理论具有禁闭型的 $U(1)$ 规范相互作用。物理激发中只出现自旋子对的束缚态。
\end{keypoints}

让我们先讨论一个具体的平均场基态——具有最近邻耦合的海森堡模型中的二聚体态（dimer state）（Majumdar and Ghosh, 1969; Affleck et al., 1987; Read and Sachdev, 1989）：
\begin{equation*}
H=J_1\sum_{i}\left(\pmb{S}_i\pmb{S}_{\pmb{i}+\pmb{x}}+\pmb{S}_i\pmb{S}_{\pmb{i}+\pmb{y}}\right).
\tag{9.1.15}
\end{equation*}

二聚体态由如下拟设描写：
\begin{equation*}
\begin{array}{l}
\chi_{i,i+\pmb{x}}=\dfrac{1}{2}\left(1+(-)^{i_x}\right)\\[6pt]
\text{其余}\;=0
\end{array}
\tag{9.1.16}
\end{equation*}

以及 $a_0(\pmb{i})=0$。可以验证，式 (9.1.6) 中 $H_{\text{mean}}$ 的平均场基态满足自洽方程 (9.1.7) 和约束 (9.1.5)。在平均场二聚体态中，自旋子能谱没有色散：
\begin{equation*}
E_k=\pm\frac{1}{2}J_1
\tag{9.1.17}
\end{equation*}

具有 $E_k=-\frac{1}{2}J_1$ 的价带被自旋子完全填满。自旋激发在二聚体态中有有限能隙。同样清楚的是，二聚体态破坏 $x$ 方向的平移对称性。

在零阶平均场理论中，二聚体基态之上的激发是自旋 $1/2$ 的自旋子。这一结果显然不对，因为从物理上看，二聚体态是由填满格子的自旋单态二聚体构成的（见图 9.1），基本激发对应于自旋三重态的二聚体，即玻色型的自旋 1 激发。在一阶平均场理论中，自旋子耦合到 $U(1)$ 规范场。在 $2+1$ 维，$U(1)$ 规范理论是禁闭的，因此单个自旋子不可观测，物理谱中只能出现玻色型的束缚态（携带整数自旋）。如果我们真的制造出两个分离开的自旋 $1/2$ 激发，那么由图 9.1 可见，它们将由一条错位二聚体构成的弦连接起来。这根弦在两个自旋 $1/2$ 激发之间导致线性禁闭相互作用，与一阶平均场理论给出的图像一致。我们看到，一阶平均场理论给出了定性上正确的结果。

\begin{figure}[H]
\centering
\includegraphics[width=0.72\textwidth]{fig_9.2.pdf}
\caption*{图 9.2\quad (a) $\pi$ 通量态的平均场拟设。沿箭头方向 $\chi_{ij}=\mathrm{i}\chi_1$。(b) $\pi$ 通量态中费米子的色散。价带是填满的。低能激发存在于价带与导带相切的四个费米点附近。}
\end{figure}

\subsection{$\pi$ 通量态}

\begin{keypoints}
\item 投影之后，平均场 $\pi$ 通量相给出平移、旋转和宇称对称的自旋液体波函数。
\item $\pi$ 通量相的低能有效理论（一阶平均场理论）包含无能隙狄拉克费米子耦合到 $U(1)$ 规范场。这些狄拉克费米子携带自旋 $1/2$ 且电中性。
\item 稳定、边缘和不稳定平均场态的概念。
\item 一阶平均场理论只对稳定平均场态或涨落微弱的边缘平均场态可靠。$\pi$ 通量相不是这些平均场态之一。
\end{keypoints}

我们要研究的第二个平均场态，是最近邻海森堡模型的 $\pi$ 通量态（$\pi$-flux state）（Affleck and Marston, 1988; Kotliar, 1988）。$\pi$ 通量态由如下拟设给出（见图 9.2）：
\begin{equation*}
\chi_{i,i+\pmb{x}}=\mathrm{i}\chi_1,\qquad \chi_{i,i+\pmb{y}}=\mathrm{i}\chi_1(-)^{i_x},\qquad a_0(\pmb{i})=0.
\tag{9.1.18}
\end{equation*}

平均场哈密顿量 (9.1.6) 等价于每个方格带 $\pi$ 通量的电子跳跃问题，因为绕一个方格一周有 $\prod\chi_{ij}\propto\mathrm{e}^{\mathrm{i}\pi}$。$\pi$ 通量相中自旋子能谱为（见图 9.2）
\begin{equation*}
E_k=\pm J_1\chi_1\sqrt{\sin(k_x)^{2}+\sin(k_y)^{2}}
\tag{9.1.19}
\end{equation*}

其中 $-\frac{\pi}{2}<k_x<\frac{\pi}{2}$，$-\pi<k_y<\pi$。$\chi_1$ 的值可以由对平均场能量
\begin{equation*}
-\sum_{\pmb{k}}^{\prime}J_1\chi_1\sqrt{\sin(k_x)^{2}+\sin(k_y)^{2}}+J_1|\chi_1|^{2}N_{\text{site}}
\end{equation*}

求极小得到。我们发现
\begin{equation*}
\chi_1=\frac{1}{4}\int\frac{\mathrm{d}^{2}\pmb{k}}{(2\pi)^{2}}\sqrt{\sin(k_x)^{2}+\sin(k_y)^{2}}
\end{equation*}

除 $\pi$ 通量相之外，还有许多其他不同的平均场拟设局部地使平均场能量取极小。不过，对于只有最近邻耦合 $J_1$ 的海森堡模型，$\pi$ 通量相在平移不变的自旋液体中平均场能量最低。

让我们讨论 $\pi$ 通量相的一些物理性质。首先，考虑平均场态的对称性。我们知道，平均场理论旨在描写一个自旋态，其波函数为 $\Psi_{\text{spin}}(\{\alpha_i\})$，其中 $\alpha_i=\pm 1/2$ 是格点 $i$ 上 $S^{z}$ 的取值。当我们说"平均场态的对称性"时，我们真正指的是"相应自旋波函数 $\Psi_{\text{spin}}(\{\alpha_i\})$ 的对称性"。相应的自旋波函数由平均场态投影到每格点一个费米子的子空间得到（见式 (9.1.12)）。

由于平均场拟设 $\chi_{ij}$ 不依赖自旋取向，平均场基态中每一个负能级都被一个自旋向上和一个自旋向下的费米子占据。因此，平均场态 $|\Psi^{(\chi_{ij})}_{\text{mean}}\rangle$ 在自旋旋转下不变；结果，投影得到的物理自旋波函数也在自旋旋转下不变。

然而，平均场拟设在沿 $x$ 方向平移一个晶格间距下并非不变。因此，平均场态 $|\Psi^{(\chi_{ij})}_{\text{mean}}\rangle$ 破坏平移对称性。这似乎暗示物理自旋波函数也破坏平移对称性。事实上，物理自旋波函数并不破坏平移对称性。原因在于，平移后的拟设 $(\chi_{i,i+\pmb{x}},\chi_{i,i+\pmb{y}})=(-\mathrm{i}\chi_1(-)^{i_x},\mathrm{i}\chi_1)$ 通过取 $\mathrm{e}^{\mathrm{i}\theta_i}=(-)^{i_x}$ 的规范变换 (9.1.11) 与原拟设 $(\chi_{i,i+\pmb{x}},\chi_{i,i+\pmb{y}})=(\mathrm{i}\chi_1(-)^{i_x},\mathrm{i}\chi_1)$ 相联系。因此，这两个拟设给出同一个物理自旋波函数，投影后的自旋波函数在平移下不变。$\pi$ 通量相是一个平移对称的自旋液体。

第二，我们考虑低能激发的性质。半填充（即平均每格点一个费米子）时的费米面是位于 $(k_x,k_y)=(0,0)$ 和 $(0,\pi)$ 的点。在平均场层次上，$\pi$ 通量态含有无能隙自旋激发，它们对应于跨费米点的粒子–空穴激发。除这些无能隙自旋激发之外，平均场通量态还含有 $U(1)$ 规范涨落 $(a_{ij},a_0)$。在一阶平均场理论中，低能激发由耦合到 $U(1)$ 规范场的无能隙自旋子描写。

如果我们忽略自旋子 $f_i$ 与规范涨落 $(a_{ij},a_0)$ 之间的相互作用，那么 $\pi$ 通量相将具有由费米型准粒子描写的无能隙中性自旋 $1/2$ 激发。然而，除非证明了在计入规范相互作用后这一结果仍然成立，我们不应相信这个惊人的结果。

现在让我们考虑自旋激发与规范涨落之间的相互作用。为方便起见，让我们在哈密顿量 (9.1.10) 中加上 Maxwell 项 $\frac{1}{8\pi g^{2}}(v^{2}f_{12}^{2}-f_{0i}^{2})$，其中 $g$ 是耦合常数，$f_{\mu\nu}$ 是 $a_{\mu}$ 的场强。原理论对应于 $g\to\infty$ 的极限。Maxwell 项是在积掉费米子的过程中产生的。设 $g(\Lambda)$ 是积掉能标 $\Lambda$ 与 $\Lambda_0$ 之间的费米子后得到的耦合常数，其中 $\Lambda_0$ 是原理论中的能量截断。由量纲分析可得 $\mathrm{d}g^{-2}(\Lambda)\sim\mathrm{d}\Lambda/\Lambda^{2}$（注意 $\mathrm{d}g^{-2}(\Lambda)\propto\mathrm{d}\Lambda$）。于是 $g^{-2}(\Lambda)\sim\Lambda^{-1}-\Lambda_0^{-1}$。

由于与规范场耦合，一个自旋子会产生规范场 $(a_0,a_{ij})$ 的"电"场 $f_{0i}$（注意自旋子携带规范场的一个单位电荷）。粒子–空穴激发中粒子–空穴对之间的势能为
\begin{equation*}
V(\pmb{r}_1-\pmb{r}_2)=2\pi g^{2}\ln|\pmb{r}_1-\pmb{r}_2|
\tag{9.1.20}
\end{equation*}

我们发现粒子与空穴在长距离上仍有相互作用。为估计这一相互作用的强度和效应，让我们比较粒子–空穴对的动能 $E_K\sim 1/|\pmb{r}_1-\pmb{r}_2|$ 与势能 $V(\pmb{r}_1-\pmb{r}_2)$。由于 $g^{2}\sim v/|\pmb{r}_1-\pmb{r}_2|$（假定 $\Lambda\sim v/|\pmb{r}_1-\pmb{r}_2|$），我们看到 $V/E_K\sim 1$。这意味着相互作用是边缘的，在低能下不能忽略。这一结果提示，$\pi$ 通量态中的量子涨落（例如规范涨落）是重要的，将剧烈改变零阶平均场理论的低能性质。在这种情况下，零阶平均场理论（具有自由自旋子激发）不能给出自旋液体低能性质的可靠图像。要得到 $\pi$ 通量态可靠的低能性质，必须处理无能隙费米子与规范场的耦合系统。忽略规范相互作用而得到的那些惊人性质可能并不成立。\footnote{当平均场涨落微弱时，边缘平均场态可以导致代数自旋液体（algebraic spin liquid）——一种在低能下没有自由准粒子的自旋液体。见 9.9.5 节和 9.10 节。}

我们看到，投影构造的怀疑者们是对的：零阶平均场理论的结果有误导性，不可轻信。然而，相信者们也是对的：一阶平均场理论确实告诉我们，零阶平均场理论的结果是不正确的。到目前为止，投影构造没有误导过我们；只要计入适当的涨落，它是可以信赖的。与自旋液体物理性质相对应的，是一阶平均场理论（而非零阶平均场理论）的结果。

为了刻画平均场态周围涨落的重要性，这里我们想引入稳定、边缘和不稳定平均场态三个概念。在稳定平均场态中，涨落是微弱的，涨落诱导的相互作用在低能下消失（即相互作用是不相关微扰）。在边缘平均场态中，相互作用与能量之比在低能极限下趋于有限常数，此时相互作用是边缘微扰。在不稳定平均场态中，相互作用与能量之比在低能极限下发散，此时相互作用是相关微扰。只有对稳定平均场态，零阶平均场理论的性质才能在涨落之下幸存。对于不稳定平均场态，涨落会在低能下驱动相变，我们无法从零阶平均场理论推出自旋液体的任何物理性质。在这种情况下，一阶平均场理论也没有用，因为它并不能帮助我们推出自旋系统的物理性质。上面讨论的 $\pi$ 通量态是一个边缘平均场态。对于边缘平均场态，如果相互作用与能量之比很小，一阶平均场理论可能有用；此时相应自旋态的物理性质可以微扰地计算。对 $\pi$ 通量态而言，相互作用与能量之比是 1 的量级，因此很难从一阶平均场理论得到 $\pi$ 通量态的低能物理性质。

由上述讨论可见，一阶平均场理论只对稳定平均场态以及涨落微弱的边缘平均场态有用。用平均场理论研究自旋液体的关键，在于找到稳定平均场态（见 9.1.6、9.2.4 和 9.2.6 节），或涨落微弱的边缘平均场态（见 9.8 节）。

\subsection*{问题 9.1.2}

$\pi$ 通量态的旋转对称性

\begin{enumerate}
\item 证明式 (9.1.14)。
\item 证明 $\pi$ 通量态不破坏 $90^{\circ}$ 旋转对称性。
\end{enumerate}

% 块边界判定：
%   起点——书页 367（p382）顶部为 Problem 9.1.3，新起条目，无 ch9_a 溢出内容，
%     故本块无 %JOIN%，也不设 %--A-TAIL-- 区（ch9_a 止于书页 366 末的问题 9.1.2）。
%   终点——书页 379（p394）止于 9.2.2 节首要点列表的前两条；第三条要点印在
%     书页 380（p395）顶部，连同其后的正文 "Just like the U(1) ..." 均属 ch9_c。
%     本块按"要点列表只收到本页末"处理：keypoints 环境在两条后闭合；
%     第三条的原文与建议译文见文件尾【块界备注】，请 ch9_c 以其文件的
%     %--B-TAIL-- 区或正文开头续排该条（环境需另行开启），勿弃勿重。

\subsection*{问题 9.1.3}

设 $|\Psi_{\text{mean}}\rangle$ 为 $\pi$ 通量拟设的平均场基态。计算 $\langle\Psi_{\text{mean}}|\pmb{S}_i\pmb{S}_{i+\pmb{y}}|\Psi_{\text{mean}}\rangle$。利用这一结果求变分平均场能量 $\langle\Psi_{\text{mean}}|H|\Psi_{\text{mean}}\rangle$，其中 $H$ 是具有近邻相互作用的海森堡模型（见式 (9.1.15)）。比较 $\pi$ 通量态与二聚体态的平均场能量。

\subsection*{问题 9.1.4}

\textbf{Dirac 费米子}

\begin{enumerate}
\item 证明：在连续极限下，$\pi$ 通量态中的低能自旋子由下式描述
\begin{equation*}
H=v\sum_{\pmb{k}\sim0}\lambda^{\dagger}_{\alpha,\pmb{k}}(k_x\Gamma_x+k_y\Gamma_y)\lambda_{\alpha,\pmb{k}}
\end{equation*}
其中 $\lambda^{\top}_{\alpha,\pmb{k}}=(f_{\alpha,\pmb{k}},f_{\alpha,\pmb{k}+\pmb{Q}})$，$\pmb{Q}=(\pi,\pi)$，且 $\Gamma_x^2=\Gamma_y^2=1$。求 $v$ 与 $\Gamma_{x,y}$。在实空间中，$H$ 可改写为 $H=\int\mathrm{d}^2x\,v\lambda^{\dagger}_{\alpha}(-\mathrm{i}\Gamma_x\partial_x-\mathrm{i}\Gamma_y\partial_y)\lambda_{\alpha}$。

\item 相应的拉格朗日量 $\mathcal{L}=\mathrm{i}\lambda^{\dagger}_{\alpha}\partial_t\lambda_{\alpha}-v\lambda^{\dagger}_{\alpha}(-\mathrm{i}\Gamma_x\partial_x-\mathrm{i}\Gamma_y\partial_y)\lambda_{\alpha}$ 可以改写为
\begin{equation*}
\mathcal{L}=\mathrm{i}\bar{\lambda}_{\alpha}(\gamma^0\partial_t+v\gamma^x\partial_x+v\gamma^y\partial_y)\lambda_{\alpha}
\end{equation*}
其中 $\bar{\lambda}_{\alpha}=\lambda^{\dagger}_{\alpha}\gamma^0$，$\gamma_0^2=1$，且 $(\Gamma_x,\Gamma_y)=(-\gamma^0\gamma^x,-\gamma^0\gamma^y)$。求 $\gamma^{0,x,y}$。由 $\mathcal{L}$ 描述的费米子称为无质量狄拉克费米子。

\item 证明 $(\gamma^0,\gamma^x,\gamma^y)$ 在 $1+2$ 维中满足下列狄拉克代数：
\begin{equation*}
\{\gamma^{\mu},\gamma^{\nu}\}=\eta^{\mu\nu},\qquad \mu,\nu=0,x,y,
\end{equation*}
其中 $\eta^{\mu\nu}$ 是对角矩阵，$\eta^{00}=1$ 且 $\eta^{xx}=\eta^{yy}=-1$。

\item 有质量的狄拉克费米子由如下拉格朗日量描述：
\begin{equation*}
\mathcal{L}=\mathrm{i}\bar{\lambda}_{\alpha}(\gamma^0\partial_t+v\gamma^x\partial_x+v\gamma^y\partial_y)\lambda_{\alpha}+m\bar{\lambda}_{\alpha}\lambda_{\alpha}.
\end{equation*}
求相应运动方程 $\mathrm{i}(\gamma^0\partial_t+v\gamma^x\partial_x+v\gamma^y\partial_y+m)\lambda_{\alpha}=0$ 的解。证明该费米子具有色散 $\omega_{\pmb{k}}=\sqrt{v^2\pmb{k}^2+m^2}$。

\item \emph{与 $U(1)$ 规范场的耦合}

用 3.7.1 节所述的程序，求最小耦合到一个 $U(1)$ 规范场的狄拉克费米子的拉格朗日量。确保所得的拉格朗日量具有 $U(1)$ 规范不变性。
\end{enumerate}

\subsection{如何消除无能隙 $U(1)$ 规范玻色子}

\begin{keypoints}
\item 如何得到稳定的平均场态：通过 Chern--Simons 项和/或 Anderson--Higgs 机制让规范涨落获得有限能隙。
\item 稳定的自旋液体总包含电中性的自旋 $1/2$ 自旋子，且自旋子之间只有短程相互作用。
\end{keypoints}

我们已经看到，在 $\pi$ 通量态中，无能隙 $U(1)$ 规范玻色子与无能隙费米子强烈相互作用，甚至直到零能。对于 $\pi$ 通量相，很难求出一阶平均场理论（即费米子--规范场耦合系统）的性质。一阶平均场理论虽然不会误导我们，但它似乎复杂得毫无用处。现在的问题是：能否构造一个稳定的平均场态，使平均场涨落在低能下相互作用微弱？如果可以，一阶平均场理论的性质就能容易地求得。

由于一阶平均场理论至少必须包含实施每格点单费米子约束的规范涨落，得到稳定平均场态的一个办法是让规范玻色子获得能隙。有能隙的规范玻色子只能传递费米子之间的短程相互作用。基于朗道费米液体理论的经验，我们知道如何处理具有短程相互作用的费米子。

首先让我们澄清"让规范玻色子获得能隙"是什么意思。我们知道，规范玻色子不过是平均场拟设 $\chi_{ij}$ 的涨落。这些涨落的动力学依赖于平均场拟设。因此，所谓"让规范玻色子获得能隙"，指的就是找到这样的平均场拟设，使拟设的集体模式涨落有能隙。

为给寻找稳定平均场拟设的探索提供动机，让我们先考虑 $U(1)$ 规范玻色子可以通过哪些方式获得能隙。事实上，在 3.7.5 节和 4.4 节中，我们已经遇到过 $U(1)$ 规范玻色子获得能隙的两种方式。第一种方式是 Anderson--Higgs 机制（Anderson--Higgs mechanism），即规范玻色子与凝聚的电荷玻色子耦合。耦合系统的动力学由
\begin{equation*}
\mathcal{L}=\frac{1}{8\pi g^2}(\mathbf{e}^2-b^2)+c_1a_0^2-c_2\pmb{a}^2
\end{equation*}
在 $2+1$ 维及连续极限下描述（见式 (3.7.17)），其中 $c_1$ 与 $c_2$ 来自凝聚的电荷玻色子。第二种方式是通过 Chern--Simons 项。如果平均场拟设使得满带具有非零的霍尔电导，那么 $U(1)$ 规范玻色子的动力学将由
\begin{equation*}
\mathcal{L}=\frac{1}{8\pi g^2}(\mathbf{e}^2-b^2)+\frac{K}{4\pi}\epsilon^{\mu\nu\lambda}a_{\mu}\partial_{\nu}a_{\lambda}.
\tag{9.1.21}
\end{equation*}
描述。在问题 9.1.5 中我们将求出，Chern--Simons 项给规范玻色子一个非零能隙。找到稳定平均场态的关键，在于找到实现上述两种机制之一的平均场拟设。

在这两种情形中，规范场都只能传递短程相互作用。其结果是，自旋子不被禁闭。自旋液体之上的准粒子由自由自旋子描述，它们携带自旋 $1/2$、电荷为零。在存在 Chern--Simons 项时，自旋子甚至可以具有分数统计。

\begin{figure}[H]
\centering
\includegraphics[width=0.72\textwidth]{fig_9.3.pdf}
\caption*{图 9.3\quad 手征自旋态的平均场拟设。$\chi_{ij}$ 在箭头方向上为 $\mathrm{i}\chi$。}
\end{figure}

\subsection*{问题 9.1.5}

证明由式 (9.1.21) 描述的涨落具有有限能隙。（可以通过计算经典运动方程来做到这一点。）

\subsection{手征自旋态}

\begin{keypoints}
\item 平均场手征自旋拟设给出一个平移和旋转对称的自旋液体态，但它破坏时间反演与宇称对称性。
\item 平均场手征自旋态是稳定的。由平均场理论得到的手征自旋液体确实代表真实自旋系统的一个稳定量子相。
\item 手征自旋液体含有分数化激发——自旋子。自旋子携带自旋 $1/2$、不带电荷，并具有分数统计。
\end{keypoints}

在这一节中，我们将讨论实现上述第二种机制的一个平均场态。该平均场态由拟设 $\chi_{ij}$ 描述，其中 $\chi_{ij}$ 是复数并产生通量。由式 (9.1.6) 描述的自旋子的行为就好像在磁场中运动。当通量与自旋子密度（即每格点一个费米子）有正确的公度关系（commensuration）时，整数个朗道能级（更确切地说，由于晶格的存在，是朗道能带）被完全填满。在这种情况下，有效格点规范理论因填满的朗道能级的有限霍尔电导而含有一个 Chern--Simons 项。我们将把这样的稳定平均场态称为手征自旋态（chiral spin state）（Wen et al., 1989; Khveshchenko and Wiegmann, 1989）。

最简单的手征自旋态由下列拟设给出（见图 9.3）：
\begin{equation*}
\begin{array}{lll}
\chi_{i,i+\pmb{x}}=\mathrm{i}\chi_1, & \chi_{i,i+\pmb{y}}=\mathrm{i}\chi_1(-)^{i_x}, & a_0(i)=0,\\[2pt]
\chi_{i,i+\pmb{x}+\pmb{y}}=-\mathrm{i}\chi_2(-)^{i_x}, & \chi_{i,i+\pmb{x}-\pmb{y}}=\mathrm{i}\chi_2(-)^{i_x} &
\end{array}
\tag{9.1.22}
\end{equation*}

上述 $\chi_{ij}$ 在每个正方形中产生 $\pi$ 通量，在每个三角形中产生 $\frac{1}{2}\pi$ 通量。上一节讨论的 $\pi$ 通量相有些特殊：由于 $\pi$ 通量等价于 $-\pi$ 通量，$\pi$ 通量相保持时间反演对称性（$T$）与宇称（$P$）。然而，由式 (9.1.22) 描述的 $\frac{1}{2}\pi$ 通量相并不等价于 $-\frac{1}{2}\pi$ 通量相。在 $T$ 或 $P$ 之下，$\frac{1}{2}\pi$ 通量变为 $-\frac{1}{2}\pi$ 通量，式 (9.1.22) 中的 $\chi_2$ 变为 $-\chi_2$。因此，具有非零 $\chi_2$ 的手征自旋态自发破坏 $T$ 与 $P$。利用恒等式
\begin{equation*}
E_{123}\equiv\pmb{S}_1\cdot(\pmb{S}_2\times\pmb{S}_3)=2\mathrm{i}(\hat{\chi}_{12}\hat{\chi}_{23}\hat{\chi}_{31}-\hat{\chi}_{13}\hat{\chi}_{32}\hat{\chi}_{21}),\qquad \hat{\chi}_{ij}=f^{\dagger}_if_j
\tag{9.1.23}
\end{equation*}
我们发现 $\langle E_{123}\rangle\sim\mathrm{Im}(\chi_{12}\chi_{23}\chi_{31})$ 在手征自旋态中不为零。注意 $E_{123}$ 在 $T$ 或 $P$ 下为奇；因此可以把 $E_{123}$ 视为 $T$ 与 $P$ 破坏的序参量。

指定了手征自旋态之后，我们想知道哪个自旋哈密顿量支持手征自旋态。让我们考虑下列阻挫自旋哈密顿量：
\begin{equation*}
H=J_1\sum_i(\pmb{S}_i\cdot\pmb{S}_{i+\pmb{x}}+\pmb{S}_i\cdot\pmb{S}_{i+\pmb{y}})+J_2\sum_i(\pmb{S}_i\cdot\pmb{S}_{i+\pmb{x}+\pmb{y}}+\pmb{S}_i\cdot\pmb{S}_{i+\pmb{x}-\pmb{y}})
\tag{9.1.24}
\end{equation*}

我们想知道平均场哈密顿量 (9.1.6) 何时支持平均场手征自旋态。在平均场哈密顿量 (9.1.6) 中，$J_{ij}$ 对近邻取 $J_1$，对次近邻取 $J_2$。

当 $\chi_2=0$ 时，由平均场哈密顿量决定的自旋子能谱由式 (9.1.19) 给出。导带与价带在 $(k_x,k_y)=(0,0)$ 与 $(0,\frac{\pi}{a})$ 两点相接触。当 $\chi_2\neq0$ 时，导带与价带之间打开一个能隙。平均场自旋子能谱由
\begin{equation*}
E^{\pm}_{\pmb{k}}=\pm\sqrt{J_1^2\chi_1^2(\sin(k_x)^2+\sin(k_y)^2)+J_2^2\chi_2^2(\cos(k_x+k_y)+\cos(k_x-k_y))^2}
\tag{9.1.25}
\end{equation*}
给出。平均场基态通过填满价带得到。由于存在能隙，自旋--自旋关联是短程的。利用基态波函数，我们可以计算 $\langle f^{\dagger}_{\alpha i}f_{\alpha i}\rangle$ 来检验自洽条件 (9.1.7)。我们发现，当 $J_2/J_1<0.49$ 时，自洽方程只支持一个 $\chi_1\neq0$、$\chi_2=0$ 的解，即 $\pi$ 通量相。当 $J_2/J_1>0.49$ 时，式 (9.1.7) 还支持一个 $\chi_2\neq0$ 的第二解。发现第二解具有更低的平均场能量。在这种情况下，$T$ 与 $P$ 被自发破坏。

注意，手征态的平均哈密顿量等价于磁场中电子跳跃的问题。从属费米子与规范场 $a_{\mu}$ 的耦合，与电子和电磁场的耦合完全相同。因此可以预期，由式 (9.1.10) 描述的从属费米子系统具有一种类似于霍尔效应的现象。这里的"霍尔效应"意味着 $a_{\mu}$ 规范场的"电"场在横向感应出一个自旋子流：
\begin{equation*}
j_x=\sigma_{xy}e_y,\qquad e_y=\partial_ta_y-\partial_ya_0
\tag{9.1.26}
\end{equation*}
其中 $\sigma_{xy}$ 是霍尔电导。有一个定理指出，满带的霍尔电导总是量子化为整数乘 $1/2\pi$（Thouless et al., 1982; Avron et al., 1983）。在我们的情形，这意味着 $\sigma_{xy}=2n/2\pi$，其中因子 2 来自自旋。"霍尔"电导的数值可以用 4.4 节讨论的方法求得。我们得到
\begin{equation*}
\sigma_{xy}=2/2\pi.
\tag{9.1.27}
\end{equation*}
理解这一结果的最简单方式，是注意到每一个自旋向上的自旋子对应一个通量量子，每一个自旋向下的自旋子也对应一个通量量子。这样，自旋向上与自旋向下的自旋子都具有填充因子 $\nu=1$。关掉晶格势之后，平均场手征自旋态中的价带就变成第一朗道能级。价带中的自旋向上与自旋向下的从属费米子各对"霍尔"电导贡献 $1/2\pi$。因此，总"霍尔"电导由式 (9.1.27) 给出。

规范涨落的有效作用量可以通过把式 (9.1.10) 中的自旋子积掉而得到。利用"霍尔"电导与 Chern--Simons 项之间的关系，我们可以容易地把连续极限下的有效作用量写成
\begin{equation*}
S=\int\mathrm{d}^3x\,\frac{1}{2}\sigma_{xy}a_{\mu}\partial_{\nu}a_{\lambda}\epsilon_{\mu\nu\lambda}+\frac{1}{8\pi g^2}(\mathbf{e}^2-v^2b^2)+\dots
\tag{9.1.28}
\end{equation*}
这里式 (9.1.28) 中的 $g^2$ 是自旋子能隙的量级，而 $v$ 是典型自旋子速度、为 $1/aJ$ 的量级。利用有效拉格朗日量，我们可以计算规范涨落的低能动力学性质。

由于非零的"霍尔"电导 $\sigma_{xy}$，我们可以\emph{不}在导带中产生自旋子、也不在价带中产生空穴而改变自旋子密度。这一性质对理解准粒子的性质十分重要。让我们缓慢地打开 $a_{\mu}$ 场的通量 $\Phi=\int\mathrm{d}^2x\,b$。如果通量分布在一个大区域，通量中的"磁"场 $b$ 就很小。此时价带与导带之间的能隙保持有限，价带中的所有能级都被一个自旋向上和一个自旋向下的自旋子填满。打开通量会感应出一个环形"电"场 $e_{\theta}$，由于非零的 $\sigma_{xy}$，它在径向 $\hat{\pmb{r}}$ 产生自旋子流。于是，一些电荷积累在原点附近。我们发现，被感应出的自旋子总数为
\begin{equation*}
N\equiv\sum_i(\langle f^{\dagger}_{\alpha i}f_{\alpha i}\rangle-1)=-\sigma_{xy}\Phi=-\frac{\Phi}{\pi}
\tag{9.1.29}
\end{equation*}

我们想强调，通量只改变自旋子密度，不感应任何自旋量子数。无论通量感应出多少自旋子，通量管携带的自旋总为零，因为价带中的每个能级都被一个自旋向上和一个自旋向下的自旋子填满。

现在我们可以讨论手征自旋态中准粒子激发的量子数了。平均场手征自旋态中最简单的激发可以通过往导带中加入一个自旋子得到。然而，这个激发不是物理的，因为加入导带的这个多余自旋子违反约束 (9.1.5)。为满足约束，我们可以加入通量来改变价带中的自旋子密度。由上面的讨论可以看到，导带中自旋子所产生的多余自旋子密度，可以通过引入 $\pi$ 通量来抵消（见式 (9.1.29)）。因此，手征自旋态中的物理准粒子是经 $\pi$ 通量\emph{缀饰}的自旋子。缀饰后的自旋子携带自旋 $1/2$，因为通量不能感应任何自旋量子数。然而，作为电荷与通量的束缚态（注意自旋子携带 $a_{\mu}$ 规范场的一个单位电荷），自旋子具有分数统计。交换两个缀饰自旋子等价于把一个自旋子绕另一个运动半周，这会产生相位 $\theta=\frac{1}{2}q\Phi$。这里 $q=1$ 是自旋子的电荷，$\Phi=\pi$ 是束缚在自旋子上的通量。于是，缀饰自旋子的统计角为 $\theta=\frac{\pi}{2}$。具有这种统计的粒子恰介于玻色子（$\theta=0$）与费米子（$\theta=\pi$）之间，称为半子（semion）。

为理解自旋子的低能动力学，我们想导出自旋子的有效拉格朗日量。先通过令 $a_{\mu}=0$ 忽略规范涨落。此时，导带或价带中的自旋子由式 (9.1.25) 给出的色散描述。这里 $E^+_{\pmb{k}}$（$E^-_{\pmb{k}}$）在 $(0,0)$ 与 $(0,\pi/a)$ 有两个极小（极大）。因此，在连续极限下我们有四种自旋子，由下列有效拉格朗日量描述：
\begin{equation*}
\mathcal{L}=\sum_{I=1,2}\left[-\mathrm{i}f^{\dagger}_{I\alpha}\partial_tf_{I\alpha}+\frac{1}{2m_s}f^{\dagger}_{I\alpha}\partial_i^2f_{I\alpha}-\mathrm{i}\bar{f}^{\dagger}_{I\alpha}\partial_t\bar{f}_{I\alpha}+\frac{1}{2m_s}\bar{f}^{\dagger}_{I\alpha}\partial_i^2\bar{f}_{I\alpha}\right]
\tag{9.1.30}
\end{equation*}
这里 $f_{1\alpha}$ 与 $f_{2\alpha}$ 对应导带两个极小附近的自旋子，$\bar{f}_{1\alpha}$ 与 $\bar{f}_{2\alpha}$ 对应价带极大附近的空穴。当 $a_{\mu}\neq0$ 时，自旋子与规范场的耦合可以把式 (9.1.30) 中的 $\partial_{\mu}$ 换成 $\partial_{\mu}\pm\mathrm{i}a_{\mu}$ 而得到。这种耦合形式由规范不变性的要求决定。把与规范场的耦合包括进来之后，总的有效拉格朗日量具有如下形式：
\begin{equation*}
\begin{aligned}
\mathcal{L}={}&\sum_{I=1,2}\left[-\mathrm{i}f^{\dagger}_{I\alpha}(\partial_t-\mathrm{i}a_0)f_{I\alpha}+\frac{1}{2m_s}f^{\dagger}_{I\alpha}(\partial_i-\mathrm{i}a_i)^2f_{I\alpha}\right]\\
&+\sum_{I=1,2}\left[-\mathrm{i}\bar{f}^{\dagger}_{I\alpha}(\partial_t+\mathrm{i}a_0)\bar{f}_{I\alpha}+\frac{1}{2m_s}\bar{f}^{\dagger}_{I\alpha}(\partial_i+\mathrm{i}a_i)^2\bar{f}_{I\alpha}\right]\\
&-\frac{2}{4\pi}a_{\mu}\partial_{\nu}a_{\lambda}\epsilon_{\mu\nu\lambda}+\frac{1}{8\pi g^2}(\mathbf{e}^2-v^2b^2)
\end{aligned}
\tag{9.1.31}
\end{equation*}
由运动方程 $\frac{\delta\mathcal{L}}{\delta a_0}=0$ 得到
\begin{equation*}
n_1+n_2=-\frac{1}{\pi}b
\tag{9.1.32}
\end{equation*}
其中 $n_I=\sum_{\alpha}(f^{\dagger}_{I\alpha}f_{I\alpha}-\bar{f}^{\dagger}_{I\alpha}\bar{f}_{I\alpha})$，$I=1,2$，是自旋子的密度。式 (9.1.32) 告诉我们，自旋子经 $\pi$ 通量缀饰，与先前的结果一致。自旋子的统计也可以直接由式 (9.1.31) 计算。

我们想指出，按照式 (9.1.31)，对每个固定的 $I$ 有两种准粒子，即导带中的自旋 $1/2$ 自旋子与价带中的自旋 $1/2$ 空穴。这个结果是\emph{不正确的}。对每个固定的 $I$ 值，只应有一种准粒子，即自旋 $1/2$ 自旋子。导带中的自旋子与价带中的空穴在投影之后给出同一个自旋子（Affleck et al., 1988; Dagotto et al., 1988）。在意识到手征自旋拟设实际上具有 $SU(2)$ 规范结构之后，这一重复计数问题即可解决（见问题 9.2.4）。

\subsection*{问题 9.1.6}

通过极小化平均场能量 (9.1.6)，求出手征自旋态拟设 (9.1.22) 中决定 $\chi_1$ 与 $\chi_2$ 的方程。

\subsection*{问题 9.1.7}

证明式 (9.1.25)。

\subsection*{问题 9.1.8}

利用 4.4 节的结果证明式 (9.1.27)。

\section{$SU(2)$ 投影构造}

\subsection{隐藏的 $SU(2)$ 规范结构}

\begin{keypoints}
\item 投影构造具有一个隐藏的 $SU(2)$ 规范结构。
\item 平均场拟设 $(U_{ij},a^l_0(\pmb{i}))$ 是物理自旋液体的多对一标记。$SU(2)$ 规范等价的拟设标记同一个物理波函数。
\item 物理波函数在一个对称变换下的不变性，只要求相应的平均场拟设在计及一个规范变换的意义下不变。
\end{keypoints}

\subsubsection{自旋子对凝聚}

在上一节中，我们用 Chern--Simons 项构造了一个稳定的平均场态。在这一节中，我们想考虑因 Anderson--Higgs 机制而稳定的另一类平均场态。为使 Anderson--Higgs 机制奏效，我们首先需要一个携带 $a_{\mu}$ 荷的玻色子。其次，带电玻色子应当具有适当的动力学，使它们能够凝聚。在由投影构造得到的平均场理论中，我们没有携带 $a_{\mu}$ 荷的玻色子，但有携带 $a_{\mu}$ 荷的费米子。于是，我们可以用带电费米子的对来构造带电玻色子，并让这些费米子对凝聚。我们看到，Anderson--Higgs 机制可以通过费米子对凝聚实现。\footnote{如果费米子是电子，那么费米子对凝聚态就是 BCS 超导态。}下面我们将把费米子对凝聚纳入我们的平均场拟设。我们希望这些平均场拟设代表稳定的平均场态。

记得用自旋子算符表示时，哈密顿量 (9.1.1) 可以改写为
\begin{equation*}
H=\sum_{\langle ij\rangle}-\frac{1}{2}J_{ij}\left(f^{\dagger}_{i\alpha}f_{j\alpha}f^{\dagger}_{j\beta}f_{i\beta}+\frac{1}{2}f^{\dagger}_{i\alpha}f_{i\alpha}f^{\dagger}_{j\beta}f_{j\beta}\right)
\tag{9.2.1}
\end{equation*}
其中我们加入了适当的常数项 $\sum_if^{\dagger}_{i\alpha}f_{i\alpha}$ 以得到上述结果。

为得到包含费米子对凝聚的平均场哈密顿量，我们把算符 $f^{\dagger}_{i\alpha}f_{j\beta}$ 与 $f_{i\alpha}f_{i\beta}$ 都换成它们的基态期望值
\begin{equation*}
\begin{array}{ll}
\eta_{ij}\epsilon_{\alpha\beta}=-2\langle f_{i\alpha}f_{j\beta}\rangle, & \eta_{ij}=\eta_{ji},\\
\chi_{ij}\delta_{\alpha\beta}=2\langle f^{\dagger}_{i\alpha}f_{j\beta}\rangle, & \chi_{ij}=\chi^{\dagger}_{ji}.
\end{array}
\tag{9.2.2}
\end{equation*}
$\chi_{ij}$ 项已包含在先前的平均场拟设中。$\eta_{ij}$ 项是新的，它描述费米子对凝聚。我们还把约束 (9.1.4) 换成其基态平均
\begin{equation*}
\langle f^{\dagger}_{i\alpha}f_{i\alpha}\rangle=1,\qquad \langle f_{i\alpha}f_{i\beta}\epsilon_{\alpha\beta}\rangle=0
\tag{9.2.3}
\end{equation*}
这一约束可以通过引入依赖格点且不依赖时间的拉格朗日乘子（Lagrangian multiplier）$a^l_0(\pmb{i})$，$l=1,2,3$，来实施。这样，我们得到下列包含费米子对凝聚的零阶平均场哈密顿量：
\begin{equation*}
\begin{aligned}
H_{\text{mean}}={}&\sum_{\langle ij\rangle}-\frac{3}{8}J_{ij}\left[(\chi_{ji}f^{\dagger}_{i\alpha}f_{j\alpha}+\eta_{ij}f^{\dagger}_{i\alpha}f^{\dagger}_{j\beta}\epsilon_{\alpha\beta}+h.c)-|\chi_{ij}|^2-|\eta_{ij}|^2\right]\\
&+\sum_i\left[a^3_0(f^{\dagger}_{i\alpha}f_{i\alpha}-1)+[(a^1_0+\mathrm{i}a^2_0)f_{i\alpha}f_{i\beta}\epsilon_{\alpha\beta}+h.c.]\right]
\end{aligned}
\tag{9.2.4}
\end{equation*}
这里，式 (9.1.6) 中的 $\chi_{ij}$ 与 $\eta_{ij}$ 必须满足自洽条件 (9.2.2)，而依赖格点的场 $a^l_0(\pmb{i})$ 的选取须使平均场基态满足式 (9.2.3)。这样的 $\chi_{ij}$、$\eta_{ij}$ 与 $a^l_0$ 给出一个平均场解。

对于具有近邻自旋耦合的海森堡模型（见式 (9.1.15)），下列带有费米子对凝聚的平均场拟设是平均场方程 (9.2.2) 的一个解：
\begin{equation*}
\begin{array}{ll}
\chi_{i,i+\pmb{x}}=\chi, & \chi_{i,i+\pmb{y}}=\chi,\\
\eta_{i,i+\pmb{x}}=\eta, & \eta_{i,i+\pmb{y}}=-\eta,\\
a^l_0=0. &
\end{array}
\tag{9.2.5}
\end{equation*}
这样的拟设实际上对应一个 \emph{d} 波 BCS 态。我们称之为 \emph{d} 波态。\footnote{我们注意到，费米子配对序参量 $\eta_{ij}$ 在 $90^{\circ}$ 旋转下变号。这类似于携带角动量 2 的波函数。这就是我们把这个态称为 \emph{d} 波态的原因。}由于费米子对凝聚，我们预期由上述拟设描述的平均场态因 Anderson--Higgs 机制而不含有无能隙 $U(1)$ 规范玻色子。

好，尽管一切显得非常契合、物理图像也显得非常合理，上述结果却被证明是不正确的。我们在哪里犯了错误？这里给出的物理图像与推理是正确的，错误原来出在数学上。我们曾宣称，平均场理论 (9.1.10)（或更一般地，式 (9.2.4)）具有 $U(1)$ 规范结构，使得由 $U(1)$ 规范变换 (9.1.13) 相联系的两个拟设在投影之后对应同一个物理自旋波函数（见式 (9.1.14)）。事实上，平均场理论具有更大的 $SU(2)$ 规范结构。由 $SU(2)$ 规范变换相联系的两个拟设对应同一个物理自旋波函数。不同的规范结构对我们理解平均场态的性质有深远的影响。特别是，$\chi_{ij}$、$\eta_{ij}$ 与 $a^l_0(\pmb{i})$ 的涨落描述 $SU(2)$ 规范涨落。

\subsubsection{$SU(2)$ 形式}

为理解平均场哈密顿量 (9.2.4) 与约束 (9.1.4) 中的 $SU(2)$ 规范结构（Affleck et al., 1988; Dagotto et al., 1988），我们引入二重态（doublet）
\begin{equation*}
\psi=\binom{\psi_1}{\psi_2}=\binom{f_{\uparrow}}{f^{\dagger}_{\downarrow}}
\tag{9.2.6}
\end{equation*}
和一个矩阵
\begin{equation*}
U_{ij}=\begin{pmatrix}\chi^{\dagger}_{ij} & \eta_{ij}\\ \eta^{\dagger}_{ij} & -\chi_{ij}\end{pmatrix}=U^{\dagger}_{ji}
\tag{9.2.7}
\end{equation*}
利用式 (9.2.6) 与 (9.2.7)，我们可以把式 (9.2.3) 与 (9.2.4) 改写为
\begin{equation*}
\langle\psi^{\dagger}_{\pmb{i}}\tau^l\psi_{\pmb{i}}\rangle=0,\qquad l=1,2,3
\tag{9.2.8}
\end{equation*}
\begin{equation*}
H_{\text{mean}}=\sum_{\langle ij\rangle}\frac{3}{8}J_{ij}\left[\frac{1}{2}\operatorname{Tr}(U^{\dagger}_{ij}U_{ij})-(\psi^{\dagger}_iU_{ij}\psi_j+h.c.)\right]+\sum_i a^l_0\psi^{\dagger}_i\tau^l\psi_i
\tag{9.2.9}
\end{equation*}
其中 $\tau^l$，$l=1,2,3$，是泡利矩阵。由式 (9.2.9) 可以清楚地看到，哈密顿量在局域 $SU(2)$ 变换 $W_i$ 下不变：
\begin{equation*}
\begin{array}{c}
\psi_i\rightarrow W_i\psi_i\\
U_{ij}\rightarrow W_iU_{ij}W^{\dagger}_j
\end{array}
\tag{9.2.10}
\end{equation*}

$SU(2)$ 规范结构实际上起源于式 (9.1.2)。这里的 $SU(2)$ 是使物理自旋算符保持不变的自旋子之间最一般的变换。由于物理哈密顿量是自旋算符的函数，这些变换就成为规范变换。

\subsubsection{$SU(2)$ 规范变换的含义}

与 $U(1)$ 规范变换一样，$SU(2)$ 规范变换有如下含义：由 $SU(2)$ 规范变换相联系的两个拟设对应同一个物理自旋波函数。为看到这一点，我们注意到，每个格点上的 $f$ 费米子态与 $\psi$ 费米子态有下列关系：
\begin{equation*}
\begin{array}{ll}
|0_f\rangle=\psi^{\dagger}_2|0_{\psi}\rangle, & f^{\dagger}_{\uparrow}f^{\dagger}_{\downarrow}|0_f\rangle=\psi^{\dagger}_1|0_{\psi}\rangle,\\
f^{\dagger}_{\downarrow}|0_f\rangle=|0_{\psi}\rangle, & f^{\dagger}_{\uparrow}|0_f\rangle=\psi^{\dagger}_1\psi^{\dagger}_2|0_{\psi}\rangle
\end{array}
\end{equation*}
其中 $|0_{\psi}\rangle$ 是没有 $\psi$ 费米子的态，即 $\psi_a|0_{\psi}\rangle=0$。于是，物理的每格点单 $f$ 费米子态对应每格点 $\psi$ 费米子数目为偶数的态。这些每格点偶 $\psi$ 费米子态是局域 $SU(2)$ 单态（即它们在每个格点上都是 $SU(2)$ 单态）。空的 $\psi$ 费米子态对应自旋向下，双占据的 $\psi$ 费米子态对应自旋向上。有了这一理解，我们就可以把平均场态投影到每格点偶 $\psi$ 费米子子空间，得到自旋系统的合法波函数。设 $\pmb{i}_1,\pmb{i}_2,\dots$ 为自旋向上的位置。物理自旋波函数可以写成 $\pmb{i}_1,\pmb{i}_2,\dots$ 的函数，即 $\Psi_{\text{spin}}(\pmb{i}_1,\pmb{i}_2,\dots)$。由于 $\pmb{i}_1,\pmb{i}_2,\dots$ 是仅有的具有两个 $\psi$ 费米子的格点，我们得到
\begin{equation*}
\Psi_{\text{spin}}(\{\pmb{i}_n\})=\langle 0_{\psi}|\prod_n\psi_{1,\pmb{i}_n}\psi_{2,\pmb{i}_n}|\Psi^{(U_{ij})}_{\text{mean}}\rangle.
\tag{9.2.11}
\end{equation*}
由于 $\langle 0_{\psi}|$ 与 $\psi_{1,\pmb{i}}\psi_{2,\pmb{i}}$ 在局域 $SU(2)$ 变换下不变，由 $SU(2)$ 规范变换 $U'_{ij}=W_iU_{ij}W^{\dagger}_j$ 相联系的两个拟设 $U_{ij}$ 与 $U'_{ij}$ 只是标记\emph{同一物理态}的两个不同标签：
\begin{equation*}
\langle 0_{\psi}|\prod_n\psi_{1,\pmb{i}_n}\psi_{2,\pmb{i}_n}|\Psi^{(U_{ij})}_{\text{mean}}\rangle=\langle 0_{\psi}|\prod_n\psi_{1,\pmb{i}_n}\psi_{2,\pmb{i}_n}|\Psi^{(W_iU_{ij}W^{\dagger}_j)}_{\text{mean}}\rangle
\tag{9.2.12}
\end{equation*}

平均场态与物理自旋波函数之间的关系 (9.2.11) 使我们能够从平均场拟设的变换构造出物理自旋波函数的变换。例如，具有平移后拟设 $U'_{ij}=U_{i-l,j-l}$ 的平均场态 $|\Psi^{(U'_{ij})}_{\text{mean}}\rangle$ 在投影之后产生平移了的物理自旋波函数。

显然，平移不变的拟设将导致平移不变的物理自旋波函数。然而，投影后物理波函数的平移对称性并不要求相应拟设的平移不变性。物理态是平移对称的，当且仅当平移后的拟设 $U'_{ij}$ 与原拟设 $U_{ij}$ 规范等价。我们看到，规范结构可能使我们对称性的分析复杂化，因为物理自旋波函数 $\Psi_{\text{spin}}(\{\alpha_i\})$ 可能比投影前的平均场态 $|\Psi^{(U_{ij})}_{\text{mean}}\rangle$ 具有更多的对称性。

\subsubsection{自旋旋转不变性}

我们注意到，$\psi$ 的两个分量都携带自旋向上。因此，自旋旋转对称性在我们的形式体系中不是显式的，很难判断式 (9.2.9) 是否描述一个自旋旋转不变的态。事实上，对满足 $U_{ij}=U^{\dagger}_{ji}$ 的一般 $U_{ij}$，式 (9.2.9) 未必描述自旋旋转不变的态。然而，如果 $U_{ij}$ 具有如下形式
\begin{equation*}
\begin{array}{l}
U_{ij}=\chi^{\mu}_{ij}\tau^{\mu},\qquad \mu=0,1,2,3,\\
\chi^0_{ij}=\text{imaginary},\qquad \chi^l_{ij}=\text{real},\qquad l=1,2,3,
\end{array}
\tag{9.2.13}
\end{equation*}
那么式 (9.2.9) 将描述一个自旋旋转不变的态。这是因为上述 $U_{ij}$ 可以改写成式 (9.2.7) 的形式；此时式 (9.2.9) 可以改写为式 (9.2.4)，其中自旋旋转不变性是显式的。在式 (9.2.13) 中，$\tau^0$ 是单位矩阵，$\tau^{1,2,3}$ 是泡利矩阵。

\subsubsection{变分方法}

我们想说明，除了自洽方程 (9.2.2) 之外，还有另一种求平均场解的办法。我们可以把 $H_{\text{mean}}$（见式 (9.2.9)）的平均场基态，即 $|\Psi^{(U_{ij})}_{\text{mean}}\rangle$，当作尝试波函数，把 $U_{ij}$ 当作变分参数。引入
\begin{equation*}
(\tilde{U}_{ij})_{\alpha\beta}\equiv-2\langle\Psi^{(U_{ij})}_{\text{mean}}|\psi_{i,\alpha}\psi^{\dagger}_{j,\beta}|\Psi^{(U_{ij})}_{\text{mean}}\rangle
\end{equation*}
我们发现 $|\Psi^{(U_{ij})}_{\text{mean}}\rangle$ 的平均场能量为
\begin{equation*}
E_{\text{mean}}(\{U_{ij}\})=-\sum_{\langle ij\rangle}\frac{3}{16}J_{ij}\operatorname{Tr}(\tilde{U}^{\dagger}_{ij}\tilde{U}_{ij})
\tag{9.2.14}
\end{equation*}
这里 $E_{\text{mean}}(\{U_{ij}\})$ 是 $U_{ij}$ 的泛函，具有 $SU(2)$ 规范不变性：
\begin{equation*}
E_{\text{mean}}(\{U_{ij}\})=E_{\text{mean}}(\{W_iU_{ij}W^{\dagger}_j\})
\end{equation*}
平均场解 $U_{ij}$ 现在可以通过极小化 $E_{\text{mean}}$ 得到。

\subsubsection{一阶平均场理论}

为得到一阶平均场理论，我们从由平均场哈密顿量
\begin{equation*}
H_{\text{mean}}=\sum_{\langle ij\rangle}-\frac{3}{8}J_{ij}(\psi^{\dagger}_i\bar{U}_{ij}\psi_j+h.c.)+\sum_i a^l_0\psi^{\dagger}_i\tau^l\psi_i
\tag{9.2.15}
\end{equation*}
描述的零阶平均场理论出发，其中 $\bar{U}_{ij}$ 是平均场解，满足自洽条件
\begin{equation*}
\bar{\chi}_{ij}=\langle\psi^{\dagger}_{i\alpha}\psi_{j\alpha}\rangle,\qquad \bar{\eta}_{ij}=-\langle\psi_{i\alpha}\psi_{i\beta}\epsilon_{\alpha\beta}\rangle
\tag{9.2.16}
\end{equation*}
（译注：式 (9.2.16) 第二个等式中两处下标原书均印作 $i$；按式 (9.2.2)，第二个下标应为 $j$。）

式 (9.2.15) 中的 $a^l_0$ 的选取须使式 (9.2.3) 得到满足。平均场基态附近的重要涨落是 $U_{ij}$ 的下列"相位"涨落：
\begin{equation*}
U_{ij}=\bar{U}_{ij}\mathrm{e}^{\mathrm{i}a_{ij}}
\tag{9.2.17}
\end{equation*}
其中 $a_{ij}=a^l_{ij}\tau^l$ 是一个 $2\times2$ 无迹厄米矩阵。由于 $SU(2)$ 规范结构，这些涨落就是 $SU(2)$ 规范涨落。为了在低能下得到自旋液体态定性正确的结果，我们必须把这些涨落包括进来。这引导我们走向一阶平均场理论
\begin{equation*}
H_{\text{mean}}=\sum_{\langle ij\rangle}-\frac{3}{8}J_{ij}(\psi^{\dagger}_i\bar{U}_{ij}\mathrm{e}^{\mathrm{i}a^l_{ij}\tau^l}\psi_j+h.c.)+\sum_i a^l_0\psi^{\dagger}_i\tau^l\psi_i
\tag{9.2.18}
\end{equation*}
它描述耦合到 $SU(2)$ 格点规范场的自旋子。

\begin{smallnote}
我们想指出，自旋液体的平均场拟设 $U_{ij}$ 可以分为两类：非阻挫拟设——$U_{ij}$ 只把偶格点连接到奇格点；和阻挫拟设——$U_{ij}$ 在两个偶格点之间和/或两个奇格点之间非零。非阻挫拟设的每个方格只有纯 $SU(2)$ 通量，而阻挫拟设在某些方格中，除 $SU(2)$ 通量之外还有 $\pi/2$ 整数倍的 $U(1)$ 通量。
\end{smallnote}

\subsection*{问题 9.2.1}

用式 (9.2.2) 由式 (9.2.1) 得到式 (9.2.4)。

\subsection*{问题 9.2.2}

$\pi$ 通量态与 \emph{d} 波态的等价性

\begin{enumerate}
\item 证明 $\pi$ 通量态 (9.1.18) 由 $SU(2)$ 链变量
\begin{equation*}
U_{i,i+\pmb{x}}=-\mathrm{i}\chi_1\tau^0,\qquad U_{i,i+\pmb{y}}=-\mathrm{i}\chi_1(-)^{i_x}\tau^0
\end{equation*}
描述。
\item 证明当 $\eta=\chi$ 时，\emph{d} 波态 (9.2.5) 由 $SU(2)$ 链变量
\begin{equation*}
U_{i,i+\pmb{x}}=\chi\tau^3+\chi\tau^1,\qquad U_{i,i+\pmb{y}}=\chi\tau^3-\chi\tau^1
\end{equation*}
描述。
\item 证明当 $\chi$ 与 $\chi_1$ 以某种方式相联系时，上述两个拟设规范等价。求把一个拟设变成另一个的 $SU(2)$ 规范变换 $W_i$。（提示：考虑 $SU(2)$ 规范变换 $(\mathrm{i}\tau^1)^{i_x}(\mathrm{i}\tau^3)^{i_y}$。）求 $\chi$ 与 $\chi_1$ 之间的关系。
\item 证明当 $\chi$ 与 $\chi_1$ 以该方式相联系时，$\pi$ 通量态与 \emph{d} 波态具有相同的自旋子谱。
\end{enumerate}

\subsection{$SU(2)$ 规范涨落的动力学}

\begin{keypoints}
\item Anderson--Higgs 机制可以通过凝聚 $SU(2)$ 规范通量来实现，而无需使用希格斯玻色子。
\item 共线的 $SU(2)$ 通量把 $SU(2)$ 规范结构破缺为 $U(1)$ 规范结构。
% 【块界备注】书页 379（p394）止于此。上述 keypoints 列表的第三条要点印于书页 380（p395）顶部：
%   "Non-collinear SU(2) flux breaks the SU(2) gauge structure down to a Z2 gauge structure.
%    In this case, all of the SU(2) gauge bosons gain a gap."
%   建议译文："非共线的 $SU(2)$ 通量把 $SU(2)$ 规范结构破缺为 $Z_2$ 规范结构。此时，所有的
%   $SU(2)$ 规范玻色子都获得能隙。"
%   该条要点与其后的正文（"Just like the $U(1)$ projective construction, ..."）均属 ch9_c
%   （书页 380 起）。请 ch9_c 在其文件头 %--B-TAIL-- 区或正文开头另起 keypoints 环境续排
%   该条要点（本块环境已闭合，勿并入），勿弃勿重。

% ------------------ 新增术语（ch9_b，供术语表.md "新增术语"区追加） ------------------
% | 英文 | 中文 |
% |---|---|
% | chiral spin state | 手征自旋态 |
% | chiral spin liquid | 手征自旋液体 |
% | semion | 半子 |
% | slave fermion | 从属费米子 |
% | valence band | 价带 |
% | conduction band | 导带（ch5_a 已收） |
% | Landau band | 朗道能带 |
% | (un)frustrated ansatz | （非）阻挫拟设 |
% | Lagrangian multiplier | 拉格朗日乘子 |
% | commensuration | 公度 |
% | pi-flux phase / state | $\pi$ 通量相 / $\pi$ 通量态 |
% | d-wave state | d 波态 |
% | doublet | 二重态 |
% | dressed spinon | 缀饰自旋子 |
% | flux tube | 通量管（ch7_a 的"磁通管"为磁场情形） |
% | Anderson–Higgs mechanism | Anderson–Higgs 机制 |
% | over-counting problem | 重复计数问题 |
% | Dirac algebra | 狄拉克代数 |
% | massless Dirac fermion | 无质量狄拉克费米子 |
% | traceless hermitian matrix | 无迹厄米矩阵 |
\item 非共线的 $SU(2)$ 通量把 $SU(2)$ 规范结构破缺到 $Z_2$ 规范结构。此时，所有的 $SU(2)$ 规范玻色子都获得能隙。
\end{keypoints}

正如 $U(1)$ 投影构造一样，为了在 $SU(2)$ 投影构造中得到一个稳定的平均场态，我们需要找到消去无能隙 $SU(2)$ 规范玻色子的办法。$SU(2)$ Chern--Simons 项是赋予 $SU(2)$ 规范玻色子非零能隙的一种方式。下面我们将讨论如何用 Anderson--Higgs 机制消去无能隙的 $SU(2)$ 规范玻色子。

我们知道，$U(1)$ 规范玻色子不携带自身的规范电荷，因此需要额外的带电玻色子来实现 Anderson--Higgs 机制。与之相反，$SU(2)$ 规范玻色子自身就携带非零的规范电荷，所以我们不需要额外的希格斯玻色子来实现 Anderson--\allowbreak Higgs 机制。$SU(2)$ 规范玻色子有能力杀死它们自己。

为了看清 $SU(2)$ 规范玻色子如何自杀，我们考虑一个格点 $SU(2)$ 规范理论。格点 $SU(2)$ 规范场由链变量 $U_{ij}\in SU(2)$ 给出。一阶平均场理论 (9.2.18) 就是一个 $SU(2)$ 格点规范理论。一个位形的能量是 $U_{ij}$ 的函数，即 $E(U_{ij})$。该能量在 $SU(2)$ 规范变换
\begin{equation*}
E(\tilde{U}_{ij})=E(U_{ij}),\qquad \tilde{U}_{ij}=W_i(U_{ij})W_j,\qquad W_i\in SU(2)
\tag{9.2.19}
\end{equation*}
下不变。

为了理解格点 $SU(2)$ 规范涨落的动力学，我们把 $U_{ij}$ 写成 $U_{ij}=\bar{U}_{ij}\mathrm{e}^{\mathrm{i}a_{ij}^{l}\tau^{l}}$，其中链接上的 $2\times 2$ 矩阵 $a_{ij}$ 描写规范涨落。于是能量可以写成 $E(\bar{U}_{ij},\mathrm{e}^{\mathrm{i}a_{ij}^{l}\tau^{l}})$。要判断 $SU(2)$ 规范涨落是否获得能隙，我们需要考察在 $a_{ij}^{l}$ 小的极限下，$E(\bar{U}_{ij},\mathrm{e}^{\mathrm{i}a_{ij}^{l}\tau^{l}})$ 中是否含有质量项 $(a_{ij}^{l})^{2}$。

为了理解平均场拟设 $\bar{U}_{ij}$ 如何影响规范涨落的动力学，引入平均场解的下列圈变量（loop variable）是方便的：
\begin{equation*}
P(C_i)=\bar{U}_{ij}\bar{U}_{jk}\dots\bar{U}_{li}
\tag{9.2.20}
\end{equation*}
这里 $P(C_i)$ 称为通过回路 $C_i$（由 $i\to j\to k\to\dots\to l\to i$ 给出，基点（base point）为 $i$）的 $SU(2)$ 通量，我们也称之为 $SU(2)$ 通量算符。在连续极限下，该圈变量对应于规范场强。在规范变换下，$P(C_i)$ 按如下方式变换：
\begin{equation*}
P(C_i)=W_iP(C_i)W_i
\tag{9.2.21}
\end{equation*}

我们注意到，$SU(2)$ 通量具有 $P(C)=\chi^{0}(C)\tau^{0}+\mathrm{i}\chi^{l}(C)\tau^{l}$ 的形式。因此，当 $\chi^{l}\neq 0$ 时，$SU(2)$ 通量在 $SU(2)$ 空间中有一个方向，该方向由 $\chi^{l}$ 指示。由式 (9.2.21) 可见，局域 $SU(2)$ 规范变换会转动 $SU(2)$ 通量的方向。由于具有不同基点的回路的 $SU(2)$ 通量方向可以被局域 $SU(2)$ 规范变换独立地转动，直接比较不同基点的 $SU(2)$ 通量方向是没有意义的。但是，比较具有同一基点的回路的 $SU(2)$ 通量方向却十分有意义。基于通过具有同一基点的回路的 $SU(2)$ 通量，我们可以把不同的 $SU(2)$ 通量位形分成以下三类：平凡 $SU(2)$ 通量，其中所有 $P(C)\propto\tau^{0}$；共线 $SU(2)$ 通量，其中所有 $SU(2)$ 通量都指向同一方向；非共线 $SU(2)$ 通量，其中具有同一基点的各回路的 $SU(2)$ 通量方向不同。

首先，让我们考虑对所有回路都具有平凡 $SU(2)$ 通量的拟设 $\bar{U}_{ij}$。$SU(2)$ 通量在 $SU(2)$ 规范变换下不变。我们可以（通过施行规范变换）选出一个平均场拟设，使得所有 $\bar{U}_{ij}\propto\tau^{0}$。此时，能量的规范不变性意味着
\begin{equation*}
E(\bar{U}_{ij},\mathrm{e}^{\mathrm{i}a_{ij}^{l}\tau^{l}})=E(\bar{U}_{ij},\mathrm{e}^{\mathrm{i}\theta_i^{l}\tau^{l}}\mathrm{e}^{\mathrm{i}a_{ij}^{l}\tau^{l}}\mathrm{e}^{-\mathrm{i}\theta_j^{l}\tau^{l}}).
\tag{9.2.22}
\end{equation*}
于是，在 $E$ 的展开中，质量项 $(a_{ij}^{1})^{2}$、$(a_{ij}^{2})^{2}$ 和 $(a_{ij}^{3})^{2}$ 一个都不允许出现。\footnote{在规范变换 $\mathrm{e}^{\mathrm{i}\theta_i^{1}\tau^{1}}$ 下，$a_{ij}^{1}$ 按 $a_{ij}^{\prime 1}=a_{ij}^{1}+\theta_i^{1}-\theta_j^{1}$（译注：原文左端无撇号）变换。质量项 $(a_{ij}^{1})^{2}$ 在这样的变换下不是不变的，因而不被允许。}因此，$SU(2)$ 规范涨落是无能隙的，并出现在低能处。由于拟设 $\bar{U}_{ij}\propto\tau^{0}$ 在整体 $SU(2)$ 规范变换下不变，我们说 $SU(2)$ 规范结构未破缺。

第二，让我们假设 $SU(2)$ 通量是共线的。这意思是，具有同一基点的不同回路的 $SU(2)$ 通量都指向同一方向。不过，具有不同基点的回路的 $SU(2)$ 通量仍可能指向不同方向（即使对共线 $SU(2)$ 通量也是如此）。利用局域 $SU(2)$ 规范变换，我们总能把不同基点的 $SU(2)$ 通量转到同一方向，并且可以选这个方向为 $\tau^{3}$ 方向。此时，所有 $SU(2)$ 通量都具有 $P(C)\propto\chi^{0}(C)+\mathrm{i}\chi^{3}(C)\tau^{3}$ 的形式。我们可以（通过施行 $SU(2)$ 规范变换）选出一个平均场拟设，使得所有 $\bar{U}_{ij}$ 都具有 $\mathrm{i}\,\mathrm{e}^{\mathrm{i}\phi_{ij}\tau^{3}}$ 的形式。由于该拟设在整体 $U(1)$ 规范变换 $\mathrm{e}^{\mathrm{i}\theta\tau^{3}}$ 下不变，而在 $\mathrm{e}^{\mathrm{i}\theta\tau^{1,2}}$ 下不然，我们说 $SU(2)$ 规范结构破缺到了 $U(1)$ 规范结构。能量的规范不变性意味着
\begin{equation*}
E(\bar{U}_{ij},\mathrm{e}^{\mathrm{i}a_{ij}^{l}\tau^{l}})=E(\bar{U}_{ij},\mathrm{e}^{\mathrm{i}\theta_i\tau^{3}}\mathrm{e}^{\mathrm{i}a_{ij}^{l}\tau^{l}}\mathrm{e}^{-\mathrm{i}\theta_j\tau^{3}}).
\tag{9.2.23}
\end{equation*}
当 $a_{ij}^{1,2}=0$ 时，上式约化为
\begin{equation*}
E(\bar{U}_{ij},\mathrm{e}^{\mathrm{i}a_{ij}^{3}\tau^{3}})=E(\bar{U}_{ij},\mathrm{e}^{\mathrm{i}(a_{ij}^{3}+\theta_i-\theta_j)\tau^{3}}).
\tag{9.2.24}
\end{equation*}
我们发现，质量项 $(a^{3})^{2}$ 与式 (9.2.24) 不相容，因此至少规范玻色子 $a^{3}$ 是无质量的。那么 $a^{1}$ 和 $a^{2}$ 规范玻色子又如何？令 $P_A(\pmb{i})$ 为通过基点为 $\pmb{i}$ 的回路的 $SU(2)$ 通量。如果我们假设一切能够出现在能量函数中的规范不变项都确实出现在能量函数中，则 $E(U_{ij})$ 将含有下列规范不变项：
\begin{equation*}
E=a\,\mathrm{Tr}[P_A(\pmb{i})U_{\pmb{i},\pmb{i}+\pmb{x}}P_A(\pmb{i}+\pmb{x})U_{\pmb{i}+\pmb{x},\pmb{i}}]+\dots
\tag{9.2.25}
\end{equation*}
如果把 $U_{\pmb{i},\pmb{i}+\pmb{x}}$ 写成 $\chi\,\mathrm{e}^{\mathrm{i}\phi_{ij}\tau^{3}}\mathrm{e}^{\mathrm{i}a_{x}^{l}\tau^{l}}$，利用 $U_{\pmb{i},\pmb{i}+\pmb{x}}=-U_{\pmb{i}+\pmb{x},\pmb{i}}^{\dagger}$ 这一事实（见式 (9.2.13)），并展开到 $(a_{x}^{l})^{2}$ 阶，则式 (9.2.25) 变为
\begin{equation*}
E=-\frac{1}{2}a\chi^{2}\mathrm{Tr}([P_A,a_{x}^{l}\tau^{l}]^{2})+\dots
\tag{9.2.26}
\end{equation*}
由式 (9.2.26) 我们看到，若 $P_A\propto\tau^{3}$，则 $a^{1}$ 与 $a^{2}$ 的质量项会被产生。

总结起来，我们发现：若 $SU(2)$ 通量是共线的，则拟设在 $U(1)$ 转动 $\mathrm{e}^{\mathrm{i}\theta\pmb{n}\cdot\pmb{\tau}}$（$\pmb{n}$ 为 $SU(2)$ 通量的方向）下不变。于是 $SU(2)$ 规范结构破缺到 $U(1)$ 规范结构，相应的平均场态将具有无能隙的 $U(1)$ 规范涨落。

第三，我们考虑 $SU(2)$ 通量非共线的情形。在上面我们已经证明，一个 $SU(2)$ 通量 $P_A$ 可以诱生 $\mathrm{Tr}([P_A,a_{x}^{l}\tau^{l}]^{2})$ 形式的质量项。对于非共线的 $SU(2)$ 通量位形，可以有两个指向不同方向的 $SU(2)$ 通量 $P_A$ 与 $P_B$，质量项还将含有 $\mathrm{Tr}([P_B,a_{x}^{l}\tau^{l}]^{2})$ 项。这样一来，$SU(2)$ 规范场所有的质量项，即 $(a_{ij}^{1})^{2}$、$(a_{ij}^{2})^{2}$ 和 $(a_{ij}^{3})^{2}$，都会被产生。所有的 $SU(2)$ 规范玻色子都将获得能隙（见 3.7.5 节），拟设所描写的平均场态将是一个稳定态。由于该拟设在 $Z_2$ 变换 $W_i=-\tau^{0}$ 下不变，而在其他更一般的整体 $SU(2)$ 规范变换下不然，非共线 $SU(2)$ 通量把 $SU(2)$ 规范结构破缺到 $Z_2$ 规范结构。于是我们可以猜测，低能有效理论是一个 $Z_2$ 规范理论。在 9.2.4 节和 9.3 节中，我们将研究具有非共线 $SU(2)$ 通量的态的低能性质。我们将证明，这类态的低能性质——例如 $Z_2$ 涡旋的存在和基态简并——确实与 $Z_2$ 规范理论的相同。因此，我们将把具有非共线 $SU(2)$ 通量的平均场态称为平均场 $Z_2$ 态。

在平均场 $Z_2$ 态中，所有规范涨落都有能隙。这些涨落只能在自旋子之间传递短程相互作用。低能自旋子相互作用很弱，表现得像自由自旋子。因此，把平均场涨落包括进来并不会定性地改变平均场态的性质。平均场态在低能下是稳定的。

一个稳定的平均场自旋液体态暗示着一个真实的物理自旋液体的存在。稳定平均场态的物理性质适用于该物理自旋液体。如果我们相信这两条论断，那么就可以通过研究相应的稳定平均场态来研究物理自旋液体的性质。由于在平均场 $Z_2$ 态中自旋子不被禁闭，由平均场 $Z_2$ 态得到的物理自旋液体含有中性的自旋-1/2 费米子作为其激发。意识到自旋模型是一个纯玻色模型之后，这是一个十分惊人的结果。

\subsection*{问题 9.2.3}

(a) 证明拟设 (9.2.5) 是一个 $SU(2)$ 共线态。

(b) 找出把 $U_{ij}$ 变换到 $\mathrm{e}^{\mathrm{i}\phi_{ij}\tau^{3}}$ 形式的 $SU(2)$ 规范变换。（提示：试试 $W_i=(\mathrm{i}\tau^{3})^{i_x+i_y}$。）

\subsection*{问题 9.2.4}

(a) 证明由式 (9.1.22) 得到的手征自旋拟设 $U_{ij}$ 在整体 $SU(2)$ 规范变换 $U_{ij}=WU_{ij}W^{\dagger}$ 下不变。

(b) 证明 $a_0^{l}=0$ 的手征自旋拟设 $U_{ij}$ 满足约束 (9.2.8)。

(c) 证明式 (9.1.22) 所描写的拟设与下列平移不变拟设
\begin{equation*}
\begin{aligned}
u_{\pmb{i},\pmb{i}+\pmb{x}} &= -\chi\tau^{3}-\chi\tau^{1}, &\qquad u_{\pmb{i},\pmb{i}+\pmb{y}} &= -\chi\tau^{3}+\chi\tau^{1},\\
u_{\pmb{i},\pmb{i}+\pmb{x}+\pmb{y}} &= \eta\tau^{2}, &\qquad u_{\pmb{i},\pmb{i}-\pmb{x}+\pmb{y}} &= -\eta\tau^{2},\\
a_0^{l} &= 0. &&
\end{aligned}
\end{equation*}
规范等价。用 $f$ 费米子的语言（见式 (9.2.4) 与 (9.2.7)），上述拟设描写一个 $d_{x^{2}-y^{2}}+\mathrm{i}d_{xy}$“超导”态。

由于 $SU(2)$ 规范结构未破缺，手征自旋态的低能有效理论实际上是一个 $SU(2)$ Chern--Simons 理论（level 为 1）。

\subsection{由平移不变拟设得到的自旋液体}

\begin{keypoints}
\item 均匀 RVB 态、$\pi$ 通量态与交错通量态（staggered-flux state）。
\item 投影后的物理波函数的对称性。
\item 低能 $SU(2)$ 或 $U(1)$ 规范涨落。
\end{keypoints}

作为 $SU(2)$ 投影构造的一个应用，我们将研究由平移不变拟设
\begin{equation*}
U_{\pmb{i}+\pmb{l},\pmb{j}+\pmb{l}}=U_{\pmb{i}\pmb{j}},\qquad a_0^{l}(\pmb{i})=a_0^{l},
\end{equation*}
得到的自旋液体。这样的拟设将给出具有平移对称性的物理波函数。\footnote{这里我们要区分拟设的不变性与拟设的对称性。当拟设自身在平移下不改变时，我们说拟设具有平移不变性；当由拟设得到的物理自旋波函数具有平移对称性时，我们说拟设具有平移对称性。由于规范结构的存在，一个不具有平移不变性的拟设也可以具有平移对称性。}首先，让我们引入
\begin{equation*}
u_{ij}=\frac{3}{8}J_{ij}U_{ij}=u_{ij}^{\mu}\tau^{\mu}
\end{equation*}
其中 $u_{\pmb{l}}^{1,2,3}$ 为实数，$u_{\pmb{l}}^{0}$ 为虚数。在零阶平均场理论中，自旋子谱由哈密顿量（见式 (9.2.9)）
\begin{equation*}
H=-\sum_{\langle ij\rangle}\left(\psi_i^{\dagger}u_{ij}\psi_j+h.c.\right)+\sum_i\psi_i^{\dagger}a_0^{l}\tau^{l}\psi_i
\tag{9.2.27}
\end{equation*}
决定。在 $\pmb{k}$ 空间中，我们有
\begin{equation*}
H=-\sum_{\pmb{k}}\psi_{\pmb{k}}^{\dagger}(u^{\mu}(\pmb{k})-a_0^{\mu})\tau^{\mu}\psi_{\pmb{k}}
\end{equation*}
其中 $\mu=0,1,2,3$（译注：原文此处印作“0, 1, 2,”），
\begin{equation*}
u^{\mu}(\pmb{k})=\sum_{\pmb{l}}u_{\pmb{i},\pmb{i}+\pmb{l}}^{\mu}\mathrm{e}^{\mathrm{i}\pmb{l}\cdot\pmb{k}},
\end{equation*}
$a_0^{0}=0$，而 $N$ 为格点总数。费米子谱有两支，由
\begin{align*}
E_{\pm}(\pmb{k})&=u^{0}(\pmb{k})\pm E_0(\pmb{k})\\
E_0(\pmb{k})&=\sqrt{\sum_{\pmb{l}}(u^{l}(\pmb{k})-a_0^{l})^{2}}
\tag{9.2.28}
\end{align*}
给出。约束可以由 $\partial E_{\mathrm{ground}}/\partial a_0^{l}=0$ 得到，其形式为
\begin{equation*}
N\langle\psi_i^{\dagger}\tau^{l}\psi_i\rangle=\sum_{\pmb{k},E_-(\pmb{k})<0}\frac{u^{l}(\pmb{k})-a_0^{l}}{E_0(\pmb{k})}-\sum_{\pmb{k},E_+(\pmb{k})<0}\frac{u^{l}(\pmb{k})-a_0^{l}}{E_0(\pmb{k})}=0
\tag{9.2.29}
\end{equation*}
由此可以定出 $a_0^{l}$，$l=1,2,3$。

让我们尝试下列简单拟设\footnote{该拟设实际上就是前面引入的 $d$ 波拟设 (9.2.5)。}（Baskaran et al., 1987; Affleck and Marston, 1988; Kotliar and Liu, 1988）：
\begin{equation*}
a_0^{l}=0,\qquad u_{\pmb{i},\pmb{i}+\pmb{x}}=\chi\tau^{3}+\eta\tau^{1},\qquad u_{\pmb{i},\pmb{i}+\pmb{y}}=\chi\tau^{3}-\eta\tau^{1}.
\tag{9.2.30}
\end{equation*}
首先，我们要证明相应的自旋液体具有全部对称性。

为了研究相应物理自旋态的对称性，我们需要求出它的波函数 $|\Psi_{\mathrm{phy}}\rangle$。利用该拟设，可以得到平均场哈密顿量 (9.2.27) 的平均场基态 $|\Psi_{\mathrm{mean}}\rangle$。把 $|\Psi_{\mathrm{mean}}\rangle$ 投影到每格点 $\psi$ 费米子数为偶数的子空间，就得到物理波函数 $|\Psi_{\mathrm{phy}}\rangle$，即 $|\Psi_{\mathrm{phy}}\rangle=\mathcal{P}|\Psi_{\mathrm{mean}}\rangle$。

该拟设已经在这两个平移 $T_x:\,\pmb{i}\to\pmb{i}+\pmb{x}$ 与 $T_y:\,\pmb{i}\to\pmb{i}+\pmb{y}$ 下不变。拟设也在两个宇称变换 $P_x:\,x\to-x$ 与 $P_y:\,y\to-y$ 下不变。宇称变换 $P_{xy}:\,(x,y)\to(y,x)$ 把 $u_{\pmb{i},\pmb{i}+\pmb{x}}$ 变为 $u_{\pmb{i},\pmb{i}+\pmb{y}}$、把 $u_{\pmb{i},\pmb{i}+\pmb{y}}$ 变为 $u_{\pmb{i},\pmb{i}+\pmb{x}}$。变换后的拟设
\begin{equation*}
a_0^{l}=0,\qquad u_{\pmb{i},\pmb{i}+\pmb{x}}=\chi\tau^{3}-\eta\tau^{1},\qquad u_{\pmb{i},\pmb{i}+\pmb{y}}=\chi\tau^{3}+\eta\tau^{1}.
\end{equation*}
与原来的拟设不同。但是请注意，一个整体规范变换 $G_{P_{xy}}(\pmb{i})=\mathrm{i}\tau^{3}$ 把 $u_{ij}$ 变为 $G_{P_{xy}}u_{ij}G_{P_{xy}}^{\dagger}$，它诱导 $(\tau^{1},\tau^{2},\tau^{3})\to(-\tau^{1},-\tau^{2},\tau^{3})$ 的改变。因此，变换后的拟设通过规范变换 $G_{P_{xy}}$ 与原来的拟设规范等价。换句话说，$P_{xy}$ 宇称变换再接着一个规范变换 $G_{P_{xy}}$，保持拟设不变。于是，物理波函数在 $P_{xy}$ 宇称变换下不变。$90^{\circ}$ 转动 $R_{90}$ 由 $P_xP_{xy}$ 生成，故物理波函数也有 $90^{\circ}$ 转动对称性。上述拟设还有时间反演对称性，因为时间反演变换 $u_{ij}\to-u_{ij}$ 再接着规范变换 $G_T(i)=\mathrm{i}\tau^{2}$，保持拟设不变。

总而言之，该拟设在下列对称变换与规范变换的组合 $G_xT_x$、$G_yT_y$、$G_{P_x}P_x$、$G_{P_y}P_y$、$G_{P_{xy}}P_{xy}$ 与 $G_TT$ 下不变，相应的规范变换由
\begin{equation*}
\begin{array}{ccc}
G_x=\tau^{0}, & \quad G_y=\tau^{0}, & \quad G_T=\mathrm{i}\tau^{2},\\
G_{P_x}=\tau^{0}, & \quad G_{P_y}=\tau^{0}, & \quad G_{P_{xy}}=\mathrm{i}\tau^{3}.
\end{array}
\tag{9.2.31}
\end{equation*}
给出。所以，由拟设 (9.2.30) 描写的自旋液体具有全部对称性。我们把这样的态称为对称自旋液体态。

我们注意到，式 (9.2.30) 中的链变量 $u_{ij}$ 在 $\tau^{1,2,3}$ 空间中指向不同的方向，因而不是共线的。这样一来，该拟设在整体 $SU(2)$ 规范变换下并非不变。人们可能因此期待 $SU(2)$ 规范结构破缺到 $Z_2$ 规范结构，从而拟设 (9.2.30) 描写一个 $Z_2$ 自旋液体。实际上这个推理是错误的。正如我们在 9.2.2 节中指出的，要判断 $SU(2)$ 规范结构是否破缺到 $Z_2$ 规范结构，需要检查的是 $SU(2)$ \emph{通量}的共线性，而不是 $SU(2)$ 链变量的共线性。

当 $\chi=\eta$ 或 $\eta=0$ 时，所有回路的 $SU(2)$ 通量 $P_C$ 都是平凡的，即 $P_C\propto\tau^{0}$。此时 $SU(2)$ 规范结构未破缺，$SU(2)$ 规范涨落是无能隙的（在弱耦合极限下）。$\eta=0$ 所描写的自旋液体中的自旋子具有谱 $E_{\pmb{k}}=\pm 2|\chi(\cos(k_x)+\cos(k_y))|$，它有一个大费米面。我们将把这个态称为 $SU(2)$ 无能隙态（文献中称其为均匀 RVB 态）。$\chi=\eta$ 的态只在孤立的 $\pmb{k}$ 点有无能隙自旋子，自旋子谱为 $E_{\pmb{k}}=\pm 2\sqrt{\chi^{2}+\eta^{2}}\sqrt{\cos^{2}(k_x)+\cos^{2}(k_y)}$（见图 9.2）。我们将把这样的态称为 $SU(2)$ 线性态，以强调费米点附近 $E\propto|\pmb{k}|$ 的线性色散。（这种态在文献中以及在 9.1.1 节中称为 $\pi$ 通量态。）$SU(2)$ 线性态的低能有效理论由无质量狄拉克费米子（即自旋子）耦合到 $SU(2)$ 规范场所描写（见问题 9.1.4）。

经过适当的规范变换，$SU(2)$ 无能隙拟设可以改写为
\begin{equation*}
u_{\pmb{i},\pmb{i}+\pmb{x}}=\mathrm{i}\chi,\qquad\qquad u_{\pmb{i},\pmb{i}+\pmb{y}}=\mathrm{i}\chi,
\tag{9.2.32}
\end{equation*}
而 $SU(2)$ 线性拟设可以改写为
\begin{equation*}
u_{\pmb{i},\pmb{i}+\pmb{x}}=\mathrm{i}\chi,\qquad\qquad u_{\pmb{i},\pmb{i}+\pmb{y}}=\mathrm{i}(-)^{i_x}\chi.
\tag{9.2.33}
\end{equation*}
在这些形式中，由于 $u_{ij}\propto\mathrm{i}\tau^{0}$，$SU(2)$ 规范结构明显未破缺。我们可以容易地看到，所有的 $SU(2)$ 通量 $P_{ij\dots k}$ 都是平凡的。\footnote{在后面将要讨论的投影对称群（projective symmetry group）分类中，$SU(2)$ 无能隙拟设 (9.2.32) 标记为 SU2An0，$SU(2)$ 线性拟设 (9.2.33) 标记为 SU2Bn0（见式 (9.4.33)）。}这些拟设也在整体 $SU(2)$ 规范变换下不变。$SU(2)$ 无能隙拟设与 $SU(2)$ 线性拟设中的自旋子在低能下都有有限的相互作用，所以这两个平均场态都是边缘平均场态。这两个平均场态是否对应于真实的物理自旋液体，目前还不清楚。

当 $\chi\neq\eta$ 且 $\chi,\eta\neq 0$ 时，通量 $P_C$ 是非平凡的，但 $SU(2)$ 通量是共线的。此时，$SU(2)$ 规范结构破缺到 $U(1)$ 规范结构。无能隙自旋子仍然只出现在孤立的 $\pmb{k}$ 点，从自旋子谱 $E_{\pmb{k}}=\pm 2\sqrt{\chi^{2}(\cos k_x+\cos k_y)^{2}+\eta^{2}(\cos k_x-\cos k_y)^{2}}$ 可以看出这一点。我们将把这样的态称为 $U(1)$ 线性态。（这个态在文献中称为交错通量态和/或 $d$ 波配对态。）经过适当的规范变换，$U(1)$ 线性态也可以用拟设（见问题 9.2.3）
\begin{equation*}
u_{\pmb{i},\pmb{i}+\pmb{x}}=\mathrm{i}\chi-(-)^{i}\eta\tau^{3},\qquad u_{\pmb{i},\pmb{i}+\pmb{y}}=\mathrm{i}\chi+(-)^{i}\eta\tau^{3},
\tag{9.2.34}
\end{equation*}
来描写，其中 $U(1)$ 规范结构是显式的。\footnote{在投影对称群分类中，这样的态标记为 U1Cn01n（见式 (9.4.25)--(9.4.29)）。}其低能有效理论由无质量狄拉克费米子（即自旋子）耦合到 $U(1)$ 规范场所描写。同样，由于无能隙 $U(1)$ 规范场在低能诱导的有限相互作用，$U(1)$ 线性态也是一个边缘平均场态。

\subsection{由平移不变拟设得到的稳定 $Z_2$ 自旋液体}

\begin{keypoints}
\item 稳定的自旋液体态——$Z_2$ 自旋液体态。
\item $Z_2$ 自旋液体含有分数化的激发——自旋子。自旋子携带自旋-1/2、不带电荷，具有费米统计。
\end{keypoints}

在学会了构造稳定平均场态的配方之后，本节我们研究一个不破坏任何对称性的稳定平均场态。该拟设具有下列形式（Wen, 1991a）：
\begin{align*}
u_{\pmb{i},\pmb{i}+\pmb{x}}=u_{\pmb{i},\pmb{i}+\pmb{y}}=-\chi\tau^{3},\qquad & a_0^{2,3}=0,\quad a_0^{1}\neq 0,\\
u_{\pmb{i},\pmb{i}+\pmb{x}+\pmb{y}}=\eta\tau^{1}+\lambda\tau^{2},\qquad & u_{\pmb{i},\pmb{i}-\pmb{x}+\pmb{y}}=\eta\tau^{1}-\lambda\tau^{2}
\tag{9.2.35}
\end{align*}
为了证明式 (9.2.35) 描写一个稳定的平均场态，我们只需证明它产生非共线的 $SU(2)$ 通量。围绕两个具有同一基点的三角形，我们得到下列 $SU(2)$ 通量：
\begin{align*}
u_{\pmb{i},\pmb{i}+\pmb{x}}u_{\pmb{i}+\pmb{x},\pmb{i}+\pmb{y}}u_{\pmb{i}+\pmb{y},\pmb{i}}&=-(\eta\tau^{1}-\lambda\tau^{2})\chi^{2}\\
u_{\pmb{i},\pmb{i}+\pmb{y}}u_{\pmb{i}+\pmb{y},\pmb{i}-\pmb{x}}u_{\pmb{i}-\pmb{x},\pmb{i}}&=-(\eta\tau^{1}+\lambda\tau^{2})\chi^{2}
\end{align*}
我们看到，若 $\chi\eta\lambda\neq 0$，则 $SU(2)$ 通量非共线，式 (9.2.35) 中的拟设描写一个稳定的平均场态，更确切地说，一个 $Z_2$ 态。

由于拟设在平移下不变，物理自旋态在平移下也不变，于是物理波函数 $|\Psi_{\mathrm{phy}}\rangle$ 具有平移对称性。在 $P_x$ 宇称变换下，拟设变为
\begin{align*}
u_{\pmb{i},\pmb{i}+\pmb{x}}=u_{\pmb{i},\pmb{i}+\pmb{y}}=-\chi\tau^{3},\qquad & a_0^{2,3}=0,\quad a_0^{1}\neq 0,\\
u_{\pmb{i},\pmb{i}+\pmb{x}+\pmb{y}}=\eta\tau^{1}-\lambda\tau^{2},\qquad & u_{\pmb{i},\pmb{i}-\pmb{x}+\pmb{y}}=\eta\tau^{1}+\lambda\tau^{2}
\tag{9.2.36}
\end{align*}
（注意，对平移不变拟设，$P_x$ 把 $u_{\pmb{i},\pmb{i}+\pmb{x}}$ 变为 $u_{\pmb{i},\pmb{i}-\pmb{x}}=u_{\pmb{i},\pmb{i}+\pmb{x}}^{\dagger}$）。变换后的拟设在 $SU(2)$ 规范变换 $G_x(i)=\mathrm{i}\tau^{1}(-)^{i}$ 下与原来的拟设规范等价。\footnote{我们定义 $(-)^{i}\equiv(-)^{i_x+i_y}$。规范变换 $W_i=\mathrm{i}\tau^{1}$ 把式 (9.2.36) 变换为
\begin{equation*}
\begin{aligned}
u_{\pmb{i},\pmb{i}+\pmb{x}}&=u_{\pmb{i},\pmb{i}+\pmb{y}}=\chi\tau^{3},\qquad & a_0^{2,3}&=0,\quad a_0^{1}\neq 0,\\
u_{\pmb{i},\pmb{i}+\pmb{x}+\pmb{y}}&=\eta\tau^{1}+\lambda\tau^{2},\qquad & u_{\pmb{i},\pmb{i}-\pmb{x}+\pmb{y}}&=\eta\tau^{1}-\lambda\tau^{2}
\end{aligned}
\end{equation*}
然后，规范变换 $W_i=(-)^{i}$ 再把上式变为式 (9.2.35)。}尽管 $P_x$ 变换后的拟设与原来的拟设不同，两者在投影后给出同一个物理波函数。物理自旋态具有 $P_x$ 宇称对称性。

我们还注意到，若 $(a_0^{2},a_0^{3})$ 不为零，上述组合变换 $G_xP_x$ 会把它们变成 $(-a_0^{2},-a_0^{3})$。由于平均场哈密顿量具有 $P_x$ 对称性，非零 $(a_0^{2},a_0^{3})$ 的平均场基态能量满足性质 $E_{\mathrm{ground}}(a_0^{2},a_0^{3})=E_{\mathrm{ground}}(-a_0^{2},-a_0^{3})$。因此，对我们的拟设而言，$a_0^{2}=a_0^{3}=0$ 总能满足约束 $\langle\psi_i^{\dagger}\tau^{2}\psi_i\rangle\allowbreak=N_{\mathrm{site}}^{-1}\partial E_{\mathrm{ground}}/\partial a_0^{2}\allowbreak=0$ 与 $\langle\psi_i^{\dagger}\tau^{3}\psi_i\rangle\allowbreak=N_{\mathrm{site}}^{-1}\partial E_{\mathrm{ground}}/\partial a_0^{3}\allowbreak=0$。我们只需调节 $a_0^{1}$ 使 $\langle\psi_i^{\dagger}\tau^{1}\psi_i\rangle=0$。

$P_y$ 宇称变换的作用与 $P_x$ 相同，所以物理波函数也有 $P_y$ 宇称对称性。在时间反演变换下，$T$ 把 $(u_{ij};a_0)$ 变为 $(-u_{ij};-a_0)$。如果我们选 $G_T=\mathrm{i}\tau^{3}(-)^{i}$，时间反演变换再接着规范变换 $G_T$ 将保持拟设 (9.2.35) 不变。因此，$|\Psi_{\mathrm{phy}}\rangle$ 具有时间反演对称性。该拟设在 $(x,y)\to(y,x)$ 宇称 $P_{xy}$ 下不变，所以物理波函数具有 $P_{xy}$ 宇称对称性。我们看到，拟设 (9.2.35) 在下列组合变换 $G_xT_x$、$G_yT_y$、$G_{P_x}P_x$、$G_{P_y}P_y$、$G_{P_{xy}}P_{xy}$ 与 $G_TT$ 下不变，其中规范变换为
\begin{equation*}
\begin{array}{ccc}
G_x=\tau^{0}, & \quad G_y=\tau^{0}, & \quad G_T=\mathrm{i}(-)^{i}\tau^{3},\\
G_{P_x}=\tau^{0}, & \quad G_{P_y}=\tau^{0}, & \quad G_{P_{xy}}=\mathrm{i}\tau^{3}
\end{array}
\tag{9.2.37}
\end{equation*}
该拟设描写一个具有全部对称性的自旋液体。我们把这样的态称为 $Z_2$ 对称自旋液体态。

$Z_2$ 态中的自旋子谱由
\begin{align*}
E_{\pm}(\pmb{k})&=\pm\sqrt{\epsilon_1^{2}+\epsilon_2^{2}+\epsilon_3^{2}}\\
\epsilon_1&=2\chi(\cos k_x+\cos k_y)\\
\epsilon_2&=2\eta[\cos(k_x+k_y)+\cos(k_x-k_y)]+a_0^{1}\\
\epsilon_3&=2\lambda[\cos(k_x+k_y)-\cos(k_x-k_y)]
\end{align*}
给出。我们发现自旋子是完全有能隙的。我们将把态 (9.2.35) 称为 $Z_2$ 有能隙自旋液体。\footnote{这里我们想指出，在 9.4.3 节将要讨论的投影对称群分类中，态 (9.2.35) 所标记的投影对称群为 Z2Axx0z。}$Z_2$ 有能隙自旋液体对应于 Kivelson et al. (1987) 以及 Rokhsar and Kivelson (1988) 提出的短程共振价键（sRVB）态。

由于式 (9.2.35) 是一个稳定拟设，规范涨落全都有能隙，只在自旋子之间产生短程相互作用。由于自旋子有能隙，$Z_2$ 有能隙态的低能激发对应于自旋子的稀薄气体。由于相互作用是短程的，这些自旋子表现得像自由费米子。因此，$Z_2$ 有能隙态中的激发是自由的中性自旋-1/2 费米子。

\begin{figure}[H]
\centering
\includegraphics[width=0.72\textwidth]{fig_9.4.pdf}
\caption*{图 9.4\quad 沿粗线所示的链接翻转 $U_{ij}$ 的符号，就产生一个 $Z_2$ 涡旋。这里 $\Theta_{ij}=1$，唯有粗线所示的链接上 $\Theta_{ij}=-1$。}
\end{figure}

\subsection{$Z_2$ 自旋液体中的 $Z_2$ 涡旋}

\begin{keypoints}
\item $Z_2$ 自旋液体态含有一种 $Z_2$ 涡旋激发，它既不携带自旋也不携带电荷。
\item 把一个 $Z_2$ 涡旋绑定到自旋子上，会使自旋子的统计性从玻色性变为费米性，或从费米性变为玻色性。
\end{keypoints}

除自旋子外，$Z_2$ 自旋液体态中还有另一类准粒子激发（Wen, 1991a）。这种激发在平均场理论中以拓扑孤子（topological soliton）的形式出现，对应于 $Z_2$ 规范场的 $\pi$ 通量，因此我们将把这种新激发称为 $Z_2$ 涡旋。它通过沿一串链接翻转 $u_{ij}$ 的符号产生（见图 9.4），由下列拟设描写：
\begin{equation*}
\tilde{u}_{ij}=\bar{u}_{ij}\Theta_{ij}
\tag{9.2.38}
\end{equation*}
其中 $\Theta_{ij}$ 已在图 9.4 中说明。$Z_2$ 涡旋位于弦的端点。注意，在远离 $Z_2$ 涡旋处，$\tilde{u}_{ij}$ 局域地规范等价于 $\bar{u}_{ij}$，因此符号翻转的 $u_{ij}$ 弦不耗费任何能量，也观测不到。拟设 $\tilde{u}_{ij}$ 实际上描写的是位于弦端点的一个局域激发。

由于 $u_{ij}$ 决定自旋子的跳跃幅，当自旋子绕 $Z_2$ 涡旋跳行一周时会多出一个负号。因此，对自旋子而言 $Z_2$ 涡旋表现得像一个 $\pi$ 通量。根据 7.1.2 节的讨论，我们发现把一个 $Z_2$ 涡旋绑定到一个自旋子上，会把自旋子的统计性从费米性变为玻色性。因此，$Z_2$ 自旋液体既含有玻色统计的中性自旋-1/2 激发，也含有费米统计的中性自旋-1/2 激发。

\subsection*{问题 9.2.5}

证明：对于式 (9.2.38) 的 $Z_2$ 涡旋拟设，通过回路 $C$ 的 $SU(2)$ 通量与基态拟设 $u_{ij}$ 通过同一回路的 $SU(2)$ 通量相同，除非回路 $C$ 包围着涡旋。由于平均场能量是 $SU(2)$ 通量的函数，上述结果表明符号翻转的 $u_{ij}$ 弦不耗费能量。

\subsection{稳定的无能隙 $Z_2$ 自旋液体}

\begin{keypoints}
\item 存在许多不同的稳定自旋液体态——$Z_2$ 自旋液体态。其中一些可以具有无能隙的自旋子激发。
\item 这些态具有完全相同的对称性。没有办法利用对称性来刻画这些不同的自旋液体态。
\end{keypoints}

除了式 (9.2.35) 之外，对于 $Z_2$ 对称自旋液体我们还可以写下另一个拟设（Wen, 2002c）：
\begin{equation*}
u_{\pmb{i},\pmb{i}+\pmb{x}}=\mathrm{i}\eta\tau^{0}-\chi(\tau^{3}-\tau^{1}),\qquad u_{\pmb{i},\pmb{i}+\pmb{y}}=\mathrm{i}\eta\tau^{0}-\chi(\tau^{3}+\tau^{1}),\qquad a_0^{l}=0,
\tag{9.2.39}
\end{equation*}
其中 $\chi$ 与 $\eta$ 均不为零。当 $\chi=0$ 时，上述拟设成为 $SU(2)$ 无能隙自旋液体；当 $\eta=0$ 时，成为 $SU(2)$ 线性自旋液体。为了证明式 (9.2.39) 的拟设描写一个对称自旋液体，我们必须证明它在下列对称变换与规范变换的组合 $G_xT_x$、$G_yT_y$、$G_{P_x}P_x$、$G_{P_y}P_y$、$G_{P_{xy}}P_{xy}$ 与 $G_TT$ 下不变。我们发现，如果选取规范变换为
\begin{equation*}
\begin{array}{lll}
G_x=\tau^{0}, & \quad G_y=\tau^{0}, & \quad G_T=(-)^{i}\tau^{0},\\[2pt]
G_{P_x}=\mathrm{i}(-)^{i_x}\dfrac{\tau^{1}+\tau^{3}}{\sqrt{2}}, & \quad G_{P_y}=\mathrm{i}(-)^{i_y}\dfrac{\tau^{1}-\tau^{3}}{\sqrt{2}}, & \quad G_{P_{xy}}=\mathrm{i}\tau^{3},
\end{array}
\tag{9.2.40}
\end{equation*}
则对称变换再接着相应的规范变换，将保持拟设不变。

利用时间反演对称性，我们可以证明拟设 (9.2.39) 中取零值的 $a_0^{l}$ 确实满足约束 (9.2.29)。这是因为在组合时间反演变换 $G_TT$ 下 $a_0^{l}\to-a_0^{l}$。由于 $a_0^{l}=0$ 的拟设在 $G_TT$ 下不变，非零 $a_0^{l}$ 的平均场基态能量满足 $E_{\mathrm{ground}}(a_0^{l})=E_{\mathrm{ground}}(-a_0^{l})$，于是当 $a_0^{l}=0$ 时 $\partial E_{\mathrm{ground}}/\partial a_0^{l}=0$。事实上，任何只在两个不交叠子格之间有链接的拟设（即非阻挫拟设），只要 $a_0^{l}=0$，就是时间反演对称的。对这类拟设（包括拟设 (9.2.39)），取零值的 $a_0^{l}$ 满足约束 (9.2.29)。

自旋子谱由下式给出（见图 9.10(a)）
\begin{equation*}
E_{\pm}=2\eta(\sin(k_x)+\sin(k_y))\pm 2|\chi|\sqrt{2\cos^{2}(k_x)+2\cos^{2}(k_y)}
\end{equation*}
自旋子有两个费米点和两个小的费米口袋（Fermi pocket）（当 $\eta$ 小时）。$SU(2)$ 通量是非平凡的。而且，$SU(2)$ 通量 $P_{C_1}$ 与 $P_{C_2}$ 互不对易，其中 $C_1=\pmb{i}\to\pmb{i}+\pmb{x}\to\pmb{i}+\pmb{x}+\pmb{y}\to\pmb{i}+\pmb{y}\to\pmb{i}$ 与 $C_2=\pmb{i}\to\pmb{i}+\pmb{y}\to\pmb{i}-\pmb{x}+\pmb{y}\to\pmb{i}-\pmb{x}\to\pmb{i}$ 是两个具有同一基点的回路。非共线 $SU(2)$ 通量把 $SU(2)$ 规范结构破缺到 $Z_2$ 规范结构。我们将把式 (9.2.39) 描写的自旋液体称为 $Z_2$ 无能隙自旋液体。\footnote{$Z_2$ 无能隙自旋液体是 9.4.3 节中分类的 $Z_2$ 自旋液体之一，其投影对称群标记为 Z2Ax12(12)n（见 9.4.3 节及式 (9.4.16)）。}其低能有效理论由无质量狄拉克费米子和具有小费米面的费米子耦合到 $Z_2$ 规范场所描写。由于没有无能隙规范玻色子，规范涨落只能在自旋子之间传递短程相互作用。短程相互作用在低能下是不相关的，自旋子在低能下是自由费米子。因此，平均场 $Z_2$ 无能隙态是一个稳定态。这个稳定的平均场态暗示存在第二种真实的物理自旋液体（除 9.2.4 节讨论的 $Z_2$ 有能隙自旋液体之外）。这种自旋液体的物理性质可以由平均场理论得到。我们发现，$Z_2$ 无能隙自旋液体中的激发由中性的自旋-1/2 费米型准粒子描写！

第三种 $Z_2$ 对称自旋液体可以由下列阻挫拟设得到（Balents et al., 1998; Senthil and Fisher, 2000）：
\begin{align*}
a_0^{3}\neq 0,\qquad & a_0^{1,2}=0,\\
u_{\pmb{i},\pmb{i}+\pmb{x}}=\chi\tau^{3}+\eta\tau^{1},\qquad & u_{\pmb{i},\pmb{i}+\pmb{y}}=\chi\tau^{3}-\eta\tau^{1},\\
u_{\pmb{i},\pmb{i}+\pmb{x}+\pmb{y}}=+\gamma\tau^{3},\qquad & u_{\pmb{i},\pmb{i}-\pmb{x}+\pmb{y}}=+\gamma\tau^{3}.
\tag{9.2.41}
\end{align*}
该拟设具有平移、转动、宇称与时间反演对称性。它在组合变换 $G_{x,y}T_{x,y}$、$G_{P_x,P_y,P_{xy}}P_{x,y,xy}$ 与 $G_TT$ 下不变，其中
\begin{equation*}
\begin{array}{ccc}
G_x=\tau^{0}, & \quad G_y=\tau^{0}, & \quad G_T=\mathrm{i}\tau^{2},\\
G_{P_x}=\tau^{0}, & \quad G_{P_y}=\tau^{0}, & \quad G_{P_{xy}}=\mathrm{i}\tau^{3}.
\end{array}
\tag{9.2.42}
\end{equation*}
当 $a_0^{3}\neq 0$、$\chi\neq\pm\eta$ 且 $\chi\eta\neq 0$ 时，我们发现 $a_0^{l}\tau^{l}$ 与 $SU(2)$ 通量算符——例如穿过一个方格的 $SU(2)$ 通量 $P(C_{\mathrm{square}})$——不对易。\footnote{这里 $\tau^{l}a_0^{l}(\pmb{i})$ 可以看作通过一个零尺寸回路的 $SU(2)$ 通量。}因此，该拟设把 $SU(2)$ 规范结构破缺到 $Z_2$ 规范结构。自旋子谱由下式给出（见图 9.8(a)）
\begin{align*}
E_{\pm}&=\pm\sqrt{\epsilon^{2}(\pmb{k})+\Delta^{2}(\pmb{k})}\\
\epsilon(\pmb{k})&=2\chi(\cos(k_x)+\cos(k_y))+2\gamma(\cos(k_x+k_y)+\cos(k_x-k_y))+a_0^{3}\\
\Delta(\pmb{k})&=2\eta(\cos(k_x)-\cos(k_y))+a_0^{3}
\end{align*}
它只在该四个 $\pmb{k}$ 点无能隙，且具有线性色散。因此，式 (9.2.41) 描写的自旋液体是一个 $Z_2$ 线性自旋液体。\footnote{$Z_2$ 线性自旋液体由投影对称群 Z2A0032 或等价的 Z2A0013 描写（见 9.4.3 节）。}

$Z_2$ 有能隙拟设 (9.2.35)、$Z_2$ 无能隙拟设 (9.2.39) 与 $Z_2$ 线性拟设 (9.2.41) 在投影后给出三个自旋液体态。这三个态具有完全相同的对称性，所以我们不能用对称性来区分它们。为了找到一组新的量子数来刻画这三种不同的自旋液体，我们注意到：虽然三种 $Z_2$ 自旋液体具有相同的对称性，但它们的拟设分别在同样的对称平移之后跟着\emph{不同的}规范变换下不变（见式 (9.2.37)、(9.2.40) 与 (9.2.42)）。因此，三种自旋液体的平均场拟设的不变群各不相同，我们可以用这样的不变群来刻画不同的自旋液体。更详细和系统的讨论将在 9.4 节给出。

\subsection*{问题 9.2.6}

证明：如果规范变换由式 (9.2.40) 给出，则式 (9.2.39) 中的拟设在组合变换 $G_xT_x$、$G_yT_y$、$G_{P_x}P_x$、$G_{P_y}P_y$、$G_{P_{xy}}P_{xy}$ 与 $G_TT$ 下不变。

\subsection*{问题 9.2.7}

证明手征自旋态具有平移对称性和 $90^{\circ}$ 转动对称性。证明手征自旋态还具有组合对称性 $TP_x$、$TP_y$ 与 $TP_{xy}$。

\subsection{评注：平均场理论中的时间反演变换}

在时间反演变换（或自旋反转变换）下，$\pmb{S}_i\to-\pmb{S}_i$。由于本章只考虑自旋旋转不变的态，为方便起见，我们把 $T$ 变换定义为时间反演变换与自旋旋转变换的组合：$(S^{x},S^{y},S^{z})\to(S^{x},-S^{y},S^{z})$。我们将不严格地把它称为时间反演变换。

这样的变换可以由
\begin{equation*}
\psi\to\psi'=-\mathrm{i}\tau^{2}\psi^{*}
\tag{9.2.43}
\end{equation*}
生成。让我们用一个算符 $\mathcal{T}$ 来表示上述变换，即 $\mathcal{T}\psi\mathcal{T}^{-1}=-\mathrm{i}\tau^{2}\psi^{*}$。注意 $\mathcal{T}$ 把 $\mathrm{i}$ 变成 $-\mathrm{i}$，即 $\mathcal{T}^{-1}\mathrm{i}\mathcal{T}=-\mathrm{i}$。用 $f$ 费米子 (9.2.6) 表示，式 (9.2.43) 变为
\begin{equation*}
f\to\mathcal{T}f\mathcal{T}^{-1}=-\mathrm{i}\tau^{2}f^{*}
\end{equation*}
由 $\pmb{S}=\frac{1}{2}f^{\dagger}\pmb{\sigma}f$，我们发现 $\mathcal{T}$ 把 $\pmb{S}$ 变为
\begin{equation*}
\mathcal{T}\pmb{S}\mathcal{T}^{-1}=-\frac{1}{2}f^{\top}\pmb{\sigma}f^{*}=\frac{1}{2}f^{\dagger}\pmb{\sigma}^{\top}f=(S^{x},-S^{y},S^{z})
\end{equation*}
在平均场理论中，式 (9.2.43) 诱导平均场拟设的下列变换：
\begin{equation*}
\begin{aligned}
U_{ij}&\to U'_{ij}=(-\mathrm{i}\tau^{2})U_{ij}^{*}(\mathrm{i}\tau^{2})=-U_{ij}\\
a_0^{l}(\pmb{i})\tau^{l}&\to a_0^{\prime l}(\pmb{i})\tau^{l}=(-\mathrm{i}\tau^{2})(a_0^{l}(\pmb{i})\tau^{l})^{\top}(\mathrm{i}\tau^{2})=-a_0^{l}(\pmb{i})\tau^{l}
\end{aligned}
\tag{9.2.44}
\end{equation*}

% ------------------------------------------------------------
% 新增术语（ch9_c，书页 380–392，供术语表收录）：
%   loop variable 圈变量
%   SU(2)-flux operator SU(2) 通量算符
%   base point 基点
%   trivial / collinear / non-collinear SU(2) flux 平凡 / 共线 / 非共线 SU(2) 通量
%   mass term 质量项（ch3_e 已收）
%   uniform RVB state 均匀 RVB 态
%   staggered-flux state 交错通量态
%   d-wave pairing state d 波配对态
%   symmetric spin-liquid state 对称自旋液体态
%   Z₂-symmetric spin-liquid state Z₂ 对称自旋液体态
%   Z₂-gapped spin liquid Z₂ 有能隙自旋液体
%   Z₂-gapless spin liquid Z₂ 无能隙自旋液体
%   Z₂-linear spin liquid Z₂ 线性自旋液体（SU(2)/U(1)-linear state 同式：线性态）
%   sRVB (short-ranged resonating valence bond) state 短程共振价键（sRVB）态
%   Fermi pocket 费米口袋
%   topological soliton 拓扑孤子
%   unfrustrated / frustrated ansatz 非阻挫 / 阻挫拟设
%   invariant group 不变群
%   level (of Chern–Simons theory) level（Chern–Simons 理论的 level，不译）
%   spin–rotation transformation 自旋旋转变换
%   spin-reversal transformation 自旋反转变换
% ------------------------------------------------------------

% 块边界判定：
%   起点——书页 393（p408）顶部为图 9.5（属 9.2.7 节 T 变换讨论，前块止于书页 392 末的
%     式 (9.2.44)），页面首段 "The above describes ..." 为完整新句，无前块溢出内容，
%     故本块不设 %JOIN%，亦无需 %--C-TAIL-- 区。
%   终点——书页 405（p420）末句 "The key to obtaining the most general Z2A0013 ansatz
%     is to" 延续至书页 406（p421）首行 "find out which $u_m^\mu$ must vanish ..."。
%     按约定，半句译文（"……的关键是"）留在本文件正文最后，不写入任何 TAIL 区，
%     由 ch9_e 在其正文开头以 %JOIN% 续接。
% 蓝本问题记录：式 (9.4.11) 蓝本漏排 $G_y(\pmb{i})=\tau^0$ 一栏且上标误作 i_y（以图为准）；
%   式 (9.4.13) 蓝本作 $(-)^{m_y i_x}$，图上为 $(-)^{m i_x}$；书页 402 "±G_yT_x" 为
%   "±G_yT_y" 之误（据脚注 70 与式 (9.4.6)），已按本意译出并加译注。

\begin{figure}[H]
\centering
\includegraphics[width=0.72\textwidth]{fig_9.5.pdf}
\caption*{图 9.5\quad 跨越 $x$ 线与 $y$ 线的链获得一个额外的负号。}
\end{figure}

上述结果描述了平均场拟设在 $T$ 变换下的变换方式。自旋哈密顿量 (9.1.1) 与平均场哈密顿量 (9.2.4) 在时间反演变换下都是不变的。

\subsection*{问题 9.2.8}

试求在真实的时间反演变换 $S\to-S$ 下，$U_{ij}$ 与 $a_0^l(\pmb{i})$ 如何变换。

\section{有能隙自旋液体态中的拓扑序}

\begin{keypoints}
\item 自旋液体基态的拓扑简并。
\item 稳健的拓扑简并基态的存在意味着拓扑序的存在。
\end{keypoints}

我们用平均场理论构造了 $Z_2$ 有能隙态与手征自旋态。这两个自旋液体态的所有激发都有有限的能隙，并且保持平移对称性。我们想指出，这里研究的手征自旋液体与 $Z_2$ 有能隙自旋液体包含一些新型的序，它们既不能用对称性、也不能用与破缺对称性相联系的序参量来刻画。对 $Z_2$ 有能隙态而言，这一点相当明显，因为它不破缺任何对称性。手征自旋液体自发破缺了时间反演对称性与宇称对称性，并有相应的局域序参量来描述这种对称性破缺。正如我们后面将看到的，手征自旋液体还包含一些无法用对称性破缺描述的额外结构。由于手征自旋液体与 $Z_2$ 有能隙自旋液体都有有限的能隙，这类新型的序就是拓扑序。

正如 8.2.1 节所指出的，证明 $Z_2$ 有能隙态与手征自旋态具有非平庸拓扑序的一种办法，是证明这两个态具有拓扑简并的基态。让我们先考虑 $Z_2$ 有能隙态及其基态简并度。

我们把 $Z_2$ 有能隙态放到一个环面上，$x$ 与 $y$ 两个方向的格点数都取偶数。考虑由一个平均场解 $\bar{u}_{\pmb{ij}}$ 构造出的下列四个拟设：
\begin{equation*}
u_{\pmb{ij}}^{(m,n)}=(-)^{m s_x(\pmb{ij})}(-)^{n s_y(\pmb{ij})}\bar{u}_{\pmb{ij}}
\tag{9.3.1}
\end{equation*}
其中 $m,n=0,1$。这里 $s_{x,y}(\pmb{ij})$ 取 0 或 1：若链 $\pmb{ij}$ 跨越 $x$ 线，则 $s_x(\pmb{ij})=1$（见图 9.5），否则 $s_x(\pmb{ij})=0$；类似地，若链 $\pmb{ij}$ 跨越 $y$ 线，则 $s_y(\pmb{ij})=1$，否则 $s_y(\pmb{ij})=0$。

我们注意到，不同 $m$、$n$ 的 $u_{\pmb{ij}}^{(m,n)}$ 是局域规范等价的。这是因为，在无限系统上，改变 $u_{\pmb{ij}}\to(-)^{m s_x(\pmb{ij})}(-)^{n s_y(\pmb{ij})}u_{\pmb{ij}}$ 可以由一个 $SU(2)$ 规范变换 $u_{\pmb{ij}}\to W_{\pmb{i}}u_{\pmb{ij}}W_{\pmb{j}}^{\dagger}$ 生成，其中 $W_{\pmb{i}}=(-)^{m\Theta(i_x)}(-)^{n\Theta(i_y)}$，而 $\Theta(n)=1$ 若 $n>0$，$\Theta(n)=0$ 若 $n\leqslant0$。结果，不同的拟设 $u_{\pmb{ij}}^{(m,n)}$ 穿过每个方格的 $SU(2)$ 通量都相同。

对大系统而言，能量是 $u_{\pmb{ij}}$ 的局域函数，并且若 $u_{\pmb{ij}}$ 与 $\tilde{u}_{\pmb{ij}}$ 规范等价，则有 $F(u_{\pmb{ij}})=F(\tilde{u}_{\pmb{ij}})$。因此，不同 $u_{\pmb{ij}}^{(m,n)}$ 的能量相同（在热力学极限下）。另一方面，不同 $m$、$n$ 的 $u_{\pmb{ij}}^{(m,n)}$ 在整体意义上并不规范等价。一个自旋子沿 $x$（或 $y$）方向绕环面一整圈，会获得相位 $\mathrm{e}^{\mathrm{i}m\pi}$（或 $\mathrm{e}^{\mathrm{i}n\pi}$）。因此，$u_{\pmb{ij}}^{(mn)}\big|_{m,n=0,1}$ 描述不同的、相互正交的简并基态（另见问题 9.3.1）。这样我们就发现，$Z_2$ 有能隙态在环面上有四个简并基态（Read and Chakraborty, 1989; Wen, 1991a）。从自旋子绕环面一圈所获得的相位可以看出，$m$ 与 $n$ 标记穿过环面两个孔的 $\pi$ 通量的数目（见图 9.6）。四个简并基态对应于把 $\pi$ 通量穿入两个孔的不同方式。这一图像可以推广到更高亏格的黎曼面。简并基态可以这样构造：让 $\pi$ 通量以零个或一个单位穿过亏格为 $g$ 的黎曼面上的 $2g$ 个孔（见图 9.6）。这给出亏格为 $g$ 的黎曼面上 $2^{2g}$ 个简并基态。这种在不同黎曼面上的简并图案是 $Z_2$ 规范理论的标志。这些简并度与 $Z_2$ 涡旋激发的存在，以一种物理的方式表明，$Z_2$ 有能隙态的低能有效理论确实是一个 $Z_2$ 规范理论（见问题 9.3.1）。把自旋子包括进来之后，$Z_2$ 有能隙态的一阶平均场理论描述的是耦合到 $Z_2$ 规范理论上的费米子。$Z_2$ 有能隙态的物理性质可以从这样一个有效理论得到。

\begin{figure}[H]
\centering
\includegraphics[width=0.72\textwidth]{fig_9.6.pdf}
\caption*{图 9.6\quad (a) 亏格 $g=1$ 的黎曼面（环面）上有 $2g=2$ 个孔；(b) 亏格 $g=2$ 的黎曼面上有 $2g=4$ 个孔。我们可以往这些孔里穿入零通量或 $\pi$ 通量。这给出 $Z_2$ 有能隙自旋液体的 $2^{2g}=4$ 或 16 个简并基态。}
\end{figure}

在一个有限的紧致化格点上，系统从一个基态隧穿到另一个基态的唯一途径是如下隧穿过程：先产生一对 $Z_2$ 涡旋，其中一个 $Z_2$ 涡旋绕环面一整圈，然后与另一个 $Z_2$ 涡旋湮没。这样的过程实质上给环面的孔添加了一个单位的 $\pi$ 通量，并使 $m$ 或 $n$ 改变 1。由于通量守恒，不同的基态无法通过任何局域涨落相互隧穿。这一结果的直接推论是，有限格点上不同基态之间的能量劈裂预期为 $\mathrm{e}^{-L/\xi}$ 的量级，其中 $1/\xi$ 与 $Z_2$ 涡旋的有限能隙相关。

在平均场理论中，基态的简并是规范不变性的一个后果。即使我们在原来的自旋哈密顿量中加入下列任意微扰，规范不变性仍然严格成立：
\begin{equation*}
\delta H=\sum\delta J_{ij}\,\pmb{S}_i\cdot\pmb{S}_j+\cdots
\tag{9.3.2}
\end{equation*}
其中 $\delta H$ 可以破坏平移对称性、旋转对称性等。因此上述论证依然成立，即使对修改后的哈密顿量，平均场基态仍保持四重简并。我们预期这一结果在平均场理论之外也成立。只要微扰足够弱、不至于关闭 $Z_2$ 涡旋的能隙，任何微扰都无法改变 $Z_2$ 有能隙态的基态简并度。因此，基态简并度是刻画自旋液体态不同相的一个普适量子数。这也表明 $Z_2$ 有能隙态中存在非平庸的拓扑序。

现在让我们考虑手征自旋态中的拓扑序。同样，我们想计算手征自旋态在比如环面上的基态简并度。计算基态简并度最容易的办法，是使用有效 Chern--Simons 理论 (9.1.28)，它描述 $a_{\mu}$ 集体模式的动力学。$U(1)$ Chern--Simons 理论的基态简并度已在 FQH 态的情形中讨论过。重复 8.2.1 节的计算，我们发现基态简并度是 2。我们想指出，这一二重简并是针对破缺 $T$ 与 $P$ 的基态的某一个扇区（sector）而言的，例如 $\langle\pmb{S}_1\cdot(\pmb{S}_2\times\pmb{S}_3)\rangle>0$ 的扇区。另一个扇区 $\langle\pmb{S}_1\cdot(\pmb{S}_2\times\pmb{S}_3)\rangle<0$ 也有两个简并基态。因此，总的基态简并度是 $4=2\times2$：一个因子 2 来自 $T$、$P$ 破缺，另一个因子 2 来自规范涨落。

由规范涨落引起的二重基态简并，源于手征自旋态中的非平庸拓扑序。这一简并依赖于空间的拓扑。在亏格为 $g$ 的黎曼面上，规范涨落给出 $2^g$ 个简并基态。这一简并度同样直接与规范结构相关，并且可以证明它对自旋哈密顿量的任意局域微扰都是稳健的。正如 FQH 态一样，手征自旋态的基态简并度与自旋子的分数统计密切相关。一般地可以证明，若自旋子具有分数统计 $\theta=\frac{p\pi}{q}$，则手征自旋态在亏格为 $g$ 的黎曼面上将有 $q^g$ 重简并的基态（Wen, 1990）。此外，作为非平庸拓扑序的另一个结果，手征自旋态具有无能隙的手征边缘激发（Wen, 1992）。

为了看出拓扑序概念的用处，让我们考虑下列物理问题：手征自旋态与 $Z_2$ 有能隙态之间的差别是什么？有人可能马上会说，这两个态有不同的对称性。然而，如果我们修改哈密顿量使 $T$ 与 $P$ 对称性破缺，两个态就会有相同的对称性。这时我们仍然可以问这两个态是否相同——即一个态能否不经相变而连续变形为另一个态。当 $T$ 与 $P$ 对称性被显式破缺时，手征自旋态有两个简并基态，而 $Z_2$ 有能隙态在环面上仍有四个简并基态。因此，即使具有相同的对称性，手征自旋态与 $Z_2$ 有能隙态也是不同的。两态的差别在于它们有不同的基态简并度和不同的拓扑序。

\subsection*{问题 9.3.1}

\textbf{非共线 $SU(2)$ 通量态中的 $Z_2$ 规范结构}

设 $u_{\pmb{ij}}=s_{ij}\bar{u}_{\pmb{ij}}$，其中 $s_{ij}=\pm1$，$\bar{u}_{\pmb{ij}}$ 是一个具有非共线 $SU(2)$ 通量的拟设。令 $|\{s_{ij}\}\rangle$ 为 $|\Psi_{\text{mean}}^{(u_{\pmb{ij}})}\rangle$ 的投影（见式 (9.2.11)），即
\begin{equation*}
|\{s_{ij}\}\rangle=\langle0_{\psi}|\prod_n\psi_{1,\pmb{i}_n}\psi_{2,\pmb{i}_n}|\Psi_{\text{mean}}^{(u_{\pmb{ij}})}\rangle.
\end{equation*}
\begin{enumerate}[label=(\alph*)]
\item 证明：与 $Z_2$ 规范理论中的态一样，若 $s_{ij}$ 与 $\tilde{s}_{ij}$ 通过一个 $Z_2$ 规范变换相联系，则 $|\{s_{ij}\}\rangle=|\{\tilde{s}_{ij}\}\rangle$。
\item 证明：式 (9.3.1) 所定义的四个 $u_{\pmb{ij}}^{(m,n)}$（其中 $\bar{u}_{\pmb{ij}}$ 为具有非共线 $SU(2)$ 通量的一般拟设）彼此不 $SU(2)$ 等价。（提示：考虑通过三个具有同一基点的回路 $C_{1,2,3}$ 的 $SU(2)$ 通量 $P(C_1)$、$P(C_2)$ 与 $P(C_3)$。这里 $C_1$ 与 $C_2$ 是两个小回路，满足 $[P(C_1),P(C_2)]\neq0$；$C_3$ 是绕环面一圈的大回路。）
\end{enumerate}

\section{对称自旋液体中的量子序}

\begin{keypoints}
\item 量子序是量子相中的非对称性破缺序。
\item 量子序既适用于有能隙量子态，也适用于无能隙量子态。
\end{keypoints}

在这一节中，我们将基于平均场拟设的不变群发展一种量子序理论。我们将把这个不变群称为投影对称群（projective symmetry group，PSG）。我们的理论使我们能够刻画并对数百种具有相同对称性的自旋液体态进行分类。与拓扑序的基态简并度不同，投影对称群既能刻画有能隙自旋液体，也能刻画无能隙自旋液体。

我们已经看到，可以有许多具有相同对称性的不同自旋液体。其中许多自旋液体占据相空间中的一个有限区域，代表稳定的量子相。于是我们在这里面临的处境与量子霍尔效应中的情形类似：存在许多无法用对称性和序参量区分的量子相。量子霍尔液体有有限的能隙，我们可以用基态简并度或更一般的拓扑序来描述这些不同态的内部序。这里我们仍然可以用拓扑序来描述有能隙自旋液体的内部序。然而，我们还有许多其他稳定的量子自旋液体，它们具有无能隙激发。

为了描述无能隙量子自旋液体（以及有能隙自旋液体）的内部序，我们在第 8 章中引入了一个新概念——量子序（quantum order）——它描述量子相中的非对称性破缺序（见图 8.7）。引入量子序的关键点是：量子相一般不能完全用破缺的对称性和局域序参量来刻画。量子霍尔态以及前几节构造的稳定自旋液体都说明了这一点。然而，要使量子序的概念有用，我们需要找到量子序的具体数学刻画。由于量子序既不由对称性也不由序参量描述，我们需要找到一种全新的刻画方式。Wen (2002c) 提出用投影对称群来刻画量子自旋液体中的量子（或拓扑）序。投影对称群的提出受到 9.2.6 节如下观察的启发：不同对称自旋液体的拟设，在同样的对称性变换之后跟上\emph{不同}的规范变换下不变。我们可以利用这些不同的规范变换来区分具有相同对称性的不同自旋液体。在接下来两节中，我们将在一般而形式化的框架下引入投影对称群。

\subsection{量子序与普适性质}

\begin{keypoints}
\item 定义一种新型的序，就是要找到一种新的普适性质。
\item （某一类）量子序可以用投影对称群描述。
\item 投影对称群是一个平均场拟设的不变群（或“对称性”群）。
\end{keypoints}

我们知道，定义一个量子相，就是要找出基态波函数的普适性质。我们可以利用普适性质把基态波函数归入若干普适类，使得每一类中的波函数具有相同的普适性质。这些普适类将对应于量子相。

然而，要找出多体波函数的普适性质是非常困难的。让我们考虑自由费米系统中的量子序（或普适类），以对这种困难获得一些直观认识。我们知道，自由费米子基态由一个 $N$ 个变量的反对称波函数描述。该反对称函数具有 Slater 行列式的形式，即 $\Psi(x_1,\ldots,x_N)=\det(M)$，其中 $M$ 的矩阵元由 $M_{mn}=\psi_n(x_m)$ 给出，而 $\psi_n$ 是单费米子波函数。寻找自由费米系统中量子序的第一步，是找到一种合理的方式把 Slater 行列式波函数归类。如果我们只知道实空间的多体函数 $\Psi(x_1,\ldots,x_N)$，这将是非常困难的。然而，如果我们用傅里叶变换把实空间波函数变到动量空间的波函数，就可以按照费米面拓扑把不同的波函数归类。要真正证明上述分类是合理的并对应于不同的量子相，我们需要证明：当基态波函数从一个类变到另一个类时，基态能量在转变点处总会出现奇异性。这就引出了我们对自由费米系统中量子序的理解（见 8.3.2 节）。费米面拓扑是一种普适性质，它使我们能够对不同的费米子波函数进行分类，并刻画自由费米子的不同量子相。这里我们想强调，如果没有傅里叶变换，就很难从实空间多体函数 $\Psi(x_1,\ldots,x_N)$ 中看出费米面拓扑。

为了理解自旋液体中的量子序，我们需要找到一种办法把自旋液体波函数 $\Psi_{\text{spin}}(\pmb{i}_1,\ldots,\pmb{i}_{N_t})$ 归类，其中 $\pmb{i}_m$ 标记自旋朝上的位置。这里缺少的是相应的“傅里叶”变换。正如费米面的拓扑一样，很难直接从实空间波函数中看出普适性质（如果有的话）。目前，我们还不知道如何对自旋液体波函数 $\Psi_{\text{spin}}$ 进行分类。

不过，我们不想就此放弃。受投影构造的启发，这里我们考虑一个较简单的问题来入手。我们把自己限制在可以由拟设 $(u_{\pmb{ij}},a_0^l\tau^l)$ 经式 (9.2.11) 得到的那一类多体波函数之内。我们不去寻找一般多体波函数的普适性质，而是设法找出这一子类中多体波函数的普适性质。由于这一子类中的多体波函数由拟设 $(u_{\pmb{ij}},a_0^l\tau^l)$ 标记，波函数的普适性质实际上就是拟设的普适性质，这大大简化了问题。诚然，有人可能会反驳说，拟设（或该子类波函数）的普适性质未必是自旋量子相的普适性质。对某些拟设来说情况确实如此。然而，如果拟设 $(u_{\pmb{ij}},a_0^l\tau^l)$ 所描述的平均场态对任何涨落都是稳定的（即在平均场态附近的任何涨落都不会引起红外发散），那么该平均场态就忠实地描述一个自旋量子态，拟设的普适性质也就是相应自旋量子相的普适性质。这样便在拟设的性质与物理自旋液体的性质之间建立起联系。

那么，拟设的普适性质应当是什么呢？受朗道对称性破缺序理论的启发，我们想在此提出：拟设的不变群（或“对称性”群）是拟设的一种普适性质。这个群将被称为投影对称群（PSG）。我们将在后面论证，某些 PSG 确实是量子相的普适性质，并且这些 PSG 可以用来刻画量子自旋液体中的量子序。

\subsection{投影对称群}

\begin{keypoints}
\item 投影对称群是对称性群的一个扩张。
\item 投影对称群对所有平均场相进行分类。
\item 只有当相应的平均场态稳定时，投影对称群才刻画物理自旋液体态中的量子序。
\end{keypoints}

让我们给出 PSG 的一个详细定义。PSG 是拟设的一种属性。它由保持拟设不变的所有变换构成。每个变换（即 PSG 中的每个元素）都可以写成一个对称性变换 $U$（如平移）与一个规范变换 $G_U$ 的组合。拟设在其 PSG 下的不变性可以表示为
\begin{align*}
G_U U(u_{\pmb{i}\pmb{j}}) &= u_{\pmb{i}\pmb{j}},\tag{9.4.1}\\
U(u_{\pmb{i}\pmb{j}}) &\equiv u_{U(\pmb{i}),U(\pmb{j})},\qquad G_U(u_{\pmb{i}\pmb{j}})\equiv G_U(\pmb{i})u_{\pmb{i}\pmb{j}}G_U^{\dagger}(\pmb{j}),\qquad G_U(\pmb{i})\in SU(2),
\end{align*}
对每个 $G_UU\in\text{PSG}$ 上式都成立。

每个 PSG 都包含一个特殊的子群，我们称之为不变规范群（invariant gauge group，IGG）。一个拟设的 IGG（记作 $\mathcal{G}$）由保持该拟设不变的所有纯规范变换构成：
\begin{equation*}
\mathcal{G}=\{W_{\pmb{i}}\,|\,W_{\pmb{i}}u_{\pmb{i}\pmb{j}}W_{\pmb{j}}^{\dagger}=u_{\pmb{i}\pmb{j}},\ W_{\pmb{i}}\in SU(2)\}
\tag{9.4.2}
\end{equation*}
如果我们想把 IGG 与一个对称性变换联系起来，那么相应的变换就是恒等对称性变换。

如果 IGG 是非平庸的，那么对一个固定的对称性变换 $U$，可以存在许多规范变换 $G_U$ 使得 $G_UU$ 保持拟设不变。若 $G_UU$ 在 $u_{\pmb{i}\pmb{j}}$ 的 PSG 中，则 $GG_UU$ 也在 PSG 中的充要条件是 $G\in\mathcal{G}$。这样，对每个对称性变换 $U$，$G_U$ 的不同选取与 IGG 中的元素一一对应。从上述定义我们看到，一个拟设的 PSG、IGG 与对称性群（symmetry group，SG）之间有如下关系：
\begin{equation*}
SG=PSG/IGG
\end{equation*}
这一关系告诉我们，PSG 是对称性群的一个投影表示，或称一个扩张（extension）。\footnote{更一般地，如果群 PSG 包含一个正规子群 IGG，使得 $PSG/IGG=SG$，我们就说群 PSG 是群 SG 的一个扩张。}

当然，两个规范等价拟设 $u_{\pmb{i}\pmb{j}}$ 与 $W(\pmb{i})u_{\pmb{i}\pmb{j}}W^{\dagger}(\pmb{j})$ 的 PSG 彼此相关。由 $WG_UU(u_{\pmb{i}\pmb{j}})=W(u_{\pmb{i}\pmb{j}})$，其中 $W(u_{\pmb{i}\pmb{j}})\equiv W(\pmb{i})u_{\pmb{i}\pmb{j}}W^{\dagger}(\pmb{j})$，我们发现
\begin{equation*}
WG_UUW^{-1}W(u_{\pmb{i}\pmb{j}})=WG_UW_U^{-1}UW(u_{\pmb{i}\pmb{j}})=W(u_{\pmb{i}\pmb{j}})
\end{equation*}
其中 $W_U\equiv UWU^{-1}$ 由 $W_U(\pmb{i})=W(U(\pmb{i}))$ 给出。这样，若 $G_UU$ 在拟设 $u_{\pmb{i}\pmb{j}}$ 的 PSG 中，则 $(WG_UW_U)U$ 在规范变换后的拟设 $W(\pmb{i})u_{\pmb{i}\pmb{j}}W^{\dagger}(\pmb{j})$ 的 PSG 中。我们看到，在规范变换 $W(\pmb{i})$ 之后，与对称性变换 $U$ 相联系的规范变换 $G_U$ 按如下方式改变：
\begin{equation*}
G_U(\pmb{i})\to W(\pmb{i})G_U(\pmb{i})W^{\dagger}(U(\pmb{i}))
\tag{9.4.3}
\end{equation*}

由于 PSG 是拟设的一种属性，我们可以把共享同一 PSG 的所有拟设归成一类。我们宣称，这样的类是一个普适类，对应于一个量子相。\footnote{更精确地说，这样的类由一个或若干个对应于量子相的普适类构成。（关于这一点的更详细讨论见 9.9 节。）}正是在这个意义上，我们说量子序由 PSG 刻画。

把量子序的 PSG 刻画与对称性破缺序的对称性刻画加以比较是很有启发的。我们知道，对称性破缺序可以由它的对称性性质描述。在数学上，我们说对称性破缺序由它的对称性群刻画。类似地，量子序也由群刻画。差别在于，量子序由 PSG——对称性群的扩张——刻画。我们看到，用投影对称群描述量子序，与用对称性群描述对称性破缺序，在概念上是相似的。我们还看到，具有相同对称性的量子态可以有许多不同的量子序，因为一个对称性群通常可以有许多不同的扩张。

除了能够刻画对称性破缺序之外，对称性破缺序的对称性描述还非常有用，因为它使我们无需了解系统的细节就能得到许多普适性质，例如无能隙 Nambu--Goldstone 模的数目。类似地，知道了一个量子序的 PSG，也使我们无需了解量子系统的细节就能得到它的低能性质。在 9.10 节中我们将证明，PSG 能够产生并保护低能规范涨落。事实上，低能规范涨落的规范群就是拟设的 IGG。这推广了 9.2.2 节得到的结果。除了无能隙规范玻色子之外，PSG 还能产生并保护无能隙费米子激发及其晶体动量。这些无能隙费米子甚至可以在纯玻色模型中产生。我们看到，无能隙规范玻色子与无能隙费米子可以有一个统一的起源——量子序。

这里我们想强调，PSG 真正分类的是不同的平均场相。连续的平均场相变总是伴随着 PSG 的改变。因此，严格说来，上述关于 PSG 及其物理含义的讨论只适用于平均场理论。然而，有些平均场态对涨落是稳定的。这些稳定平均场态的 PSG 就可以应用于相应的物理自旋液体，并描述这些自旋液体中真实的量子序。另一方面，也存在对涨落不稳定的平均场态。在这种情形下，相应的 PSG 可能不描述任何真实的量子序。

\subsection{对称 $Z_2$ 自旋液体的分类}

\begin{keypoints}
\item 我们分类了其 PSG 为格点对称性群（SG）之 $Z_2$ 扩张的所有平均场拟设，即 $PSG/Z_2=SG$。
\item 存在超过 100 种对称 $Z_2$ 自旋液体。
\item 不变 PSG 与代数 PSG 的概念。
\end{keypoints}

作为自旋液体量子序 PSG 理论的一个应用，我们想对与平移变换相联系的 PSG 进行分类。为简单起见，我们把 IGG $\mathcal{G}$ 限制为 $Z_2$。也就是说，我们想找出平移群的所有 $Z_2$ 扩张。

当 $\mathcal{G}=Z_2$ 时，它只含两个元素——规范变换 $G_1$ 与 $G_2$：
\begin{equation*}
\mathcal{G}=\{G_1,G_2\},\qquad G_1(\pmb{i})=\tau^0,\qquad G_2(\pmb{i})=-\tau^0.
\end{equation*}
一个拟设的 IGG 要为 $\mathcal{G}=Z_2$，该拟设就必须产生非共线的 $SU(2)$ 通量；否则 IGG 将大于 $Z_2$。拟设中的非共线 $SU(2)$ 通量把 $SU(2)$ 规范结构破缺到 $Z_2$ 规范结构，相应的自旋液体是一个 $Z_2$ 自旋液体。我们看到，平移群的 $Z_2$ 扩张在物理上对所有\emph{只}具有平移对称性的 $Z_2$ 自旋液体进行了分类。

考虑一个具有平移对称性的 $Z_2$ 自旋液体。作为平移群的 $Z_2$ 扩张，它的 PSG 由四个元素 $\pm G_xT_x$ 与 $\pm G_yT_y$ 生成\footnote{译注：原文作“$\pm G_yT_x$”，据式 (9.4.6) 及脚注（原书脚注 70）应为 $\pm G_yT_y$ 之误，此处按本意译出。}，其中
\begin{equation*}
T_x(u_{\pmb{ij}})=u_{\pmb{i}-\pmb{x},\,\pmb{j}-\pmb{x}},\qquad T_y(u_{\pmb{ij}})=u_{\pmb{i}-\pmb{y},\,\pmb{j}-\pmb{y}}.
\end{equation*}
由于拟设的平移对称性，我们可以选取一个规范，使拟设的所有 $SU(2)$ 通量算符都是平移不变的。也就是说，若两个回路 $C_1$ 与 $C_2$ 通过一个平移相联系，则 $P_{C_1}=P_{C_2}$。我们把这样的规范称为均匀规范（uniform gauge）。

在变换 $G_xT_x$ 下，以 $\pmb{i}$ 为基点的 $SU(2)$ 通量算符 $P_C$ 按如下方式变换：$P_C\to G_x(\pmb{i}')P_{T_x(C)}G_x^{\dagger}(\pmb{i}')\allowbreak=G_x(\pmb{i}')P_CG_x^{\dagger}(\pmb{i}')$，其中 $\pmb{i}'=T_x\pmb{i}$ 是平移后的回路 $T_x(C)$ 的基点。我们看到，$P_C$ 在均匀规范下的平移不变性要求：对任何回路 $C$ 都有 $G_x(\pmb{i}')P_CG_x^{\dagger}(\pmb{i}')=P_C$。同样，$G_y$ 也满足类似的条件。由于以同一基点为基点的不同 $SU(2)$ 通量算符互不对易，$G_{x,y}(\pmb{i})$ 在每个格点只能取 $\pm\tau^0$ 两个值之一。

我们注意到，形如 $W(\pmb{i})=\pm\tau^0$ 的依赖格点的规范变换不改变 $SU(2)$ 通量算符的平移不变性质。因此，我们可以利用这类规范变换，通过式 (9.4.3) 进一步简化 $G_{x,y}$。首先，我们可以选取规范使（见问题 9.4.1）
\begin{equation*}
G_y(\pmb{i})=\tau^0.
\tag{9.4.4}
\end{equation*}
我们注意到，满足 $W(\pmb{i})=W(i_x)$ 的规范变换不改变条件 $G_y(\pmb{i})=\tau^0$。我们可以利用这类规范变换使
\begin{equation*}
G_x(i_x,\ i_y=0)=\tau^0.
\tag{9.4.5}
\end{equation*}
由于 $x$ 与 $y$ 方向的平移相互对易，$G_{x,y}$ 必须满足（对任何拟设，无论是不是 $Z_2$）
\begin{equation*}
G_xT_xG_yT_y(G_xT_x)^{-1}(G_yT_y)^{-1}=G_xT_xG_yT_yT_x^{-1}G_x^{-1}T_y^{-1}G_y^{-1}\in\mathcal{G}.
\tag{9.4.6}
\end{equation*}
这意味着
\begin{equation*}
G_x(\pmb{i})G_y(\pmb{i}-\pmb{x})G_x^{-1}(\pmb{i}-\pmb{y})G_y(\pmb{i})^{-1}\in\mathcal{G}
\tag{9.4.7}
\end{equation*}
对 $Z_2$ 自旋液体，由于式 (9.4.4)，式 (9.4.7) 约化为
\begin{equation*}
G_x(\pmb{i})G_x^{-1}(\pmb{i}-\pmb{y})=+\tau^0
\tag{9.4.8}
\end{equation*}
或
\begin{equation*}
G_x(\pmb{i})G_x^{-1}(\pmb{i}-\pmb{y})=-\tau^0
\tag{9.4.9}
\end{equation*}
与式 (9.4.5) 和 (9.4.4) 联立，我们发现式 (9.4.8) 与 (9.4.9) 成为
\begin{equation*}
G_x(\pmb{i})=\tau^0,\qquad G_y(\pmb{i})=\tau^0
\tag{9.4.10}
\end{equation*}
以及
\begin{equation*}
G_x(\pmb{i})=(-)^{i_x}\tau^0,\qquad G_y(\pmb{i})=\tau^0
\tag{9.4.11}
\end{equation*}
这样，平移群只有两个规范不等价的 $Z_2$ 扩张。这两个 PSG 由 $G_xT_x$ 与 $G_yT_y$ 生成，其中的 $G_{x,y}$ 由式 (9.4.10) 与 (9.4.11) 给出。\footnote{我们想指出，若 $G_x$ 与 $G_y$ 由式 (9.4.8) 给出，则 $G_xT_x$ 与 $G_yT_y$ 满足通常的平移代数 $G_xT_xG_yT_y=G_yT_yG_xT_x$。当 $G_x$ 与 $G_y$ 由式 (9.4.9) 给出时，$G_xT_x$ 与 $G_yT_y$ 满足磁平移代数 $G_xT_xG_yT_y=-G_yT_yG_xT_x$。}在 PSG 分类下，\emph{只}具有平移对称性而没有其他对称性的 $Z_2$ 自旋液体只有两种类型。在 PSG (9.4.10) 下不变的拟设具有形式
\begin{equation*}
u_{\pmb{i},\pmb{i}+\pmb{m}}=u_{\pmb{m}}
\tag{9.4.12}
\end{equation*}
而在 PSG (9.4.11) 下不变的拟设具有形式
\begin{equation*}
u_{\pmb{i},\pmb{i}+\pmb{m}}=(-)^{m i_x}u_{\pmb{m}}
\tag{9.4.13}
\end{equation*}

通过上面的例子我们看到，PSG 是一个非常强有力的工具。对于具有给定对称性和低能规范结构的（平均场）自旋液体，它可以给出完全的分类。

在上文中，我们研究了只具有平移对称性而无其他对称性的 $Z_2$ 自旋液体，发现这样的自旋液体只有两种类型。然而，如果自旋液体具有更多对称性，类型就会多得多。Wen (2002c) 利用 PSG 得到了对称 $Z_2$ 自旋液体的分类。这一分类来自如下观察：规范变换 $G_{x,y}$、$G_{P_x,P_y,P_{xy}}$ 与 $G_T$ 必须满足某些与式 (9.4.7) 类似的代数关系。求解这些代数关系并剔除规范等价的解，我们发现，由平移、宇称与时间反演变换 $\{T_{x,y},P_{x,y,xy},T\}$ 生成的对称性群共有 272 个不同的 $Z_2$ 扩张。这里我们仅列出这 272 个 $Z_2$ PSG。这些 PSG 由 $(G_xT_x,\ G_yT_y,\ G_TT,\ G_{P_x}P_x,\ G_{P_y}P_y,\ G_{P_{xy}}P_{xy})$ 生成。PSG 可以分成两类。第一类为
\begin{equation*}
\begin{array}{rcl}
G_x(\pmb{i})=\tau^0, & & G_y(\pmb{i})=\tau^0\\[2pt]
G_{P_x}(\pmb{i})=\eta_{xpx}^{i_x}\eta_{xpy}^{i_y}g_{P_x}, & & G_{P_y}(\pmb{i})=\eta_{xpy}^{i_x}\eta_{xpx}^{i_y}g_{P_y}\\[2pt]
G_{P_{xy}}(\pmb{i})=g_{P_{xy}}, & & G_T(\pmb{i})=\eta_t^{\pmb{i}}g_T
\end{array}
\tag{9.4.14}
\end{equation*}
第二类为
\begin{equation*}
\begin{array}{rcl}
G_x(\pmb{i})=(-)^{i_x}\tau^0, & & G_y(\pmb{i})=\tau^0\\[2pt]
G_{P_x}(\pmb{i})=\eta_{xpx}^{i_x}\eta_{xpy}^{i_y}g_{P_x}, & & G_{P_y}(\pmb{i})=\eta_{xpy}^{i_x}\eta_{xpx}^{i_y}g_{P_y}\\[2pt]
G_{P_{xy}}(\pmb{i})=(-)^{i_x i_y}g_{P_{xy}}, & & G_T(\pmb{i})=\eta_t^{\pmb{i}}g_T
\end{array}
\tag{9.4.15}
\end{equation*}
这里三个 $\eta$ 可以各自独立地取 $\pm1$ 两个值。$g$ 有 17 种不同的选择，列于表 9.1 之中。因此，共有 $2\times17\times2^{3}=272$ 个不同的 PSG。它们可能给出二维正方格点上 272 种不同类型的对称 $Z_2$ 自旋液体。

为了标记这 272 个 PSG，我们提出如下方案：
\begin{equation*}
\text{Z2A}(g_{px})_{\eta_{xpx}}(g_{py})_{\eta_{xpy}}g_{pxy}(g_t)_{\eta_t},
\tag{9.4.16}
\end{equation*}
\begin{equation*}
\text{Z2B}(g_{px})_{\eta_{xpx}}(g_{py})_{\eta_{xpy}}g_{pxy}(g_t)_{\eta_t}.
\tag{9.4.17}
\end{equation*}
标记 Z2A… 对应式 (9.4.14) 的情形，标记 Z2B… 对应式 (9.4.15) 的情形。一个典型的标记形如 $\text{Z2A}\tau_{+}^{1}\tau_{-}^{2}\tau^{12}\tau_{-}^{3}$。我们还将使用一种缩写记号：把 $(\tau^{0},\tau^{1},\tau^{2},\tau^{3})$ 或 $(\tau_{+}^{0},\tau_{+}^{1},\tau_{+}^{2},\tau_{+}^{3})$ 换成 $(0,1,2,3)$，把 $(\tau_{-}^{0},\tau_{-}^{1},\tau_{-}^{2},\tau_{-}^{3})$ 换成 $(n,x,y,z)$。例如，$\text{Z2A}\tau_{+}^{1}\tau_{-}^{0}\tau^{12}\tau_{-}^{3}$ 可缩写为 Z2A1n(12)z。

严格说来，这 272 个不同的 $Z_2$ PSG 是所谓的代数 PSG（algebraic PSG）。代数 PSG 定义为对称性群的扩张。代数 PSG 不同于不变 PSG（invariant PSG），后者定义为保持一个拟设 $u_{\pmb{ij}}$ 不变的所有变换的集合。虽然不变 PSG 必定是代数 PSG，但代数 PSG 未必是不变 PSG。这是因为某些代数 PSG 具有如下性质：在代数 PSG $G_1$ 下不变的\emph{任何}拟设 $u_{\pmb{ij}}$，实际上可能在一个更大的 PSG $G_2$ 下也不变。这时，原来的 PSG $G_1$ 不可能是该拟设的不变 PSG，拟设的不变 PSG 实际上由更大的 PSG $G_2$ 给出。如果我们把自己限制在通过拟设 $u_{\pmb{ij}}$ 构造的自旋液体上，那么就应当舍弃那些不是不变 PSG 的代数 PSG，因为这些代数 PSG 并不刻画平均场自旋液体。

\begin{center}
\small 表 9.1\quad $Z_2$ PSG 的 $g_{P_{xy}}$、$g_{P_x}$、$g_{P_y}$ 与 $g_T$ 的 17 种选择。这里 $\tau^{ab}=(\tau^{a}+\tau^{b})/\sqrt{2}$，$\tau^{a\bar{b}}=(\tau^{a}-\tau^{b})/\sqrt{2}$。\\[6pt]
\begin{tabular}{cccccccc}
\toprule
$g_{P_{xy}}$ & $g_{P_x}$ & $g_{P_y}$ & $g_T$ & $g_{P_{xy}}$ & $g_{P_x}$ & $g_{P_y}$ & $g_T$ \\
\midrule
$\tau^0$ & $\tau^0$ & $\tau^0$ & $\tau^0$ & $\tau^0$ & $\mathrm{i}\tau^3$ & $\mathrm{i}\tau^3$ & $\tau^0$ \\
$\mathrm{i}\tau^3$ & $\tau^0$ & $\tau^0$ & $\tau^0$ & $\mathrm{i}\tau^3$ & $\mathrm{i}\tau^3$ & $\mathrm{i}\tau^3$ & $\tau^0$ \\
$\mathrm{i}\tau^3$ & $\mathrm{i}\tau^1$ & $\mathrm{i}\tau^1$ & $\tau^0$ & $\tau^0$ & $\tau^0$ & $\tau^0$ & $\mathrm{i}\tau^3$ \\
$\tau^0$ & $\mathrm{i}\tau^3$ & $\mathrm{i}\tau^3$ & $\mathrm{i}\tau^3$ & $\tau^0$ & $\mathrm{i}\tau^1$ & $\mathrm{i}\tau^1$ & $\mathrm{i}\tau^3$ \\
$\mathrm{i}\tau^3$ & $\tau^0$ & $\tau^0$ & $\mathrm{i}\tau^3$ & $\mathrm{i}\tau^3$ & $\mathrm{i}\tau^3$ & $\mathrm{i}\tau^3$ & $\mathrm{i}\tau^3$ \\
$\mathrm{i}\tau^3$ & $\mathrm{i}\tau^1$ & $\mathrm{i}\tau^1$ & $\mathrm{i}\tau^3$ & $\mathrm{i}\tau^1$ & $\tau^0$ & $\tau^0$ & $\mathrm{i}\tau^3$ \\
$\mathrm{i}\tau^1$ & $\mathrm{i}\tau^3$ & $\mathrm{i}\tau^3$ & $\mathrm{i}\tau^3$ & $\mathrm{i}\tau^1$ & $\mathrm{i}\tau^1$ & $\mathrm{i}\tau^1$ & $\mathrm{i}\tau^3$ \\
$\mathrm{i}\tau^1$ & $\mathrm{i}\tau^2$ & $\mathrm{i}\tau^2$ & $\mathrm{i}\tau^3$ & $\mathrm{i}\tau^{12}$ & $\mathrm{i}\tau^1$ & $\mathrm{i}\tau^2$ & $\mathrm{i}\tau^0$ \\
$\mathrm{i}\tau^{12}$ & $\mathrm{i}\tau^1$ & $\mathrm{i}\tau^2$ & $\mathrm{i}\tau^3$ & & & & \\
\bottomrule
\end{tabular}
\end{center}

我们发现，在 272 个代数 $Z_2$ PSG 中，至少有 76 个不是不变 PSG。因此，272 个代数 $Z_2$ PSG 至多能给出 196 种可能的 $Z_2$ 自旋液体。

\subsection{一个 $Z_2$ 与一个 $U(1)$ 投影对称群及其拟设}

\begin{keypoints}
\item 对给定的 PSG，我们可以找出其最普遍的拟设。这些拟设共有的物理性质，就是与该 PSG 相联系的普适性质。
\end{keypoints}

作为 $Z_2$ PSG 的一个应用，让我们构造在 Z2A0013 PSG 下不变的最普遍拟设：
\begin{equation*}
\begin{array}{rcl}
G_x(\pmb{i})=\tau^0, & & G_y(\pmb{i})=\tau^0\\[2pt]
G_{P_x}(\pmb{i})=\tau^0, & & G_{P_y}(\pmb{i})=\tau^0\\[2pt]
G_{P_{xy}}(\pmb{i})=\mathrm{i}\tau^1, & & G_T(\pmb{i})=\mathrm{i}\tau^3
\end{array}
\tag{9.4.18}
\end{equation*}
由 $G_xT_x$ 与 $G_yT_y$ 下的不变性，再考虑到 $G_x$ 与 $G_y$ 是平庸的，Z2A0013 拟设直接就是平移不变的，形式为 $u_{\pmb{i},\pmb{i}+\pmb{m}}=u_{\pmb{m}}^{\mu}\tau^{\mu}$。求出最普遍 Z2A0013 拟设的关键是

% ---- 块界备注：上句延续至书页 406（p421）首行 "find out which $u_m^\mu$ must vanish."
% ---- （"……找出哪些 $u_{\pmb{m}}^{\mu}$ 必须为零。"），请 ch9_e 以 %JOIN% 续接，勿重复勿遗漏。

% ---------------- 新增术语（ch9_d） ----------------
% | extension (of a group) | 扩张 |
% | projective representation | 投影表示 |
% | normal subgroup | 正规子群 |
% | invariant group | 不变群 |
% | symmetry group (SG) | 对称性群（SG） |
% | invariant PSG | 不变 PSG |
% | algebraic PSG | 代数 PSG |
% | uniform gauge | 均匀规范 |
% | magnetic translational algebra | 磁平移代数（ch7_a 已收“磁平移”） |
% | sector | 扇区 |
% | pure gauge transformation | 纯规范变换 |
% | collinearness | 共线性 |
% | thermodynamic limit | 热力学极限 |
% | flux conservation | 通量守恒 |
% | infra-red divergence | 红外发散 |
% | quantum phase | 量子相 |
% | universal property | 普适性质 |
% | topologically-degenerate ground state | 拓扑简并基态 |
% | first-order mean-field theory | 一阶平均场理论 |
% | Z2-gapped / Z2-gapless / Z2-linear spin liquid | Z2 有能隙 / 无能隙 / 线性自旋液体 |
% | SU(2)-gapless / SU(2)-linear state | SU(2) 无能隙态 / SU(2) 线性态 |
% | U(1)-linear state | U(1) 线性态 |
% | invariant gauge group (IGG) | 不变规范群（IGG）（术语表已收，本块首次正文出现） |
% | compactified lattice | 紧致化格点（compactification 紧致化已收） |

% 块边界判定：
%   起点——书页 406（p421）顶部为承接书页 405 的半句 "find out which $u_m^\mu$ must
%     vanish."（ch9_d 末句停在"……的关键是"），故本块正文首行以 %JOIN% 续接，
%     随后接本页其余正文。
%   终点——书页 418（p433）末段止于完整句 "The gapless spinons appear only at
%     $(\pi/2,\pm\pi/2)$."；书页 419（p434）顶部另起新段 "Knowing the translational
%     symmetry ..."（连同其后的式 (9.6.9)、9.6.2 节与问题 9.6.1–9.6.3），均属下一块。
%     本块自然结束，不设 TAIL 区。
% 蓝本问题记录（均已按图改正）：
%   1) 书页 406 蓝本作 $u_{ij}=\mathrm{i}u_{ij}^0\tau^0+w_{ij}^3\tau^3$，图上为
%      $u_{ij}^3\tau^3$（蓝本 OCR 之误）。
%   2) 式 (9.4.23)、(9.4.24) 蓝本行尾残留 "\tag {9}" 杂项，已删。
%   3) 式 (9.4.28) 蓝本把最后一个规范变换误作 $G_\Gamma$，图上为 $G_T$。
%   4) 9.4.5 节 U1A 标记行蓝本残损（"U1A$_{a_{\eta_x px}}$b$_{\eta_y px}$cd$_y$,"
%      与 "$G_{F_{xy}}$"、"If the corresponding G contains a $\tau^1$, and equal
%      to $\tau^0$ otherwise"），图上为 "U1A$a_{\eta_{xpx}}b_{\eta_{ypx}}cd_{\eta_t}$、
%      $G_{P_{xy}}$"，句子为 "They are equal to $\tau^1$ if the corresponding $G$
%      contains a $\tau^1$, and equal to $\tau^0$ otherwise."，已按图译出。
%   5) 式 (9.6.8) 第二行右列图上印作 $u_{\pmb{i},\pmb{i}+-\pmb{x}+2\pmb{y}}$（多印一
%      个加号），按本意译为 $u_{\pmb{i},\pmb{i}-\pmb{x}+2\pmb{y}}$。
%   6) 书页 418 的 $\Gamma$ 矩阵定义，蓝本排成两栏且重复 $\Gamma_4$；图上为两行
%      （$\Gamma_0,\Gamma_1,\Gamma_2$ / $\Gamma_3,\Gamma_4$），已按图重排。
%   7) 书页 411–412 "contruct"、"we would like to the study"、书页 406–407 跨页
%      "under under"、书页 415 "(9.2.34"（缺右括号）等为原书排印错误，按本意译出。

%JOIN%找出哪些 $u_{\pmb{m}}^{\mu}$ 必须为零。时间反演变换 $G_{T}T$ 与 $180^{\circ}$ 旋转 $G_{P_{x}}P_{x}G_{P_{y}}P_{y}$ 在这里起关键作用。$G_{T}T$ 下的不变性要求 $-u_{\pmb{m}}=\tau^{3}u_{\pmb{m}}\tau^{3}$，因此只有 $u_{\pmb{m}}^{1,2}$ 可以非零。$G_{P_{x}}P_{x}G_{P_{y}}P_{y}$ 下的不变性要求 $u_{\pmb{m}}=u_{-\pmb{m}}$，而由于 $u_{\pmb{m}}^{1,2}$ 为实数且 $u_{-\pmb{m}}=u_{\pmb{m}}^{\dagger}$，$u_{\pmb{m}}=u_{\pmb{m}}^{1}\tau^{1}+u_{\pmb{m}}^{2}\tau^{2}$ 总能满足这一要求，所以 $180^{\circ}$ 旋转不附加任何新条件。其余变换则把不同的 $u_{\pmb{m}}^{1,2}$ 联系起来。$G_{P_{xy}}P_{xy}$ 下的不变性要求 $u_{P_{xy}\pmb{m}}=\tau^{1}u_{\pmb{m}}\tau^{1}$，其中 $P_{xy}\pmb{m}=P_{xy}(m_{x},m_{y})=(m_{y},m_{x})$。于是，最普遍的 Z2A0013 拟设具有如下形式：
\begin{equation*}
\begin{array}{l}
u_{\pmb{m}}=u_{\pmb{m}}^{1}\tau^{1}+u_{\pmb{m}}^{2}\tau^{2}\\[2pt]
u_{P_{x}\pmb{m}}^{1,2}=u_{\pmb{m}}^{1,2},\qquad u_{P_{y}\pmb{m}}^{1,2}=u_{\pmb{m}}^{1,2}\\[2pt]
u_{P_{xy}\pmb{m}}^{1}=u_{\pmb{m}}^{1},\qquad u_{P_{xy}\pmb{m}}^{2}=-u_{\pmb{m}}^{2}
\end{array}
\end{equation*}
把式 (9.2.41) 中的 $(\tau^{3},\tau^{1})$ 换成 $(\tau^{1},\tau^{2})$，并把 $a_{0}$ 认同于 $u_{\pmb{m}=0}$，我们发现 $Z_{2}$ 线性自旋液体的拟设 (9.2.41) 是上述拟设的一个特例。因此，拟设 (9.2.41) 的 PSG 就是 Z2A0013 PSG，相应的自旋液体是一个 Z2A0013 自旋液体。

除了上面研究的 $Z_{2}$ 对称自旋液体之外，还可能存在低能规范群为 $U(1)$ 或 $SU(2)$ 的对称自旋液体。这类 $U(1)$ 对称或 $SU(2)$ 对称的自旋液体由 IGG 为 $U(1)$ 或 $SU(2)$ 的 PSG 分类。我们把这些 PSG 称为 $U(1)$ PSG 与 $SU(2)$ PSG。$U(1)$ PSG 与 $SU(2)$ PSG 已由 Wen (2002c) 算出，9.4.5 节将总结这些结果。

$U(1)$ PSG 中的一个由下式给出：
\begin{equation*}
\begin{array}{rcl}
G_{x}=g_{3}(\theta_{x})\mathrm{i}\tau^{1}, & & G_{y}=g_{3}(\theta_{y})\mathrm{i}\tau^{1},\\[2pt]
G_{P_{x}}=(-)^{i_{x}}g_{3}(\theta_{px}), & & G_{P_{y}}=(-)^{i_{y}}g_{3}(\theta_{py}),\\[2pt]
G_{P_{xy}}=g_{3}(\theta_{pxy})\mathrm{i}\tau^{1}, & & G_{T}=(-)^{\pmb{i}}g_{3}(\theta_{t}).
\end{array}
\tag{9.4.19}
\end{equation*}
其中
\begin{equation*}
g_{a}(\theta)\equiv\mathrm{e}^{\mathrm{i}\theta\tau^{a}}.
\end{equation*}
IGG 由 $\{g_{3}(\phi)\,|\,\phi\in[0,2\pi)\}$ 构成。这样的 PSG 标记为 U1Cn01n。在式 (9.4.27) 中取 $\eta_{xpx}=-1$、$\eta_{ypx}=+1$、$\eta_{pxy}=+1$ 与 $\eta_{t}=-1$ 即得到它。

让我们求出在 U1Cn01n PSG 下不变的最普遍拟设。要在 IGG 下不变，拟设必须对任意 $\theta$ 满足 $g_{3}(\theta)u_{\pmb{ij}}g_{3}^{\dagger}(\theta)=u_{\pmb{ij}}$，因此 $u_{\pmb{ij}}$ 必须具有 $u_{\pmb{ij}}=\mathrm{i}u_{\pmb{ij}}^{0}\tau^{0}+u_{\pmb{ij}}^{3}\tau^{3}$ 的形式。在平移变换 $G_{x}T_{x}$ 与 $G_{y}T_{y}$ 下的不变性要求 $u_{\pmb{ij}}^{0}$ 平移不变，而 $u_{\pmb{ij}}^{3}$ 在两个平移 $T_{x,y}$ 下变号。于是 $u_{\pmb{ij}}$ 具有形式 $u_{\pmb{i},\pmb{i}+\pmb{m}}=\mathrm{i}u_{\pmb{m}}^{0}\tau^{0}+(-)^{i}u_{\pmb{m}}^{3}\tau^{3}$。时间反演变换 $G_{T}T$ 下的不变性要求 $-u_{\pmb{m}}=(-)^{\pmb{m}}u_{\pmb{m}}$。我们发现，若 $\pmb{m}$ 为偶，则 $u_{\pmb{m}}^{0,3}=0$。\footnote{这里 $\pmb{m}$ 为偶（奇），若 $m_{x}+m_{y}$ 为偶（奇）。}$180^{\circ}$ 旋转 $G_{P_{x}}P_{x}G_{P_{y}}P_{y}$ 下的不变性要求 $u_{-\pmb{i},-\pmb{i}-\pmb{m}}=u_{-\pmb{i}-\pmb{m},-\pmb{i}}^{\dagger}=(-)^{\pmb{m}}u_{\pmb{i},\pmb{i}+\pmb{m}}$，这意味着 $-\mathrm{i}u_{\pmb{m}}^{0}\tau^{0}+(-)^{\pmb{i}+\pmb{m}}u_{\pmb{m}}^{3}\tau^{3}=(-)^{\pmb{m}}\big(\mathrm{i}u_{\pmb{m}}^{0}\tau^{0}+(-)^{\pmb{i}}u_{\pmb{m}}^{3}\tau^{3}\big)$。由此得到 $u_{\pmb{m}}^{0}=-(-)^{\pmb{m}}u_{\pmb{m}}^{0}$ 与 $u_{\pmb{m}}^{3}=u_{\pmb{m}}^{3}$，对 $u_{\pmb{m}}^{0,3}$ 没有给出新的约束。厄米关系 $u_{\pmb{i},\pmb{i}+\pmb{m}}^{\dagger}=u_{\pmb{i}+\pmb{m},\pmb{i}}$ 给出 $u_{\pmb{m}}^{0}=-u_{-\pmb{m}}^{0}$ 与 $u_{\pmb{m}}^{3}=(-)^{\pmb{m}}u_{-\pmb{m}}^{3}$。汇总上述条件，并计入 $G_{P_{x}}P_{x}$ 等的作用，我们发现最普遍的 U1Cn01n 拟设具有如下形式：
\begin{equation*}
\begin{array}{l}
u_{\pmb{i},\pmb{i}+\pmb{m}}=\mathrm{i}u_{\pmb{m}}^{0}\tau^{0}+(-)^{i}u_{\pmb{m}}^{3}\tau^{3}\\[2pt]
\qquad u_{\pmb{m}}^{0,3}=0,\ \text{若 }\pmb{m}\text{ 为偶},\qquad u_{\pmb{m}}^{0,3}=-u_{-\pmb{m}}^{0,3}\\[2pt]
\qquad u_{P_{x}\pmb{m}}^{0,3}=(-)^{m_{x}}u_{\pmb{m}}^{0,3},\qquad u_{P_{y}\pmb{m}}^{0,3}=(-)^{m_{y}}u_{\pmb{m}}^{0,3}\\[2pt]
\qquad u_{P_{xy}\pmb{m}}^{0}=u_{\pmb{m}}^{0},\qquad u_{P_{xy}\pmb{m}}^{3}=-u_{\pmb{m}}^{3}
\end{array}
\end{equation*}
$U(1)$ 线性拟设（交错通量相）(9.2.34) 是上述拟设的一个特例。因此，交错通量相 (9.2.34) 是一个 U1Cn01n 态。

我们可以用规范变换 $W_{i}=(\mathrm{i}\tau^{1})^{i}$ 把上述拟设变成平移不变的：
\begin{equation*}
\begin{array}{rcl}
u_{\pmb{i},\pmb{i}+\pmb{m}}=\tilde{u}_{\pmb{m}}^{1}\tau^{1}+\tilde{u}_{\pmb{m}}^{2}\tau^{2} & & \\[2pt]
\tilde{u}_{\pmb{m}}^{1,2}=0,\quad \text{若 }\pmb{m}\text{ 为偶}, & & \\[2pt]
G_{x}(\pmb{i})=g_{3}((-)^{\pmb{i}}\theta_{x}), & & G_{y}(\pmb{i})=g_{3}((-)^{\pmb{i}}\theta_{y}),\\[2pt]
G_{P_{x}}(\pmb{i})=g_{3}((-)^{\pmb{i}}\theta_{px}), & & G_{P_{y}}(\pmb{i})=g_{3}((-)^{\pmb{i}}\theta_{py}),\\[2pt]
G_{P_{xy}}(\pmb{i})=\mathrm{i}g_{3}((-)^{\pmb{i}}\theta_{pxy})\tau^{1}, & & G_{T}(\pmb{i})=(-)^{\pmb{i}}g_{3}((-)^{\pmb{i}}\theta_{t});
\end{array}
\tag{9.4.20}
\end{equation*}
其中我们还列出了变换后 PSG 中的规范变换。IGG 现在具有形式 $IGG=\{g_{3}((-)^{\pmb{i}}\phi)\,|\,\phi\in[0,2\pi)\}$。如果把 $(\tau^{1},\tau^{2})$ 换成 $(\tau^{3},\tau^{1})$，则可以证明上述拟设就 $f$ 费米子而言对应一个 $d$ 波态（见式 (9.2.4) 与 (9.2.30)）。

拟设 (9.4.20) 的自旋子谱由下式给出：
\begin{equation*}
E_{\pm}(\pmb{k})=\pm\sqrt{\Bigg[\sum_{\pmb{m}=\text{奇}}\tilde{u}_{\pmb{m}}^{1}\cos(\pmb{m}\cdot\pmb{k})\Bigg]^{2}+\Bigg[\sum_{\pmb{m}=\text{奇}}\tilde{u}_{\pmb{m}}^{2}\cos(\pmb{m}\cdot\pmb{k})\Bigg]^{2}}
\end{equation*}
我们注意到，对奇的 $\pmb{m}$，$\cos(\pmb{m}\cdot\pmb{k})$ 在 $\pmb{k}=(\pm\frac{\pi}{2},\pm\frac{\pi}{2})$ 处为零，因此 $E_{\pm}(\pmb{k})$ 在 $\pmb{k}=(\pm\frac{\pi}{2},\pm\frac{\pi}{2})$ 处为零。U1Cn01n 自旋液体态在 $\pmb{k}=(\pm\frac{\pi}{2},\pm\frac{\pi}{2})$ 处总有至少四支无能隙自旋子。无能隙自旋子及其晶体动量 $(\pm\frac{\pi}{2},\pm\frac{\pi}{2})$ 是 U1Cn01n 自旋液体态的普适性质。自旋子的无能隙性及其晶体动量受 U1Cn01n PSG 保护。这是一个惊人的结果。作为一个纯玻色模型，U1Cn01n 自旋液体态不破缺任何对称性，然而它却包含一种特殊的量子序，保证了无能隙\emph{费米子}的存在！由于 $\pmb{k}=(\pm\frac{\pi}{2},\pm\frac{\pi}{2})$ 处的无能隙自旋子是 U1Cn01n 态的普适性质，我们可以利用它在实验上探测 U1Cn01n 量子序。

\subsection*{问题 9.4.1}

试通过式 (9.4.3) 找出把 $G_{y}(\pmb{i})=\Theta(\pmb{i})\tau^{0}$ 变为 $G_{y}(\pmb{i})=\tau^{0}$ 的规范变换。这里 $\Theta(\pmb{i})$ 是只取 $\pm1$ 两个值的任意函数。

\subsection*{问题 9.4.2}

\begin{enumerate}
\item 从式 (9.4.14) 求出 Z2Azz13 PSG 的规范变换 $G_{x,y}$、$G_{P_{x},P_{y},P_{xy}}$ 与 $G_{T}$。
\item 求出在 Z2Azz13 PSG 下不变的最普遍拟设。我们将在后面证明，这样的拟设总有无能隙费米子激发。
\end{enumerate}

\subsection*{问题 9.4.3}

证明式 (9.4.20)。（提示：见式 (9.4.3)。）

\subsection{评注：对称 $U(1)$ 与 $SU(2)$ 自旋液体的分类}

\begin{keypoints}
\item 对所有平均场拟设的分类：这些拟设的 PSG 是格点对称性群的 $U(1)$ 或 $SU(2)$ 扩张，即 $PSG/U(1)=SG$ 或 $PSG/SU(2)=SG$。
\end{keypoints}

Wen (2002c) 给出了 $U(1)$ PSG 与 $SU(2)$ PSG 的分类。下面我们只总结结果。刻画平均场对称 $U(1)$ 自旋液体的 PSG 可以分成 U1A、U1B、U1C 与 U1$_{n}^{m}$ 四种类型。有下列 24 个 U1A 型 PSG：
\begin{equation*}
\begin{array}{rcl}
G_{x}=g_{3}(\theta_{x}), & & G_{y}=g_{3}(\theta_{y}),\\[2pt]
G_{P_{x}}=\eta_{ypx}^{i_{y}}g_{3}(\theta_{px}), & & G_{P_{y}}=\eta_{ypx}^{i_{x}}g_{3}(\theta_{py})\\[2pt]
G_{P_{xy}}=g_{3}(\theta_{pxy}),\ g_{3}(\theta_{pxy})\mathrm{i}\tau^{1}, & & G_{T}=\eta_{t}^{\pmb{i}}g_{3}(\theta_{t})\big|_{\eta_{t}=-1},\ \eta_{t}^{\pmb{i}}g_{3}(\theta_{t})\mathrm{i}\tau^{1}
\end{array}
\tag{9.4.21}
\end{equation*}
以及
\begin{equation*}
\begin{array}{rcl}
G_{x}=g_{3}(\theta_{x}), & & G_{y}=g_{3}(\theta_{y}),\\[2pt]
G_{P_{x}}=\eta_{xpx}^{i_{x}}g_{3}(\theta_{px})\mathrm{i}\tau^{1}, & & G_{P_{y}}=\eta_{xpx}^{i_{y}}g_{3}(\theta_{py})\mathrm{i}\tau^{1}\\[2pt]
G_{P_{xy}}=g_{3}(\theta_{pxy}),\ g_{3}(\theta_{pxy})\mathrm{i}\tau^{1}, & & G_{T}=\eta_{t}^{\pmb{i}}g_{3}(\theta_{t})\big|_{\eta_{t}=-1},\ \eta_{t}^{\pmb{i}}g_{3}(\theta_{t})\mathrm{i}\tau^{1}
\end{array}
\tag{9.4.22}
\end{equation*}
我们将用 $\text{U1A}a_{\eta_{xpx}}b_{\eta_{ypx}}cd_{\eta_{t}}$ 来标记这 24 个 PSG。这里 $a$、$b$、$c$ 与 $d$ 分别与 $G_{P_{x}}$、$G_{P_{y}}$、$G_{P_{xy}}$ 及 $G_{T}$ 相联系：若相应的 $G$ 包含一个 $\tau^{1}$，它们就取 $\tau^{1}$，否则取 $\tau^{0}$。一个典型的记号形如 $\text{U1A}\tau_{-}^{1}\tau^{1}\tau^{0}\tau_{-}^{1}$，可缩写为 $\text{U1A}x10x$。

还有下列 24 个 U1B 型 PSG：
\begin{equation*}
\begin{array}{rcl}
G_{x}=(-)^{i_{y}}g_{3}(\theta_{x}), & & G_{y}=g_{3}(\theta_{y}),\\[2pt]
G_{P_{x}}=\eta_{ypx}^{i_{y}}g_{3}(\theta_{px}), & & G_{P_{y}}=\eta_{ypx}^{i_{x}}g_{3}(\theta_{py})\\[2pt]
(-)^{i_{x}i_{y}}G_{P_{xy}}=g_{3}(\theta_{pxy}),\ g_{3}(\theta_{pxy})\mathrm{i}\tau^{1} & & G_{T}=\eta_{t}^{\pmb{i}}g_{3}(\theta_{t})\big|_{\eta_{t}=-1},\ \eta_{t}^{\pmb{i}}g_{3}(\theta_{t})\mathrm{i}\tau^{1}
\end{array}
\tag{9.4.23}
\end{equation*}
以及
\begin{equation*}
\begin{array}{rcl}
G_{x}=(-)^{i_{y}}g_{3}(\theta_{x}), & & G_{y}=g_{3}(\theta_{y}),\\[2pt]
G_{P_{x}}=\eta_{xpx}^{i_{x}}g_{3}(\theta_{px})\mathrm{i}\tau^{1}, & & G_{P_{y}}=\eta_{xpx}^{i_{y}}g_{3}(\theta_{py})\mathrm{i}\tau^{1},\\[2pt]
(-)^{i_{x}i_{y}}G_{P_{xy}}=g_{3}(\theta_{pxy}),\ g_{3}(\theta_{pxy})\mathrm{i}\tau^{1} & & G_{T}=\eta_{t}^{\pmb{i}}g_{3}(\theta_{t})\big|_{\eta_{t}=-1},\ \eta_{t}^{\pmb{i}}g_{3}(\theta_{t})\mathrm{i}\tau^{1}
\end{array}
\tag{9.4.24}
\end{equation*}
我们将用 $\text{U1B}a_{\eta_{xpx}}b_{\eta_{ypx}}cd_{\eta_{t}}$ 来标记这 24 个 PSG。

60 个 U1C 型 PSG 由下式给出：
\begin{equation*}
\begin{array}{rcl}
G_{x}=g_{3}(\theta_{x})\mathrm{i}\tau^{1}, & & G_{y}=g_{3}(\theta_{y})\mathrm{i}\tau^{1},\\[2pt]
G_{P_{x}}=\eta_{xpx}^{i_{x}}\eta_{ypx}^{i_{y}}g_{3}(\theta_{px}), & & G_{P_{y}}=\eta_{ypx}^{i_{x}}\eta_{xpx}^{i_{y}}g_{3}(\theta_{py})\\[2pt]
G_{P_{xy}}=\eta_{pxy}^{i_{x}}g_{3}\big(\eta_{pxy}^{\pmb{i}}\frac{\pi}{4}+\theta_{pxy}\big), & & G_{T}=\eta_{t}^{\pmb{i}}g_{3}(\theta_{t})\big|_{\eta_{t}=-1},\ \eta_{pxy}^{i_{x}}g_{3}(\theta_{t})\mathrm{i}\tau^{1}
\end{array}
\tag{9.4.25}
\end{equation*}
\begin{equation*}
\begin{array}{rcl}
G_{x}=g_{3}(\theta_{x})\mathrm{i}\tau^{1}, & & G_{y}=g_{3}(\theta_{y})\mathrm{i}\tau^{1},\\[2pt]
G_{P_{x}}=\eta_{xpx}^{i_{x}}g_{3}(\theta_{px})\mathrm{i}\tau^{1}, & & G_{P_{y}}=\eta_{xpx}^{i_{y}}\eta_{pxy}^{\pmb{i}}g_{3}(\theta_{py})\mathrm{i}\tau^{1},\\[2pt]
G_{P_{xy}}=\eta_{pxy}^{i_{x}}g_{3}\big(\eta_{pxy}^{\pmb{i}}\frac{\pi}{4}+\theta_{pxy}\big), & & G_{T}=\eta_{t}^{\pmb{i}}g_{3}(\theta_{t})\big|_{\eta_{t}=-1},\ \eta_{pxy}^{i_{x}}\eta_{t}^{\pmb{i}}g_{3}(\theta_{t})\mathrm{i}\tau^{1}
\end{array}
\tag{9.4.26}
\end{equation*}
\begin{equation*}
\begin{array}{rcl}
G_{x}=g_{3}(\theta_{x})\mathrm{i}\tau^{1}, & & G_{y}=g_{3}(\theta_{y})\mathrm{i}\tau^{1},\\[2pt]
G_{P_{x}}=\eta_{xpx}^{i_{x}}\eta_{ypx}^{i_{y}}g_{3}(\theta_{px}), & & G_{P_{y}}=\eta_{ypx}^{i_{x}}\eta_{xpx}^{i_{y}}g_{3}(\theta_{py}),\\[2pt]
G_{P_{xy}}=g_{3}(\theta_{pxy})\mathrm{i}\tau^{1}, & & G_{T}=\eta_{t}^{\pmb{i}}g_{3}(\theta_{t})\big|_{\eta_{t}=-1}
\end{array}
\tag{9.4.27}
\end{equation*}
\begin{equation*}
\begin{array}{rcl}
G_{x}=g_{3}(\theta_{x})\mathrm{i}\tau^{1}, & & G_{y}=g_{3}(\theta_{y})\mathrm{i}\tau^{1},\\[2pt]
G_{P_{x}}=\eta_{xpx}^{i_{x}}\eta_{ypx}^{i_{y}}g_{3}(\theta_{px}), & & G_{P_{y}}=\eta_{ypx}^{i_{x}}\eta_{xpx}^{i_{y}}g_{3}(\theta_{py}),\\[2pt]
G_{P_{xy}}=g_{3}\big(\eta_{pxy}^{\pmb{i}}\frac{\pi}{4}+\theta_{pxy}\big)\mathrm{i}\tau^{1}, & & G_{T}=\eta_{pxy}^{i_{x}}\eta_{t}^{\pmb{i}}g_{3}(\theta_{t})\mathrm{i}\tau^{1}
\end{array}
\tag{9.4.28}
\end{equation*}
以及
\begin{equation*}
\begin{array}{rcl}
G_{x}=g_{3}(\theta_{x})\mathrm{i}\tau^{1}, & & G_{y}=g_{3}(\theta_{y})\mathrm{i}\tau^{1},\\[2pt]
G_{P_{x}}=\eta_{xpx}^{i_{x}}g_{3}(\theta_{px})\mathrm{i}\tau^{1}, & & G_{P_{y}}=\eta_{xpx}^{i_{y}}\eta_{pxy}^{\pmb{i}}g_{3}(\theta_{py})\mathrm{i}\tau^{1},\\[2pt]
G_{P_{xy}}=g_{3}\big(\eta_{pxy}^{\pmb{i}}\frac{\pi}{4}+\theta_{pxy}\big)\mathrm{i}\tau^{1}, & & G_{T}=\eta_{t}^{\pmb{i}}g_{3}(\theta_{t})\big|_{\eta_{t}=-1},\ \eta_{t}^{\pmb{i}}\eta_{pxy}^{i_{x}}g_{3}(\theta_{t})\mathrm{i}\tau^{1}.
\end{array}
\tag{9.4.29}
\end{equation*}
它们将用 $\text{U1C}a_{\eta_{xpx}}b_{\eta_{ypx}}c_{\eta_{pxy}}d_{\eta_{t}}$ 来标记。

U1$_{n}^{m}$ 型 PSG 尚未被分类。不过我们知道，对每一个有理数 $m/n\in(0,1)$，至少存在一个 U1$_{n}^{m}$ 型的平均场对称自旋液体，其拟设由下式给出：
\begin{equation*}
u_{\pmb{i},\pmb{i}+\pmb{x}}=\chi\tau^{3},\qquad u_{\pmb{i},\pmb{i}+\pmb{y}}=\chi g_{3}\big(\frac{m\pi}{n}i_{x}\big)\tau^{3}
\tag{9.4.30}
\end{equation*}
每个方格穿过 $\pi m/n$ 的通量。因此存在无穷多种 U1$_{n}^{m}$ 自旋液体。

我们想指出，上面 108 个 U1A、U1B 与 U1C PSG 都是代数 PSG，它们只是所有可能的代数 $U(1)$ PSG 的一个子集。不过，它们确实包含了 U1A、U1B 与 U1C 型的所有不变 $U(1)$ PSG。我们发现 108 个 PSG 中有 46 个同时也是不变 PSG，因此共有 46 种 U1A、U1B、U1C 型的平均场 $U(1)$ 自旋液体。它们的拟设与标记见 Wen (2002c)。

为了给对称 $SU(2)$ 自旋液体分类，我们发现共有 8 个不同的 $SU(2)$ PSG，由下式给出：
\begin{equation*}
\begin{array}{rcl}
G_{x}(\pmb{i})=g_{x}, & & G_{y}(\pmb{i})=g_{y},\\[2pt]
G_{P_{x}}(\pmb{i})=\eta_{xpx}^{i_{x}}\eta_{xpy}^{i_{y}}g_{P_{x}}, & & G_{P_{y}}(\pmb{i})=\eta_{xpy}^{i_{x}}\eta_{xpx}^{i_{y}}g_{P_{y}}\\[2pt]
G_{P_{xy}}(\pmb{i})=g_{P_{xy}}, & & G_{T}(\pmb{i})=(-)^{\pmb{i}}g_{T},
\end{array}
\tag{9.4.31}
\end{equation*}
以及
\begin{equation*}
\begin{array}{rcl}
G_{x}(\pmb{i})=(-)^{i_{y}}g_{x}, & & G_{y}(\pmb{i})=g_{y},\\[2pt]
G_{P_{x}}(\pmb{i})=\eta_{xpx}^{i_{x}}\eta_{xpy}^{i_{y}}g_{P_{x}}, & & G_{P_{y}}(\pmb{i})=\eta_{xpy}^{i_{x}}\eta_{xpx}^{i_{y}}g_{P_{y}}\\[2pt]
G_{P_{xy}}(\pmb{i})=(-)^{i_{x}i_{y}}g_{P_{xy}}, & & G_{T}(\pmb{i})=(-)^{\pmb{i}}g_{T}
\end{array}
\tag{9.4.32}
\end{equation*}
其中各个 $g$ 都是 $SU(2)$ 中的矩阵。我们引入如下记号：
\begin{equation*}
\begin{array}{l}
\text{SU2A}\tau_{\eta_{xpx}}^{0}\tau_{\eta_{xpy}}^{0}\\[2pt]
\text{SU2B}\tau_{\eta_{xpx}}^{0}\tau_{\eta_{xpy}}^{0}
\end{array}
\tag{9.4.33}
\end{equation*}
用来表示上述 8 个 PSG：$\text{SU2A}\tau_{\eta_{xpx}}^{0}\tau_{\eta_{xpy}}^{0}$ 对应式 (9.4.31)，$\text{SU2B}\tau_{\eta_{xpx}}^{0}\tau_{\eta_{xpy}}^{0}$ 对应式 (9.4.32)。我们发现，8 个 $SU(2)$ PSG 中只有 4 个——即 $\text{SU2A}[n0,0n]$ 与 $\text{SU2B}[n0,0n]$——能给出 $SU(2)$ 对称自旋液体。SU2A$n0$ 态就是均匀 RVB 态 (9.2.32)，SU2B$n0$ 态就是 $\pi$ 通量态 (9.2.33)。另外两个 $SU(2)$ 自旋液体一个以 SU2A$0n$ 标记，即
\begin{equation*}
\begin{array}{rcl}
u_{\pmb{i},\pmb{i}+2\pmb{x}+\pmb{y}}=+\mathrm{i}\chi\tau^{0}, & & u_{\pmb{i},\pmb{i}-2\pmb{x}+\pmb{y}}=-\mathrm{i}\chi\tau^{0},\\[2pt]
u_{\pmb{i},\pmb{i}+\pmb{x}+2\pmb{y}}=+\mathrm{i}\chi\tau^{0}, & & u_{\pmb{i},\pmb{i}-\pmb{x}+2\pmb{y}}=+\mathrm{i}\chi\tau^{0},
\end{array}
\tag{9.4.34}
\end{equation*}
另一个以 SU2B$0n$ 标记，即
\begin{equation*}
\begin{array}{rcl}
u_{\pmb{i},\pmb{i}+2\pmb{x}+\pmb{y}}=+\mathrm{i}(-)^{i_{x}}\chi\tau^{0}, & & u_{\pmb{i},\pmb{i}-2\pmb{x}+\pmb{y}}=-\mathrm{i}(-)^{i_{x}}\chi\tau^{0},\\[2pt]
u_{\pmb{i},\pmb{i}+\pmb{x}+2\pmb{y}}=+\mathrm{i}\chi\tau^{0}, & & u_{\pmb{i},\pmb{i}-\pmb{x}+2\pmb{y}}=+\mathrm{i}\chi\tau^{0}
\end{array}
\tag{9.4.35}
\end{equation*}

上述结果给出了平均场层面对 $U(1)$ 与 $SU(2)$ 对称自旋液体的分类。我们想指出，$U(1)$ 与 $SU(2)$ 平均场态并不是稳定的平均场态，其中一些是边缘的。因此，与 $Z_{2}$ 平均场态相比，$U(1)$ 与 $SU(2)$ 平均场态是否对应于自旋-1/2 模型某种物理的 $U(1)$ 或 $SU(2)$ 对称自旋液体，就更不清楚了。物理的 $U(1)$ 或 $SU(2)$ 态可能只存在于大 $N$ 极限中（见 9.8 节）。

\section{无对称性破缺的连续相变}

\begin{keypoints}
\item 量子相之间的连续相变由下述原理支配：令 $PSG_{1}$ 与 $PSG_{2}$ 为相变两侧两个量子相的 PSG，$PSG_{cr}$ 为描述量子临界态（quantum critical state）的 PSG，则 $PSG_{1}\subseteq PSG_{cr}$ 且 $PSG_{2}\subseteq PSG_{cr}$。
\item 上述原理既适用于对称性破缺相变，也适用于不改变任何对称性的相变。
\end{keypoints}

在对平均场对称自旋液体进行了分类之后，我们想知道这些对称自旋液体彼此之间有什么关系。特别地，我们想知道哪些自旋液体可以通过\emph{连续}相变相互转变。在平均场层面上，这个问题可以完全用 PSG 来解决。其想法是：PSG 正是平均场态的“对称性”群。平均场相变可以用 PSG 的改变来描述。正如对称性破缺相变一样，PSG 为 $PSG_{1}$ 的平均场态能够通过一个连续的（平均场）相变变为 PSG 为 $PSG_{2}$ 的平均场态，当且仅当 $PSG_{2}\subset PSG_{1}$ 或 $PSG_{1}\subset PSG_{2}$。

为了理解这一结果，让我们假定由 $PSG_{1}$ 描述的平均场态有一个拟设 $u_{\pmb{ij}}$，它的邻居有一个拟设 $u_{\pmb{ij}}+\delta u_{\pmb{ij}}$，这里 $\delta u_{\pmb{ij}}$ 是一个小微扰。设该微扰把 PSG 改变为另一个不同的 PSG——$PSG_{2}$。由于 $\delta u_{\pmb{ij}}$ 是强度未定的微扰，要使 $u_{\pmb{ij}}+\delta u_{\pmb{ij}}$ 在 $PSG_{2}$ 下不变，$u_{\pmb{ij}}$ 与 $\delta u_{\pmb{ij}}$ 就都必须在 $PSG_{2}$ 下不变。因此 $PSG_{2}\subset PSG_{1}$。

利用上述结果，我们可以按下列步骤求出某些重要对称自旋液体邻域中的对称自旋液体。我们从一个 PSG 为 $PSG_{1}$ 的对称自旋液体 $u_{\pmb{ij}}$ 出发，这里 $PSG_{1}$ 是对称性群 $SG$ 被 $IGG_{1}$ 的扩张，即 $PSG_{1}/IGG_{1}=SG$。然后我们找出 $PSG_{1}$ 的所有作为同一对称性群 $SG$ 的扩张的子群，把这些子群记为 $PSG_{2}$。那么 $PSG_{2}$ 必须有一个正规子群 $IGG_{2}$ 使得 $PSG_{2}/IGG_{2}=SG$。利用这些子群，我们可以构造出所有具有相同对称性的近邻平均场态。

经过 Wen (2002c) 中的冗长计算，人们求出了式 (9.2.34) 的 $U(1)$ 线性态 U1Cn01n、式 (9.2.32) 的 $SU(2)$ 无能隙态 SU2An0 以及式 (9.2.33) 的 $SU(2)$ 线性态 SU2Bn0 周围的所有平均场对称自旋液体。例如，已经证明，在平均场层面上，$U(1)$ 线性自旋液体 U1Cn01n 可以连续变为 8 个不同的 $Z_{2}$ 自旋液体：Z2A0013、Z2Azz13、Z2A001n、Z2Azz1n、Z2B0013、Z2Bzz13、Z2B001n 与 Z2Bzz1n。

让我们讨论一个简单的例子来演示上述结果。拟设
\begin{equation*}
\begin{array}{rcl}
u_{\pmb{i},\pmb{i}+\pmb{x}}=\chi\tau^{1}-\eta\tau^{2}, & & u_{\pmb{i},\pmb{i}+\pmb{y}}=\chi\tau^{1}+\eta\tau^{2},\\[2pt]
u_{\pmb{i},\pmb{i}+\pmb{x}+\pmb{y}}=-\gamma\tau^{1}, & & u_{\pmb{i},\pmb{i}-\pmb{x}+\pmb{y}}=+\gamma\tau^{1},\qquad a_{0}^{l}=0.
\end{array}
\tag{9.5.1}
\end{equation*}
在 $\chi\neq\eta$ 时，$\gamma=0$ 的情形描述 U1Cn01n 自旋液体（交错通量态）。也就是说，该拟设在式 (9.4.20) 所给规范变换之后跟随的对称性变换下不变。当 $\gamma\neq0$ 时，拟设生成非共线的 $SU(2)$ 通量，把 $U(1)$ 规范结构破缺到 $Z_{2}$ 规范结构，拟设此时描述一个 $Z_{2}$ 态。对非零的 $\gamma$，拟设不再在式 (9.4.20) 所列 U1Cn01n 规范变换之后跟随的对称性变换下不变，而只在 U1Cn01n 规范变换的一个子集之后跟随的对称性变换下不变：
\begin{equation*}
\begin{array}{rcl}
G_{x}(\pmb{i})=g_{3}(0)=\tau^{0}, & & G_{y}(\pmb{i})=g_{3}(0)=\tau^{0},\\[2pt]
G_{P_{x}}(\pmb{i})=g_{3}((-)^{\pmb{i}}\pi/2)=\mathrm{i}(-)^{\pmb{i}}\tau^{3}, & & G_{P_{y}}(\pmb{i})=g_{3}((-)^{\pmb{i}}\pi/2)=\mathrm{i}(-)^{\pmb{i}}\tau^{3},\\[2pt]
G_{P_{xy}}(\pmb{i})=\mathrm{i}g_{3}(0)\tau^{1}=\mathrm{i}\tau^{1}, & & G_{T}(\pmb{i})=(-)^{\pmb{i}}g_{3}((-)^{\pmb{i}}\pi/2)=\mathrm{i}\tau^{3}.
\end{array}
\end{equation*}
上述规范变换之后跟随的对称性变换给出 Z2Azz13 PSG。于是，当 $\gamma$ 从零变为非零时，U1Cn01n 态连续地变为 Z2Azz13 态。

我们想强调，上述关于连续相变的结果只在平均场层面上有效。有些平均场结果能在量子涨落下保存下来，有些则不能。要知道哪些平均场结果在平均场理论之外仍然有效，需要逐例研究（Mudry and Fradkin, 1994）。

量子涨落的一个可能效果是使某些平均场态失稳。我们可以假定这种失稳来自量子涨落产生的某些相关微扰。在这一假定下，一旦计入量子涨落，一些平均场的稳定不动点对真实物理系统而言就变成了不稳定不动点。于是，平均场相图中的一些相可能收缩成临界线（critical line），这些线代表两个稳定量子相之间相变点上的临界态（critical state）（见图 9.7）。这一图像引出了本节开头要点中所述的关于支配连续量子相变的原理的猜想。

\begin{figure}[H]
\centering
\includegraphics[width=0.72\textwidth]{fig_9.7.pdf}
\caption*{图 9.7\quad (a) 由 $PSG_{1,2,3}$ 刻画三个相的平均场相图，它们之间的平均场相变是连续的。PSG 满足 $PSG_{1}\subset PSG_{3}$ 与 $PSG_{2}\subset PSG_{3}$。$PSG_{1,2,3}$ 的一种可能取法是 $PSG_{1}=$Z2A0013、$PSG_{2}=$Z2Azz13、$PSG_{3}=$U1Cn01n。(b) 计入量子涨落后相应的物理相图。这里我们假定，计入量子涨落之后，态 $PSG_{3}$ 变成一个只带一个相关微扰的不稳定不动点。平均场态 $PSG_{3}$ 收缩成一条临界线，它描述两个态 $PSG_{1}$ 与 $PSG_{2}$ 之间的相变。交错通量态 U1Cn01n 如果不稳定，可能作为一个临界态出现。}
\end{figure}

我们想强调，上述所有自旋液体都具有相同的对称性。因此，它们之间的连续相变（如果存在的话）代表了一类新的量子连续相变——不改变任何对称性的相变（Wegner, 1971; Kosterlitz and Thouless, 1973; Dasgupta and Halperin, 1981; Wen and Wu, 1993; Chen \emph{et al.}, 1993; Senthil \emph{et al.}, 1999; Read and Green, 2000; Wen, 2000）。

\subsection*{问题 9.5.1}

拟设
\begin{equation*}
\begin{array}{rcl}
u_{\pmb{i},\pmb{i}+\pmb{x}}=\chi\tau^{1}-\eta\tau^{2}, & & u_{\pmb{i},\pmb{i}+\pmb{y}}=\chi\tau^{1}+\eta\tau^{2},\qquad a_{0}^{1,2,3}=0,\\[2pt]
u_{\pmb{i},\pmb{i}+2\pmb{x}+\pmb{y}}=u_{\pmb{i},\pmb{i}-\pmb{x}+2\pmb{y}}=+\lambda\tau^{3}, & & u_{\pmb{i},\pmb{i}+2\pmb{x}-\pmb{y}}=u_{\pmb{i},\pmb{i}+\pmb{x}+2\pmb{y}}=-\lambda\tau^{3},
\end{array}
\end{equation*}
在 $\chi\neq\eta$ 时，$\lambda=0$ 的情形描述 U1Cn01n 自旋液体（交错通量态）。
\begin{enumerate}
\item 证明：当 $\lambda\neq0$ 时，该拟设描述一个 $Z_{2}$ 自旋液体。
\item 找出保持该拟设不变的 U1Cn01n PSG 的子群。
\item 找出上述 $Z_{2}$ 拟设的标记。
\end{enumerate}

\section{对称自旋液体动物园}

\begin{keypoints}
\item 交错通量态邻域中 8 个 $Z_{2}$ 自旋液体的物理性质。
\end{keypoints}

本节中，我们想研究 9.4 节所分类的对称自旋液体的物理性质。我们一直用投影对称性来刻画不同的平均场态。然而，从一个平均场态的 PSG 直接看出它的物理性质并不容易。为了找出平均场态的物理性质，我们需要构造在相应 PSG 下不变的显式拟设。

然而，把几百个对称自旋液体一一列出绝非易事，更不用说逐一构造它们的拟设并研究其理论性质了。我们这里想做的是研究一些重要的自旋液体；可是，如何判定一个自旋液体的重要性呢？

在对高 $T_{c}$ 超导体的研究中，人们发现 $SU(2)$ 线性态 SU2Bn0（$\pi$ 通量态）、$U(1)$ 线性态 U1Cn01n（交错通量/$d$ 波态）与 $SU(2)$ 无能隙态 SU2An0（均匀 RVB 态）是重要的。它们与高 $T_{c}$ 超导体中观察到的某些相密切相关。$SU(2)$ 线性态、$U(1)$ 线性态与 $SU(2)$ 无能隙态分别再现了未掺杂、欠掺杂与过掺杂样品中观察到的电子谱函数。然而在理论上，由于 $U(1)$ 或 $SU(2)$ 规范涨落，这些自旋液体在低能下是不稳定的。这些态可能会变为它们邻域中更稳定的自旋液体。这促使我们研究 $SU(2)$ 线性态、$U(1)$ 线性态与 $SU(2)$ 无能隙态邻域中这些更稳定的自旋液体。为了进一步缩小范围，本节将主要研究 $U(1)$ 线性态 U1Cn01n 邻域中的自旋液体。\footnote{我们想指出，这里我们只研究对称自旋液体。U1Cn01n 自旋液体也可能变为其他一些破缺某些对称性的态。这样的对称性破缺相变实际上已经在高 $T_{c}$ 超导体中被观察到（如向反铁磁态、$d$ 波超导态与条纹态的转变）。}

\subsection{$U(1)$ 线性自旋液体周围的对称自旋液体}

U1Cn01n 自旋液体 (9.2.34) 可以连续变为 8 个不同的自旋液体，它们把 $U(1)$ 规范结构破缺到 $Z_{2}$ 规范结构。这 8 个自旋液体尽管都具有相同的对称性，却由不同的 PSG 标记。下面我们将更详细地研究这 8 个 $Z_{2}$ 自旋液体。特别地，我们想求出它们之中的自旋子谱。

第一个以 Z2A0013 标记，具有如下形式：
\begin{equation*}
\begin{array}{rcl}
u_{\pmb{i},\pmb{i}+\pmb{x}}=\chi\tau^{1}-\eta\tau^{2}, & & u_{\pmb{i},\pmb{i}+\pmb{y}}=\chi\tau^{1}+\eta\tau^{2},\\[2pt]
u_{\pmb{i},\pmb{i}+\pmb{x}+\pmb{y}}=+\gamma\tau^{1}, & & u_{\pmb{i},\pmb{i}-\pmb{x}+\pmb{y}}=+\gamma\tau^{1},\\[2pt]
a_{0}^{1}\neq0,\qquad a_{0}^{2,3}=0. & &
\end{array}
\tag{9.6.1}
\end{equation*}
它与拟设 (9.2.41) 具有相同的量子序。标记 Z2A0013 告诉我们该拟设的 PSG——即它的“对称性”群。提醒一下，对称自旋液体的 PSG 由对称性变换与规范变换的复合 $\{G_{0},\allowbreak\,G_{x}T_{x},\allowbreak\,G_{y}T_{y},\allowbreak\,G_{P_{x}}P_{x},\allowbreak\,G_{P_{y}}P_{y},\allowbreak\,G_{P_{xy}}P_{xy},\allowbreak\,G_{T}T\}$ 生成。规范变换 $G_{0}$ 生成该 PSG 的 IGG。标记中的“Z2”告诉我们 IGG 为 $Z_{2}$ 且 $G_{0}(\pmb{i})=(-)$；“A”告诉我们 $G_{x}(\pmb{i})=G_{y}(\pmb{i})=1$；接下来的“00”意味着 $G_{P_{x}}(\pmb{i})=1$、$G_{P_{y}}(\pmb{i})=1$；“1”意味着 $G_{P_{xy}}(\pmb{i})=\mathrm{i}\tau^{1}$，而“3”意味着 $G_{T}(\pmb{i})=\mathrm{i}\tau^{3}$。第二个拟设以 Z2Azz13 标记，具有如下形式：
\begin{equation*}
\begin{array}{rcl}
u_{\pmb{i},\pmb{i}+\pmb{x}}=\chi\tau^{1}-\eta\tau^{2}, & & u_{\pmb{i},\pmb{i}+\pmb{y}}=\chi\tau^{1}+\eta\tau^{2},\\[2pt]
u_{\pmb{i},\pmb{i}+\pmb{x}+\pmb{y}}=-\gamma\tau^{1}, & & u_{\pmb{i},\pmb{i}-\pmb{x}+\pmb{y}}=+\gamma\tau^{1},\\[2pt]
u_{\pmb{i},\pmb{i}+2\pmb{x}}=u_{\pmb{i},\pmb{i}+2\pmb{y}}=0,\qquad a_{0}^{1,2,3}=0. & &
\end{array}
\tag{9.6.2}
\end{equation*}
现在标记中的“zz”告诉我们 $\eta_{xpx}=\eta_{xpy}=-1$，$g_{P_{x}}=g_{P_{y}}=\mathrm{i}\tau^{3}$，于是 $G_{P_{x}}(\pmb{i})=(-)^{i_{x}+i_{y}}\tau^{3}$、$G_{P_{y}}(\pmb{i})=(-)^{i_{x}+i_{y}}\tau^{3}$（见式 (9.4.14)）。第三个拟设以 Z2A001n 标记，具有如下形式：
\begin{equation*}
\begin{array}{rcl}
a_{0}^{l}=0, & & \\[2pt]
u_{\pmb{i},\pmb{i}+\pmb{x}}=\chi\tau^{1}+\eta\tau^{2}, & & u_{\pmb{i},\pmb{i}+\pmb{y}}=\chi\tau^{1}-\eta\tau^{2}\\[2pt]
u_{\pmb{i},\pmb{i}+2\pmb{x}+\pmb{y}}=\chi_{1}\tau^{1}+\eta_{1}\tau^{2}+\lambda_{1}\tau^{3}, & & u_{\pmb{i},\pmb{i}-\pmb{x}+2\pmb{y}}=\chi_{1}\tau^{1}-\eta_{1}\tau^{2}-\lambda_{1}\tau^{3},\\[2pt]
u_{\pmb{i},\pmb{i}+2\pmb{x}-\pmb{y}}=\chi_{1}\tau^{1}+\eta_{1}\tau^{2}+\lambda_{1}\tau^{3}, & & u_{\pmb{i},\pmb{i}+\pmb{x}+2\pmb{y}}=\chi_{1}\tau^{1}-\eta_{1}\tau^{2}-\lambda_{1}\tau^{3}.
\end{array}
\tag{9.6.3}
\end{equation*}
如果你足够细心，就会发现标记 Z2A001n 并不出现在 9.4.3 节所分类的 196 个 $Z_{2}$ 自旋液体的名单之中（见表 9.1）。不过，以 Z2A001n 标记的 PSG 与以 Z2A003n 标记的 PSG 是规范等价的，而 Z2A003n 确实出现在名单中（见表 9.1 第二行）。下面我们将把上述自旋液体称为 Z2A003n 态。第四个拟设以 Z2Azz1n 标记，具有如下形式：
\begin{equation*}
\begin{array}{rcl}
a_{0}^{l}=0, & & \\[2pt]
u_{\pmb{i},\pmb{i}+\pmb{x}}=\chi\tau^{1}+\eta\tau^{2}, & & u_{\pmb{i},\pmb{i}+\pmb{y}}=\chi\tau^{1}-\eta\tau^{2}\\[2pt]
u_{\pmb{i},\pmb{i}+2\pmb{x}+\pmb{y}}=\chi_{1}\tau^{1}+\eta_{1}\tau^{2}+\lambda\tau^{3}, & & u_{\pmb{i},\pmb{i}-\pmb{x}+2\pmb{y}}=\chi_{1}\tau^{1}-\eta_{1}\tau^{2}+\lambda\tau^{3},\\[2pt]
u_{\pmb{i},\pmb{i}+2\pmb{x}-\pmb{y}}=\chi_{1}\tau^{1}+\eta_{1}\tau^{2}-\lambda\tau^{3}, & & u_{\pmb{i},\pmb{i}+\pmb{x}+2\pmb{y}}=\chi_{1}\tau^{1}-\eta_{1}\tau^{2}-\lambda\tau^{3}.
\end{array}
\tag{9.6.4}
\end{equation*}

上面四个拟设具有平移不变性。接下来四个 $Z_{2}$ 拟设不具有平移不变性，因为它们都是 Z2B 型的。（不过，投影之后它们仍然描述平移对称的自旋液体。）这些 $Z_{2}$ 自旋液体如下：

Z2B0013：
\begin{equation*}
\begin{array}{rcl}
u_{\pmb{i},\pmb{i}+\pmb{x}}=\chi\tau^{1}-\eta\tau^{2}, & & u_{\pmb{i},\pmb{i}+\pmb{y}}=(-)^{i_{x}}(\chi\tau^{1}+\eta\tau^{2}),\\[2pt]
u_{\pmb{i},\pmb{i}+2\pmb{x}}=-\gamma_{2}\tau^{1}+\lambda_{2}\tau^{2}, & & u_{\pmb{i},\pmb{i}+2\pmb{y}}=-\gamma_{2}\tau^{1}-\lambda_{2}\tau^{2},\\[2pt]
a_{0}^{1}\neq0,\qquad a_{0}^{2,3}=0, & &
\end{array}
\tag{9.6.5}
\end{equation*}

Z2Bzz13：
\begin{equation*}
\begin{array}{rcl}
u_{\pmb{i},\pmb{i}+\pmb{x}}=\chi\tau^{1}-\eta\tau^{2}, & & u_{\pmb{i},\pmb{i}+\pmb{y}}=(-)^{i_{x}}(\chi\tau^{1}+\eta\tau^{2}),\\[2pt]
u_{\pmb{i},\pmb{i}+2\pmb{x}+2\pmb{y}}=-\gamma_{1}\tau^{1}, & & u_{\pmb{i},\pmb{i}-2\pmb{x}+2\pmb{y}}=\gamma_{1}\tau^{1},\\[2pt]
a_{0}^{1,2,3}=0, & &
\end{array}
\tag{9.6.6}
\end{equation*}

Z2B001n：
\begin{equation*}
\begin{array}{rcl}
u_{\pmb{i},\pmb{i}+\pmb{x}}=\chi\tau^{1}+\eta\tau^{2}, & & u_{\pmb{i},\pmb{i}+\pmb{y}}=(-)^{i_{x}}(\chi\tau^{1}-\eta\tau^{2}),\\[2pt]
u_{\pmb{i},\pmb{i}+2\pmb{x}+\pmb{y}}=(-)^{i_{x}}\lambda\tau^{3}, & & u_{\pmb{i},\pmb{i}-\pmb{x}+2\pmb{y}}=-\lambda\tau^{3},\\[2pt]
u_{\pmb{i},\pmb{i}+2\pmb{x}-\pmb{y}}=(-)^{i_{x}}\lambda\tau^{3}, & & u_{\pmb{i},\pmb{i}+\pmb{x}+2\pmb{y}}=-\lambda\tau^{3},\\[2pt]
a_{0}^{l}=0, & &
\end{array}
\tag{9.6.7}
\end{equation*}

以及 Z2Bzz1n：
\begin{equation*}
\begin{array}{rcl}
u_{\pmb{i},\pmb{i}+\pmb{x}}=\chi\tau^{1}+\eta\tau^{2}, & & u_{\pmb{i},\pmb{i}+\pmb{y}}=(-)^{i_{x}}(\chi\tau^{1}-\eta\tau^{2}),\\[2pt]
u_{\pmb{i},\pmb{i}+2\pmb{x}+\pmb{y}}=(-)^{i_{x}}(\chi_{1}\tau^{1}+\eta_{1}\tau^{2}+\lambda\tau^{3}), & & u_{\pmb{i},\pmb{i}-\pmb{x}+2\pmb{y}}=\chi_{1}\tau^{1}-\eta_{1}\tau^{2}+\lambda\tau^{3},\\[2pt]
u_{\pmb{i},\pmb{i}+2\pmb{x}-\pmb{y}}=(-)^{i_{x}}(\chi_{1}\tau^{1}+\eta_{1}\tau^{2}-\lambda\tau^{3}), & & u_{\pmb{i},\pmb{i}+\pmb{x}+2\pmb{y}}=\chi_{1}\tau^{1}-\eta_{1}\tau^{2}-\lambda\tau^{3},\\[2pt]
a_{0}^{l}=0. & &
\end{array}
\tag{9.6.8}
\end{equation*}

对于式 (9.6.1)、(9.6.2)、(9.6.3) 与 (9.6.4) 给出的前四个 $Z_{2}$ 自旋液体（假定式 (9.6.1) 中 $a_{0}^{1}$ 较小），自旋子在四个孤立点处无能隙，并具有线性色散（见图 9.8）。因此，这四个拟设描述对称的 $Z_{2}$ 线性自旋液体。第二个 $Z_{2}$ 自旋液体 Z2Azz13 的单自旋子（single-spinon）色散相当有趣：它不具有 $90^{\circ}$ 旋转对称性。这与基态中的 $90^{\circ}$ 对称性并不矛盾，因为奇数个自旋子的激发永远无法满足约束，因而是被禁止的。自旋子色散具有绕 $\pmb{k}=(0,\pi)$ 的 $90^{\circ}$ 旋转对称性，而偶数个自旋子的激发谱则具有 $90^{\circ}$ 旋转对称性。

\begin{figure}[H]
\centering
\includegraphics[width=0.72\textwidth]{fig_9.8.pdf}
\caption*{图 9.8\quad $Z_{2}$ 线性自旋液体的自旋子色散 $E_{+}(\pmb{k})$ 作为 $(k_{x}/2\pi,\,k_{y}/2\pi)$ 的函数的等高线图：(a) 式 (9.6.1) 的 Z2A0013 态；(b) 式 (9.6.2) 的 Z2Azz13 态；(c) 式 (9.6.3) 的 Z2A001n 态；(d) 式 (9.6.4) 的 Z2Azz1n 态。}
\end{figure}

我们想指出，当 $a_{0}^{1}$ 较大时，式 (9.6.1) 可能有有能隙的自旋子谱。因此，Z2A0013 态既可能是有能隙自旋液体，也可能是无能隙自旋液体。其余三个 Z2A 态总是无能隙的，是 $Z_{2}$ 线性态。在 9.9 节中我们将证明，所有 $Z_{2}$ 有能隙态与 $Z_{2}$ 线性态都是稳定的。这些态的平均场结果可以应用于物理自旋液体。

接下来让我们考虑式 (9.6.5) 中的拟设 Z2B0013。拟设 (9.6.5) 的自旋子谱由下式的本征值确定：
\begin{equation*}
H=-2\big[\chi\cos(k_{x})\Gamma_{0}-\eta\cos(k_{x})\Gamma_{2}-\chi\cos(k_{y})\Gamma_{1}+\eta\cos(k_{y})\Gamma_{3}\big]+\lambda\Gamma_{4}
\end{equation*}

\begin{figure}[H]
\centering
\includegraphics[width=0.72\textwidth]{fig_9.9.pdf}
\caption*{图 9.9\quad $Z_{2}$ 线性态的自旋子色散 $\min(E_{1}(\pmb{k}),E_{2}(\pmb{k}))$ 作为 $(k_{x}/2\pi,\,k_{y}/2\pi)$ 的函数的等高线图：(a) 式 (9.6.5) 的 Z2B0013 态；(b) 式 (9.6.6) 的 Z2Bzz13 态；(c) 式 (9.6.7) 的 Z2B001n 态；(d) 式 (9.6.8) 的 Z2Bzz1n 态。}
\end{figure}

其中 $k_{x}\in(0,\pi)$，$k_{y}\in(-\pi,\pi)$，而
\begin{equation*}
\begin{array}{rcl}
\Gamma_{0}=\tau^{1}\otimes\tau^{3}, & \quad\Gamma_{1}=\tau^{1}\otimes\tau^{1}, & \quad\Gamma_{2}=\tau^{2}\otimes\tau^{3},\\[2pt]
\Gamma_{3}=\tau^{2}\otimes\tau^{1}, & \quad\Gamma_{4}=\tau^{1}\otimes\tau^{0}, &
\end{array}
\end{equation*}
这里已取 $\gamma_{1,2}=\lambda_{2}=0$。自旋子色散的四支能带具有 $\pm E_{1}(\pmb{k})$、$\pm E_{2}(\pmb{k})$ 的形式。我们发现自旋子谱在 $\pmb{k}=(\pi/2,\pm\pi/2)$ 附近的 8 个孤立点处为零（见图 9.9(a)）。因此，Z2B0013 态是一个 $Z_{2}$ 线性自旋液体。

其余三个 Z2B 态的自旋子谱可以用类似方法求出，谱图画在图 9.9 中。这三个 Z2B 态也是 $Z_{2}$ 线性态。无能隙自旋子只出现在 $\big(\frac{\pi}{2},\pm\frac{\pi}{2}\big)$ 处。

% ---------------- 新增术语（ch9_e） ----------------
% | staggered-flux phase | 交错通量相（ch9_c 已收 staggered-flux state 交错通量态） |
% | U(1)/SU(2) projective symmetry group | U(1)/SU(2) 投影对称群（= U(1)/SU(2) PSG） |
% | type-U1A/U1B/U1C PSG | U1A/U1B/U1C 型 PSG |
% | label (of a PSG) | 标记（PSG 的标记） |
% | quantum critical state | 量子临界态 |
% | critical line | 临界线 |
% | unstable fixed point | 不稳定不动点 |
% | destabilize | 使……失稳 |
% | single-spinon (excitation/spectrum) | 单自旋子（激发/谱） |
% | non-collinear flux | 非共线通量（collinear 共线的，术语表已收） |
% | contour plot | 等高线图 |
% | the zoo of ... | ……动物园（章节名修辞，照译） |

% 块边界判定：
%   起点——书页 419（p434）首段 "Knowing the translational symmetry ..." 为完整新句；
%     前块止于书页 418 末 "The gapless spinons appear only at (pi/2, +/-pi/2)."，
%     无前块溢出内容，故本块不设 %JOIN%，亦无需 TAIL 区。
%   实际覆盖——书页 419 起于 9.6.1 节后部（式 (9.6.9) 之前的段落），依次为 9.6.2 节、
%     9.7 节、9.8 节（含 9.8.1、9.8.2）、9.9 节引言；9.5 节与 9.6/9.6.1 节标题、
%     式 (9.6.1)--(9.6.8)、图 9.8、图 9.9 均在前块（书页 415--418）内。
%   终点——书页 430（p445）末句 "... in the light of quantum orders and their PSG
%     characterization." 为 9.9 节引言末句，书页 431（p446）顶部起新小节 9.9.1，
%     本块自然收尾，不设 %--E-TAIL--。
% 蓝本问题记录：
%   1) 问题 9.6.3(a) 原书把 eta/chi 与两个态的对应关系印反（与 9.6.2 节正文矛盾：
%      正文为 chi=0 -> SU2An0、eta=0 -> SU2Bn0，与拟设形式逐项比对一致），
%      且将 SU2Bn0 误标为 "SU(2)-gapless"，已按正文与本意订正并加译注。
%   2) 式 (9.8.9) 末行原书漏印下标 0（作 $a^{1,2}=0$），按本意作 $a_0^{1,2}=0$。
%   3) 式 (9.6.9) 原书即印作 $k_1+k_2+\pi x$（而变换式为 $|\pmb{k}\rangle\to
%      |\pmb{k}+\pi\pmb{y}\rangle$），照原图转写。
%   4) 图 9.10 图注 "the two-spinon spectrum docs have" 中 "docs" 为 "does" 之误。

\begin{figure}[H]
\centering
\includegraphics[width=0.72\textwidth]{fig_9.10.pdf}
\caption*{图 9.10\quad $Z_2$ 自旋液体自旋子色散 $E_+(\pmb{k})$ 作为 $(k_x/2\pi,\,k_y/2\pi)$ 的函数的等高线图：(a) 式 (9.2.39) 中的 $Z_2$ 无能隙态 Z2Ax2(12)n；(b) 式 (9.6.10) 中的 $Z_2$ 二次方态 Z2Bx2(12)n。尽管 (a) 中的单自旋子色散缺乏旋转对称性和宇称对称性，双自旋子谱却具有这些对称性。}
\end{figure}

虽然上述 Z2B 自旋液体具有平移对称性，但发现自旋子谱只在格点布里渊区的一半上有定义，这一点看起来很奇怪。不过，这与物理自旋液体中的平移对称性并不矛盾，因为单自旋子激发并不是物理的。只有双自旋子激发才对应于物理激发，其谱应当定义在满布里渊区上。现在的问题是：如何从只在半布里渊区上有定义的单自旋子谱，得到定义在满布里渊区上的双自旋子谱。设 $|\pmb{k},1\rangle$ 与 $|\pmb{k},2\rangle$ 是单自旋子的两个本征态，能量为正，分别为 $E_1(\pmb{k})$ 与 $E_2(\pmb{k})$（这里 $k_x\in(-\pi/2,\pi/2)$，$k_y\in(-\pi,\pi)$）。沿 $\pmb{x}$ 的平移（随后再作一个规范变换）把 $|\pmb{k},1\rangle$ 与 $|\pmb{k},2\rangle$ 变成能量相同的另外两个本征态：
\begin{equation*}
\begin{aligned}
|\pmb{k},1\rangle &\to |\pmb{k}+\pi\pmb{y},1\rangle \\
|\pmb{k},2\rangle &\to |\pmb{k}+\pi\pmb{y},2\rangle
\end{aligned}
\end{equation*}
于是我们看到，双自旋子态 $|\pmb{k}_1,\alpha_1\rangle|\pmb{k}_2,\alpha_2\rangle\pm|\pmb{k}_1+\pi\pmb{y},\alpha_1\rangle|\pmb{k}_2+\pi\pmb{y},\alpha_2\rangle$ 的动量与能量由下式给出：
\begin{equation*}
\begin{aligned}
E_{2\text{-spinon}} &= E_{\alpha_1}(\pmb{k}_1)+E_{\alpha_2}(\pmb{k}_2) \\
\pmb{k} &= \pmb{k}_1+\pmb{k}_2,\ \ \pmb{k}_1+\pmb{k}_2+\pi\pmb{x}
\end{aligned}
\tag{9.6.9}
\end{equation*}
式 (9.6.9) 使我们能够由单自旋子谱构造出双自旋子谱。

\subsection{$SU(2)$ 自旋液体附近的一个奇怪的对称自旋液体}

在式 (9.2.32) 的均匀 RVB 态（即 $SU(2)$ 无能隙态 SU2An0）附近以及式 (9.2.33) 的 $\pi$ 通量态（即 $SU(2)$ 线性态 SU2Bn0）附近，存在许多类型的对称拟设。这里我们只考虑其中的一个——$Z_2$ 自旋液体 Z2By1(12)n（注意 Z2By1(12)n 与 Z2Bx2(12)n 规范等价），它具有如下形式：
\begin{equation*}
u_{\pmb{i},\pmb{i}+\pmb{x}}=\mathrm{i}\chi\tau^{0}+\eta\tau^{1},\qquad
u_{\pmb{i},\pmb{i}+\pmb{y}}=(-)^{i_x}(\mathrm{i}\chi\tau^{0}+\eta\tau^{2}),\qquad
a_0^{1,2,3}=0.
\tag{9.6.10}
\end{equation*}
上述拟设在 $\chi=0$ 时退化为 $SU(2)$ 无能隙态 SU2An0 的拟设，在 $\eta=0$ 时退化为 $SU(2)$ 线性态 SU2Bn0 的拟设。经过傅里叶变换，我们发现自旋子谱由下式决定：
\begin{equation*}
H=-2\chi\sin(k_x)\Gamma_0+2\eta\cos(k_x)\Gamma_2-2\chi\sin(k_y)\Gamma_1+2\eta\cos(k_y)\Gamma_3
\end{equation*}
其中 $k_x\in(-\pi/2,\pi/2)$，$k_y\in(-\pi,\pi)$，而
\begin{equation*}
\Gamma_0=\tau^{0}\otimes\tau^{3},\qquad
\Gamma_2=\tau^{1}\otimes\tau^{3},\qquad
\Gamma_1=\tau^{0}\otimes\tau^{1},\qquad
\Gamma_3=\tau^{2}\otimes\tau^{1}.
\end{equation*}
自旋子谱可以精确计算，它的四支具有 $\pm E_1(\pmb{k})$ 与 $\pm E_2(\pmb{k})$ 的形式。自旋子能量在两个孤立点 $\pmb{k}=(0,0)$ 与 $(0,\pi)$ 处为零。在 $\pmb{k}=0$ 附近，低能谱由下式给出（见图 9.10(b)）
\begin{equation*}
E=\pm\eta^{-1}\sqrt{(\chi^{2}+\eta^{2})^{2}(k_x^{2}-k_y^{2})^{2}+4\chi^{4}k_x^{2}k_y^{2}}
\end{equation*}
有趣（而且奇怪）的是，当 $\pmb{k}\to 0$ 时能量并不线性地趋于零，而是像 $k^{2}$ 那样趋于零！为了强调这种 $E\propto k^{2}$ 的二次方色散，我们将把这样的态称为 $Z_2$ 二次方自旋液体（$Z_2$-quadratic spin liquid）。

\subsection*{问题 9.6.1}

拟设 (9.6.3) 的 PSG。
\begin{enumerate}
\item 试由标记 Z2A001n 求出 Z2A001n PSG 的规范变换 $G_x$、$G_y$、$G_{P_x}$ 等。
\item 证明拟设 (9.6.3) 在 Z2A001n PSG 下不变。
\item 试求把 Z2A001n PSG 变换到 Z2A003n PSG 的规范变换。
\end{enumerate}

\subsection*{问题 9.6.2}

验证：对于由式 (9.6.10) 描述的自旋液体，只要 $\chi$ 与 $\eta$ 都不为零，绕回路 $\pmb{i}\to\pmb{i}+\pmb{x}\to\pmb{i}+\pmb{x}+\pmb{y}\to\pmb{i}+\pmb{y}\to\pmb{i}$ 与 $\pmb{i}\to\pmb{i}+\pmb{y}\to\pmb{i}-\pmb{x}+\pmb{y}\to\pmb{i}-\pmb{x}\to\pmb{i}$ 的 $SU(2)$ 通量算符（见式 (9.2.20)）就互不对易。因此，该自旋液体确实具有 $Z_2$ 规范结构。

\subsection*{问题 9.6.3}

(a) 证明：当 $\eta=0$ 时，式 (9.6.10) 中的拟设退化为 $SU(2)$ 线性态 SU2Bn0 的拟设；当 $\chi=0$ 时，退化为 $SU(2)$ 无能隙态 SU2An0 的拟设。（译注：原书此句将 $\eta$、$\chi$ 与两个态的对应关系印反，且把 SU2Bn0 误标为 “$SU(2)$-gapless”，现按 9.6.2 节正文及物理内容订正。）

(b) 试求自旋子谱 $\pm E_1$ 与 $\pm E_2$。

\section{量子序的物理测量}

\begin{keypoints}
\item 用 PSG 描述的量子序可以通过激发谱来测量。
\end{keypoints}

在用 PSG 从数学上刻画了量子序之后，我们想问：如何在实验中测量量子序？有能隙态中的量子序与拓扑序相联系，我们可以用基态简并、边缘态、准粒子统计等手段来测量拓扑序（Wen, 1990, 1995; Senthil and Fisher, 2001）。具有无能隙激发的态中的量子序则需要用不同的方式来测量。在本节中，我们想证明：一般而言，量子序可以通过无能隙激发的动力学性质来测量。然而，并非所有的动力学性质都是普适的。因此，在用无能隙激发来刻画并测量量子序之前，我们需要先甄别出无能隙激发的普适性质。量子序的 PSG 刻画使我们能够得到这些普适性质——我们只需找出具有同一 PSG 的所有拟设的无能隙激发所共有的性质。

为了演示上述想法，我们想研究双自旋子激发的谱。我们注意到，自旋子只能成对地产生，因此单自旋子谱不是物理的。我们还注意到，双自旋子谱中包含自旋 1 激发，而自旋 1 激发可以用中子散射实验来测量。

在给定动量处，双自旋子谱分布在一段或几段能量范围内。设 $E_{2s}(\pmb{k})$ 为动量 $\pmb{k}$ 处双自旋子谱的低边缘。在平均场理论中，双自旋子谱可以由单自旋子色散按如下方式构造：
\begin{equation*}
E_{2\text{-spinon}}(\pmb{k})=E_{1\text{-spinon}}(\pmb{q})+E_{1\text{-spinon}}(\pmb{k}-\pmb{q})
\end{equation*}
在图 9.11--9.14 中，我们给出了几个简单自旋液体的平均场 $E_{2s}$。如果平均场态对规范涨落是稳定的，我们预计平均场 $E_{2s}$ 应当与真实的 $E_{2s}$ 在定性上一致。

在我们的例子中，有八个稳定的 $Z_2$ 线性自旋液体（见图 9.11 与 9.12）。以这八个 $Z_2$ 自旋液体为例，我们可以演示自旋 1 激发的普适性质如何区分不同的量子序。

\begin{figure}[H]
\centering
\includegraphics[width=0.72\textwidth]{fig_9.11.pdf}
\caption*{图 9.11\quad $Z_2$ 线性自旋液体的 $E_{2s}(\pmb{k})$ 作为 $(k_x/2\pi,\,k_y/2\pi)$ 的函数的等高线图：(a) 式 (9.6.1) 中的 Z2A0013 态；(b) 式 (9.6.2) 中的 Z2Azz13 态；(c) 式 (9.6.3) 中的 Z2A001n 态；(d) 式 (9.6.4) 中的 Z2Azz1n 态。}
\end{figure}

(a) 通过考察自旋 1 激发的谱，可以把 Z2A 自旋液体与 Z2B 自旋液体区分开。对于 Z2B 自旋液体，自旋 1 谱以布里渊区的四分之一为周期（即谱在 $k_x\to k_x+\pi$ 与 $k_y\to k_y+\pi$ 下不变）。与之相反，Z2A 自旋液体中的自旋 1 谱不具有这种周期性。

(b) 自旋 1 谱的周期性还可以把 Z2A0013 与 Z2Azz13 自旋液体同 Z2A001n 与 Z2Azz1n 自旋液体区分开。对于 Z2A001n 与 Z2Azz1n 自旋液体，自旋 1 谱以布里渊区的二分之一为周期（即谱在 $(k_x,k_y)\to(k_x+\pi,k_y+\pi)$ 下不变）。此外，Z2A001n 与 Z2Azz1n 自旋液体恰好在 $(\pi,\pi)$、$(\pi,0)$ 和 $(0,\pi)$ 处有无能隙自旋 1 激发。

(c) Z2A0013 与 Z2Azz13 自旋液体可以通过考察 $(\pi,0)$ 和 $(0,\pi)$ 附近的无能隙自旋 1 激发来区分。Z2Azz13 自旋液体中的无能隙自旋 1 激发位于区边界上。

\begin{figure}[H]
\centering
\includegraphics[width=0.72\textwidth]{fig_9.12.pdf}
\caption*{图 9.12\quad $Z_2$ 线性自旋液体的 $E_{2s}(\pmb{k})$ 作为 $(k_x/2\pi,\,k_y/2\pi)$ 的函数的等高线图：(a) 式 (9.6.5) 中的 Z2B0013 态；(b) 式 (9.6.6) 中的 Z2Bzz13 态；(c) 式 (9.6.7) 中的 Z2B001n 态；(d) 式 (9.6.8) 中的 Z2Bzz1n 态。谱以布里渊区的四分之一为周期。}
\end{figure}

我们看到，自旋液体中的量子序可以通过探测自旋 1 激发的中子散射实验来测量。

接下来，让我们讨论 $U(1)$ 线性态 U1Cn01n（交错通量态）。人们曾提出用 U1Cn01n 态描述欠掺杂高 $T_c$ 超导体中的赝能隙金属态（Wen and Lee, 1996; Rantner and Wen, 2001）。U1Cn01n 态自然地解释了欠掺杂金属态中的赝能隙。作为一种代数自旋液体，U1Cn01n 态还解释了赝能隙态中的类 Luttinger 电子谱函数（Rantner and Wen, 2001; Franz and Tesanovic, 2001）以及 $(\pi,\pi)$ 自旋涨落的增强（Kim and Lee, 1999; Rantner and Wen, 2002）。

\begin{figure}[H]
\centering
\includegraphics[width=0.72\textwidth]{fig_9.13.pdf}
\caption*{图 9.13\quad $E_{2s}(\pmb{k})$ 作为 $(k_x/2\pi,\,k_y/2\pi)$ 的函数的等高线图：(a) 式 (9.2.39) 中的 $Z_2$ 无能隙态 Z2Ax2(12)n；(b) 式 (9.6.10) 中的 $Z_2$ 二次方态 Z2Bx2(12)n。}
\end{figure}

\begin{figure}[H]
\centering
\includegraphics[width=0.72\textwidth]{fig_9.14.pdf}
\caption*{图 9.14\quad 两个 $U(1)$ 线性自旋液体的 $E_{2s}(\pmb{k})$ 作为 $(k_x/2\pi,\,k_y/2\pi)$ 的函数的等高线图：(a) 式 (9.2.34) 中的 U1Cn01n 态（交错通量相）；(b) 无能隙相中式 (9.8.9) 的 U1Cn00x 态。谱以布里渊区的二分之一为周期。}
\end{figure}

从图 9.14 我们看到，U1Cn01n 态中自旋 1 激发的无能隙点总是位于 $\pmb{k}=(\pi,\pi)$、$(0,0)$、$(\pi,0)$ 和 $(0,\pi)$。在所有这四个 $\pmb{k}$ 点处，自旋 1 连续谱边缘的等能线都具有两个彼此交叠的椭圆的形状。而且，等能线不垂直于区边界。所有这些都是 U1Cn01n 态的普适性质。在中子散射实验中测量这些性质，将使我们能够判断赝能隙金属态是否由 U1Cn01n（交错通量）态描述。

由于 $(2+1)$ 维 $U(1)$ 规范理论的瞬子效应，U1Cn01n 态是不稳定的。因此，U1Cn01n 态不得不变为其他一些态，例如 9.6 节中讨论的八个 $Z_2$ 自旋液体，或这里未讨论的其他态。从图 9.11(a) 我们看到，如果观察到 $(\pi,\pi)$ 处的节点劈裂为 $(\pi\pm\delta,\pi\pm\delta)$ 处的四个节点，以及 $(\pi,0)$ 和 $(0,\pi)$ 处的节点劈裂为 $(\pi\pm\delta,0)$ 和 $(0,\pi\pm\delta)$ 处的两个节点，就可以用中子散射探测从 U1Cn01n 态到 $Z_2$ 线性态 Z2A0013 的相变。从图 9.11(b) 我们看到，对于从 U1Cn01n 态到 $Z_2$ 线性态 Z2Azz13 的相变，$(\pi,\pi)$ 处的节点仍劈裂为 $(\pi\pm\delta,\pi\pm\delta)$ 处的四个节点，但 $(\pi,0)$ 和 $(0,\pi)$ 处的节点以另一种方式劈裂为 $(\pi,\pm\delta)$ 和 $(\pm\delta,\pi)$ 处的两个节点。我们还可以研究从 U1Cn01n 态到其余六个 $Z_2$ 自旋液体的相变。我们发现，自旋 1 激发谱以某些具有特征的方式改变。因此，通过测量自旋 1 激发谱及其演化，我们不仅能够探测不改变任何对称性的量子相变，还能够辨别正在发生的是哪一个相变。

\section{大 $N$ 极限下 $J_1$--$J_2$ 模型的相图}

\subsection{大 $N$ 极限}

\begin{keypoints}
\item 自旋-1/2 模型可以推广为 $SP(2N)$ 模型。
\item 对于 $SP(2N)$ 模型，由 $SU(2)$ 投影构造得到的平均场理论涨落微弱，是一个好的近似。
\end{keypoints}

到目前为止，我们一直集中讨论如何刻画、分类和测量不同的自旋液体及其量子序，还没有讨论如何寻找物理自旋哈密顿量，来实现我们构造的数百个不同自旋液体中的一些。本节就要处理这个问题。

在平均场层次上，设计一个实现具有给定量子序的自旋液体的自旋哈密顿量并不很难；对一个给定的自旋哈密顿量求出其平均场基态也不很难。真正的问题是我们是否应当相信平均场结果。我们已经看到，如果所得到的平均场基态是不稳定的（即如果平均场涨落导致低能下发散的相互作用），那么平均场结果不可信，平均场态不对应任何真实的物理自旋态。我们也曾论证过，如果平均场基态是稳定的（即平均场涨落导致低能下消失的相互作用），那么平均场结果可信，平均场态确实对应一个真实的物理自旋态。

这里我们想指出，上面关于稳定平均场态的论述过于乐观。一个“稳定平均场态”在低能下没有发散的涨落，所以它不一定不稳定；但另一方面，它也不一定稳定。这是因为短距离涨落如果足够强，同样会导致相变和失稳。因此，要使一个平均场结果可靠，平均场态必须稳定（或边缘），\emph{而且}短距离涨落必须微弱。由于我们的自旋模型中没有任何小参数，即便是稳定的平均场态，其短距离涨落也不微弱。正因如此，即使是稳定态的平均场结果能否应用于我们的自旋-1/2 模型，也并不清楚。

下面我们把自旋模型推广到大 $N$ 模型。我们将证明，大 $N$ 模型的平均场态具有微弱的短距离涨落，因此大 $N$ 模型的稳定与边缘平均场态对应于真实的物理态，稳定态或边缘态的平均场结果可以应用于大 $N$ 模型。

让我们从 $SU(2)$ 平均场理论的如下路径积分表示出发：
\begin{equation*}
\begin{aligned}
Z &= \int \mathcal{D}\psi\,\mathcal{D}[a_0^l(\pmb{i})]\,\mathcal{D}U_{ij}\ \mathrm{e}^{\mathrm{i}\int\mathrm{d}t\,(L-\sum_i a_0^l(\pmb{i},t)\psi_i^{\dagger}\tau^{l}\psi_i)} \\
L &= \sum_i \psi_i^{\dagger}\mathrm{i}\partial_t\psi_i-\sum_{\langle ij\rangle}\frac{1}{4}J_{ij}\left[\frac{1}{2}\mathrm{Tr}\,U_{ij}U_{ij}^{\dagger}-(\psi_i^{\dagger}U_{ij}\psi_j+h.c)\right].
\end{aligned}
\tag{9.8.1}
\end{equation*}
它由 $SU(2)$ 平均场哈密顿量 (9.2.9) 得到，其中 $U_{ij}$ 具有式 (9.2.13) 给出的形式。\footnote{我们把系数从 $3/8$ 改成了 $1/4$，这样，在式 (9.8.1) 中把 $U_{ij}$ 积掉就会得到自旋哈密顿量 (9.1.1)。}把费米子积掉之后，我们得到 $U_{ij}$ 与 $a_0^l$ 的有效拉格朗日量：
\begin{equation*}
Z=\int \mathcal{D}[a_0^l(\pmb{i})]\,\mathcal{D}U_{ij}\ \mathrm{e}^{\mathrm{i}\int\mathrm{d}t\,L_{0,eff}(U_{ij},a_0^l)}
\end{equation*}
问题在于，在上述路径积分中，$U_{ij}$ 与 $a_0^l$ 的涨落并不微弱。

为了减小涨落，我们引入 $N$ 份费米子副本 $\psi_i^a$，把式 (9.8.1) 推广为如下路径积分（Ran and Wen, 2003）：
\begin{equation*}
\begin{aligned}
Z &= \int \mathcal{D}\psi\,\mathcal{D}[a_0^l(\pmb{i})]\,\mathcal{D}U_{ij}\ \mathrm{e}^{\mathrm{i}\int\mathrm{d}t\,(L-\sum_i a_0^l(\pmb{i},t)\psi_i^{a\dagger}\tau^{l}\psi_i^{a})} \\
L &= \sum_i \psi_i^{a\dagger}\mathrm{i}\partial_t\psi_i^{a}-\sum_{\langle ij\rangle}\frac{1}{4}J_{ij}\left[\frac{N}{2}\mathrm{Tr}\,U_{ij}U_{ij}^{\dagger}-(\psi_i^{a\dagger}U_{ij}\psi_j^{a}+h.c)\right]
\end{aligned}
\tag{9.8.2}
\end{equation*}
其中 $a=1,2,\dots,N$。把费米子积掉之后，我们得到
\begin{equation*}
Z=\int \mathcal{D}[a_0^l(\pmb{i})]\,\mathcal{D}U_{ij}\ \mathrm{e}^{\mathrm{i}\int\mathrm{d}t\,N L_{0,eff}(U_{ij},a_0^l)}
\end{equation*}
我们看到，在大 $N$ 极限下，$U_{ij}$ 与 $a_0^l$ 的涨落是微弱的，平均场近似是一个好的近似。同样清楚的是，大 $N$ 平均场理论是一个 $SU(2)$ 规范理论，$U_{ij}$ 与 $a_0^l$ 的涨落对应于 $SU(2)$ 规范涨落。

在 9.2.1 节中，我们从物理自旋哈密顿量出发构造了平均场哈密顿量。这里我们面临的是相反的问题：已知大 $N$ 平均场哈密顿量
\begin{equation*}
H_{\text{mean}}=\sum_{\langle ij\rangle}\frac{1}{4}J_{ij}\left[\frac{N}{2}\mathrm{Tr}\,U_{ij}U_{ij}^{\dagger}-(\psi_i^{a\dagger}U_{ij}\psi_j^{a}+h.c)\right]+\sum_i a_0^l\psi_i^{a\dagger}\tau^{l}\psi_i^{a}
\tag{9.8.3}
\end{equation*}
我们想求出与之对应的物理的大 $N$ 自旋模型。

构造物理模型最重要的一步是找到物理希尔伯特空间。物理希尔伯特空间是费米子希尔伯特空间的一个子空间，由 $SU(2)$ 规范不变态、即满足约束
\begin{equation*}
\psi_i^{a\dagger}\tau^{l}\psi_i^{a}|phy\rangle,\qquad l=1,2,3
\end{equation*}
的态构成，且该约束在每个格点 $\pmb{i}$ 上都要成立。

得到每个格点上的物理希尔伯特空间之后，我们还需要找到在物理希尔伯特空间内部起作用的物理算符。这些物理算符是 $SU(2)$ 规范不变算符（即与 $\psi_i^{a\dagger}\tau^{l}\psi_i^{a}$ 对易的算符）。让我们把每个格点上 $\psi$ 的所有 $SU(2)$ 规范不变双线性型写下来：
\begin{equation*}
\begin{aligned}
S^{ab+}&\equiv\frac{1}{2}\psi_\alpha^{a\dagger}\tilde\psi_\alpha^{b},\qquad
S^{ab-}\equiv\frac{1}{2}\tilde\psi_\alpha^{a\dagger}\psi_\alpha^{b},\\
S^{ab3}&\equiv\frac{1}{2}\left(\psi_\alpha^{a\dagger}\psi_\alpha^{b}-\delta^{ab}\right)
=\frac{1}{2}\left(\delta^{ab}-\tilde\psi_\alpha^{b\dagger}\tilde\psi_\alpha^{a}\right)
\end{aligned}
\end{equation*}
其中 $\tilde\psi^{a}\equiv\mathrm{i}\sigma_2\psi^{a*}$，且格点指标已经隐去。这些 $S$ 算符是自旋-1/2 模型的自旋算符的推广。

当 $N=1$ 时，大 $N$ 模型就成为自旋-1/2 模型，$S$ 算符生成 $SU(2)$ 自旋旋转群。当 $N>1$ 时，$S$ 算符生成的是什么群？首先，让我们数一数不同 $S$ 算符的个数。对于 $S^{ab+}$ 或 $S^{ab-}$，标记关于 $a$ 与 $b$ 是对称的，因此每种类型各有 $N(N+1)/2$ 个不同的算符。对于 $S^{ab3}$，两个标记并不对称，因此就简单地是 $N^{2}$ 个。总共有 $N(N+1)+N^{2}=2N^{2}+N$ 个不同的 $S$ 算符。

可以考察 $S$ 算符之间的如下对易关系：
\begin{equation*}
\begin{gathered}
\left[S^{ab3},S^{cd3}\right]=\frac{1}{2}\left(\delta^{bc}S^{ad3}-\delta^{ad}S^{cb3}\right)\\[2pt]
\left[S^{ab3},S^{cd+}\right]=\frac{1}{2}\left(\delta^{bc}S^{ad+}+\delta^{bd}S^{ac+}\right)\\[2pt]
\left[S^{ab3},S^{cd-}\right]=-\frac{1}{2}\left(\delta^{ad}S^{bc-}+\delta^{ac}S^{bd-}\right)\\[2pt]
\left[S^{ab+},S^{cd-}\right]=\frac{1}{2}\left(\delta^{ac}S^{bd3}+\delta^{ad}S^{bc3}+\delta^{bc}S^{ad3}+\delta^{bd}S^{ac3}\right)\\[2pt]
\left[S^{ab-},S^{cd-}\right]=0,\qquad \left[S^{ab+},S^{cd+}\right]=0
\end{gathered}
\end{equation*}
这些正是 $SP(2N)$ 代数的关系式。所以，$S$ 算符就是生成 $SP(2N)$ 群的那 $2N^{2}+N$ 个生成元。当 $N=1$ 时，$SP(2)$ 同构于 $SU(2)$ 自旋旋转群。

在式 (9.8.2) 中把 $U_{ij}$ 与 $a_0^l$ 积掉，我们得到大 $N$ 模型的如下物理哈密顿量：
\begin{equation*}
H=\sum_{\langle ij\rangle}\frac{J_{ij}}{N}S_i^{ab}\cdot S_j^{ba}
\end{equation*}
其中
\begin{equation*}
\begin{aligned}
S_i^{ab1}&=S_i^{ba1}=\frac{1}{2}\left(S_i^{ab+}+S_i^{ab-}\right)\\
S_i^{ab2}&=S_i^{ba2}=\frac{1}{2\mathrm{i}}\left(S_i^{ab+}-S_i^{ab-}\right),
\end{aligned}
\end{equation*}
而
\begin{equation*}
S_i^{ab}=\left(S_i^{ab1},S_i^{ab2},S_i^{ab3}\right)
\end{equation*}
我们发现 $H$ 与 $SP(2N)$ 生成元 $\sum_i S_i^{ab}$ 对易，因此我们的大 $N$ 模型具有 $SP(2N)$ 对称性。这里我们想指出，$S_i^{ab}$ 的三个分量实际上并非处于同等地位：前两个关于 $a$ 与 $b$ 标记是对称的，第三个则不然。

我们知道，当 $N=1$ 时，物理希尔伯特空间在每个格点上有两个态，这两个态构成 $SP(2)=SU(2)$ 的一个不可约表示。对更高的 $N$，每格点物理希尔伯特空间的维数为
\begin{equation*}
\begin{aligned}
N &=\ 2\ \ \ 3\ \ \ 4\ \ \ 5\ \ \ 6\ \ \dots\\
\text{Dimension} &=\ 5\ \ \ 14\ \ \ 43\ \ \ 142\ \ \ 429\ \ \dots
\end{aligned}
\tag{9.8.4}
\end{equation*}
结果发现，这些物理希尔伯特空间是 $SP(2N)$ 对称群的不可约表示。我们可以用不可约表示在某个特定 Cartan 基下的最高权态来标记它。$SP(2N)$ 的 Cartan 基可以取为每个 $\alpha$ 的 $z$ 分量自旋，即
\begin{equation*}
S^{\alpha\alpha 3},\qquad \alpha=1,2,\dots,N.
\tag{9.8.5}
\end{equation*}
于是，物理希尔伯特空间中的最高权态就是不含 $\psi$ 费米子的态。

\subsection{$SP(2N)$ 模型的相图}

\begin{keypoints}
\item 平均场相图就是大 $N$ 极限下 $SP(2N)$ 模型的相图。
\item 平均场理论或 $SP(2N)$ 模型中不同的相可以用 PSG 来标记。
\end{keypoints}

从式 (9.8.3) 我们看到，$SP(2N)$ 模型的 $SU(2)$ 平均场理论与自旋-1/2 模型的 $SU(2)$ 平均场理论具有相同的结构。对两个模型而言，平均场态都由拟设 $(U_{ij},a_0^l(\pmb{i}))$ 描述。特别地，$SP(2N)$ 模型的平均场能量恰好是 $SP(2)$ 模型（即自旋-1/2 模型）平均场能量的 $N$ 倍。因此，平均场相图不依赖于 $N$。正如前面讨论的自旋-1/2 模型那样，$SP(2N)$ 模型的不同平均场态用 PSG 来刻画和分类。

这里我们考虑正方格子上一个具体的 $SP(2N)$ 模型。该模型具有最近邻耦合 $J_1$ 和次近邻耦合 $J_2$。我们取 $J_1+J_2=1$。

\begin{figure}[H]
\centering
\includegraphics[width=0.72\textwidth]{fig_9.15.pdf}
\caption*{图 9.15\quad $J_1$--$J_2$ 自旋系统中各个相的平均场能量：“A” 标记 $\pi$ 通量态（$SU(2)$ 线性态 SU2Bn0）；“B” 标记式 (9.8.6) 中的 $SU(2)\times SU(2)$ 无能隙态；“C” 标记式 (9.8.7) 中的 $SU(2)\times SU(2)$ 线性态；“D” 标记手征自旋态（一个 $SU(2)$ 有能隙态）；“E” 标记破坏 $90^{\circ}$ 旋转对称性的 $U(1)$ 线性态 (9.8.8)；“F” 标记式 (9.8.9) 中的 $U(1)$ 有能隙态 U1Cn00x；“G” 标记式 (9.6.2) 中的 $Z_2$ 线性态 Z2Azz13；“H” 标记式 (9.6.1) 中的 $Z_2$ 线性态 Z2A0013；“I” 标记均匀 RVB 态（$SU(2)$ 无能隙态 SU2An0）。}
\end{figure}

在大 $N$ 极限下，$SU(2)$ 平均场理论 (9.8.3) 是一个好的近似，因此我们将用平均场理论来计算 $SP(2N)$ 模型的相图。描述平均场基态的平均场拟设通过把平均场能量极小化而求得，结果示于图 9.15。图 9.15 中的每条曲线代表一个局域极小的平均场能量，能量最低的曲线对应于平均场基态。我们发现，每条曲线上不同的拟设都具有同一个 PSG。这在预料之中，因为当我们改变 $J_2$ 时，曲线所描述的平均场能量以解析的方式变化，所以沿同一条曲线从一个点移动到另一个点并不发生量子相变。同一条曲线上的拟设属于同一个相，由同一个 PSG 描述。

然而，如果两条曲线相交，交点就代表一个量子相变，因为基态能量在交点处不是解析的。由于不同的曲线具有不同的 PSG，我们看到量子相变以 PSG 的改变为特征。下面我们将讨论图 9.15 中各平均场相的不同量子序（或 PSG）。

在图 9.15 中，相 A 是式 (9.2.33) 给出的 $\pi$ 通量态（SU2Bn0 $SU(2)$ 线性态）。相 B 是在对角链上带有两个相互独立的均匀 RVB 态的态。它在低能下具有 $SU(2)\times SU(2)$ 规范涨落，将被称为 $SU(2)\times SU(2)$ 无能隙态。其拟设由下式给出：
\begin{equation*}
u_{\pmb{i},\pmb{i}+\pmb{x}+\pmb{y}}=\chi\tau^{3},\qquad
u_{\pmb{i},\pmb{i}+\pmb{x}-\pmb{y}}=\chi\tau^{3},\qquad
a_0^{l}=0.
\tag{9.8.6}
\end{equation*}
相 C 是在对角链上带有两个相互独立的 $\pi$ 通量态的态。它在低能下具有 $SU(2)\times SU(2)$ 规范涨落，将被称为 $SU(2)\times SU(2)$ 线性态。其拟设由下式给出：
\begin{equation*}
u_{\pmb{i},\pmb{i}+\pmb{x}+\pmb{y}}=\chi(\tau^{3}+\tau^{1}),\qquad
u_{\pmb{i},\pmb{i}+\pmb{x}-\pmb{y}}=\chi(\tau^{3}-\tau^{1}),\qquad
a_0^{l}=0
\tag{9.8.7}
\end{equation*}
相 D 是手征自旋态 (9.1.22)。相 E 由拟设
\begin{equation*}
\begin{aligned}
u_{\pmb{i},\pmb{i}+\pmb{x}+\pmb{y}}&=\chi_1\tau^{1}+\chi_2\tau^{2}, &\qquad
u_{\pmb{i},\pmb{i}+\pmb{x}-\pmb{y}}&=\chi_1\tau^{1}-\chi_2\tau^{2}\\
u_{\pmb{i},\pmb{i}+\pmb{y}}&=\eta\tau^{3}, &\qquad
a_0^{l}&=0
\end{aligned}
\tag{9.8.8}
\end{equation*}
描述，它破坏 $90^{\circ}$ 旋转对称性，是一个 $U(1)$ 线性态。相 F 由如下 U1Cn00x 拟设描述：
\begin{equation*}
\begin{aligned}
u_{\pmb{i},\pmb{i}+\pmb{x}}&=\eta\tau^{1}, &\qquad
u_{\pmb{i},\pmb{i}+\pmb{y}}&=\eta\tau^{1}\\
u_{\pmb{i},\pmb{i}+\pmb{x}+\pmb{y}}&=\chi\tau^{3}, &\qquad
u_{\pmb{i},\pmb{i}+\pmb{x}-\pmb{y}}&=\chi\tau^{3}\\
a_0^{3}&=\lambda, &\qquad
a_0^{1,2}&=0.
\end{aligned}
\tag{9.8.9}
\end{equation*}
U1Cn00x 态可以是 $U(1)$ 线性态（若 $a_0^{3}$ 小），也可以是 $U(1)$ 有能隙态（若 $a_0^{3}$ 大）。相 F 的态结果是一个 $U(1)$ 有能隙态。相 G 由式 (9.6.2) 中的 Z2Azz13 拟设描述，是一个 $Z_2$ 线性态。相 H 由式 (9.6.1) 中的 Z2A0013 拟设描述，也是一个 $Z_2$ 线性态。相 I 是均匀 RVB 态（式 (9.2.32) 中的 $SU(2)$ 无能隙态 SU2An0）。

从图 9.15 我们观察到下列各对相之间（在平均场层次上）的连续相变：(A,D)、(A,G)、(B,G)、(C,E) 与 (B,H)。其中三个连续相变 (B,G)、(B,H) 与 (A,G) 不改变任何对称性。我们还注意到，在从相 A 到相 G 的连续相变中，相 A 的 $SU(2)$ 规范结构破缺到 $Z_2$；在 (B,G) 与 (B,H) 这两个相变中，相 B 的 $SU(2)\times SU(2)$ 规范结构破缺到 $Z_2$。

\section{量子序与平均场自旋液体的稳定性}

\begin{keypoints}
\item 许多无能隙平均场自旋液体对量子涨落可以是稳定的。即使存在长程规范相互作用，它们仍可以是稳定的。
\end{keypoints}

我们已经强调了平均场态对平均场涨落的稳定性的重要性。只有稳定的平均场态和边缘的平均场态才有可能描述物理自旋液体。我们在 9.1.5 节和 9.2.2 节中讨论过获得稳定平均场态的途径。这里我们将从量子序及其 PSG 刻画的角度再次讨论平均场态的稳定性。

% 块边界判定：
%   起点——书页 431（p446）顶部为 9.9.1 节标题，起新小节，无前块溢出内容，
%     故本块不设 %JOIN%，亦无需 X-TAIL 区。
%   实际覆盖——9.9.1--9.9.5 节、9.10 节（含 9.10.1、9.10.2）及问题 9.10.1、9.10.2；
%     9.7、9.8 两节与 9.9 节引言在前块 ch9_f（书页 421--430）内。
%   终点——书页 439（p454）以问题 9.10.2 (c) 收束，为第 9 章末页（p455 为第 10 章
%     章首），自然收尾，无需 TAIL。
% 蓝本问题记录：
%   1) 式 (9.10.6) 右端蓝本与原书均印作 $U_{P_x}$，按上下文应为 $U_{P_y}$，已订正并加译注；
%   2) 式 (9.10.6) 引导句蓝本与原书均印作 $G_{P_y}T_{P_y}$，应为 $G_{P_y}P_y$，已按本意
%      译出并加译注；
%   3) 9.10.1 节蓝本 OCR 作 "IGG=\{c^{1(-)^{i}\theta\tau^{3}}\}"，按图片订正为
%      $\{\mathrm{e}^{\mathrm{i}(-)^{\pmb{i}}\theta\tau^{3}}\}$；
%   4) 9.10.2 节 "there will a gapless U(1) gauge boson" 原书漏 "be"，按本意译出；
%   5) 蓝本 9.9.4 节把 Z2Bx2(12)n 误排作 "Z2Bx2(12)n"（与图片一致，无碍）；
%      蓝本 9.9.3 节段首 "wefdw..." 等为文字层 OCR 噪声，未采信。

\subsection{投影对称群——量子相的一个普适性质}

\begin{keypoints}
\item 如果相应的平均场态是稳定的（即没有红外发散），PSG 就是真实量子相的一个普适性质。此时，PSG 描述该量子相中的量子序。
\end{keypoints}

在本节中，我们想证明 PSG 可以成为量子态的一种普适性质，其含义是它对微扰涨落是稳健的。因此，PSG 作为一种普适性质，可以用来刻画一个量子相。任何与 PSG 相联系（或受 PSG 保护）的物理性质也是相的普适性质，可以用于在实验中探测和测量量子序。特别地，PSG 能够保护无能隙规范激发和费米子激发（见 9.10 节）。因此，PSG 的稳定性也蕴含着无能隙规范激发和费米子激发的稳定性。

我们知道，平均场自旋液体态由 $U_{ij}=\langle\psi_{\pmb{i}}\psi_{\pmb{j}}^{\dagger}\rangle$ 刻画。如果在平均场态附近纳入微扰涨落，我们预期 $U_{ij}$ 会获得微扰修正 $\delta U_{ij}$。这里我们想论证：微扰涨落只能以这样的方式改变 $U_{ij}$，使得 $U_{ij}$ 与 $U_{ij}+\delta U_{ij}$ 具有相同的 PSG。

首先，我们想指出以下众所周知的事实。微扰涨落不能改变对称性和规范结构。例如，如果拟设 $U_{ij}$ 与哈密顿量具有某个对称性，那么由微扰涨落生成的 $\delta U_{ij}$ 将具有同样的对称性。类似地，微扰涨落不能生成这样的 $\delta U_{ij}$——比方说，把一个 $U(1)$ 规范结构破缺到 $Z_2$ 规范结构的 $\delta U_{ij}$。

由于规范结构（由 IGG 描述）和对称性群都是 PSG 的一部分，把上述观察推广为“除 PSG 中的 IGG 与对称性群之外，整个 PSG 都不能被微扰涨落改变”是合理的。

事实上，平均场哈密顿量和平均场基态在 PSG 中的变换下是不变的。因此，在围绕一个平均场态的微扰计算中，PSG 中的变换表现得就像普通的对称性变换一样。所以，微扰涨落只能生成在 PSG 的变换下不变的 $\delta U_{ij}$。

由于微扰涨落（按定义）不改变相，$U_{ij}$ 与 $U_{ij}+\delta U_{ij}$ 描述同一个相。换言之，我们可以把 $U_{ij}$ 分成若干类（称为普适类，universality classes），使得每一类中的 $U_{ij}$ 彼此通过微扰涨落相连。每一个普适类描述一个相。我们看到，如果上述论证成立，那么同一普适类中的拟设都享有同一个 PSG。换言之，普适类或者相是由 PSG（或量子序）分类的。

我们想指出，在上面的讨论中我们假定了微扰涨落没有红外发散。红外发散意味着微扰涨落是相关微扰。这种发散的修正可能引起相变，从而使上述论证失效。因此，上述论证和结果只适用于满足下列要求的稳定自旋液体：(i) 所有平均场涨落都没有红外发散；(ii) 平均场涨落在格点尺度上足够弱。

在下文中，我们将讨论几类平均场态的稳定性。要求 (ii) 可以通过大 $N$ 极限和/或（必要时）调整自旋哈密顿量中的短程自旋耦合来满足。这里我们将主要考虑要求 (i)。我们发现，至少在某些大 $N$ 极限下，许多（但并非全部）平均场态确实对应于在低能下稳定的真实量子自旋液体。

迄今为止所研究的所有自旋液体（每个元胞具有奇数个电子）可以分为四类。下面我们将逐一研究每一类。

\subsection{刚性自旋液体}

刚性自旋液体（rigid spin liquid）是自旋子和所有其他激发都完全有能隙的态。规范场的能隙可以由 Chern--Simons 项或 Anderson--Higgs 机制产生。有能隙的规范场只在自旋子之间诱导短程相互作用。由于低能下没有激发，刚性自旋液体是稳定的态。刚性自旋液体由拓扑序刻画，并且具有未禁闭的中性自旋-1/2 激发。刚性自旋液体的低能有效理论是拓扑场论。$Z_2$ 有能隙自旋液体和手征自旋液体是刚性自旋液体的例子。

\subsection{玻色自旋液体}

U1Cn00x 态 (9.8.9) 在式 (9.8.9) 中的 $a_{0}^{3}$ 足够大时，可以是一个 $U(1)$ 有能隙自旋液体。这样的态在平均场层次上具有有能隙的自旋子和无能隙的 $U(1)$ 规范玻色子。我们将称之为玻色自旋液体（Bose spin liquid）。无能隙 $U(1)$ 规范涨落的动力学由如下低能有效理论描述：
\begin{equation*}
\mathcal{L}=\frac{1}{8\pi g^{2}}\left(f_{\mu\nu}\right)^{2}
\end{equation*}
其中 $f_{\mu\nu}$ 是 $U(1)$ 规范场的场强。然而，在 $1+2$ 维中，并且在纳入瞬子效应之后，$U(1)$ 规范涨落将具有红外发散，这导致规范玻色子获得能隙以及自旋子被禁闭（Polyakov, 1977）。因此，平均场 $U(1)$ 有能隙态在 $1+2$ 维中是不稳定的。它们的 PSG 可能并不描述物理自旋液体中任何真实的量子序。\footnote{我们想\emph{提一句}，平均场 $U(1)$ 有能隙态在 $1+3$ 维中是\emph{稳定的}。}

\subsection{费米自旋液体}

费米自旋液体（Fermi spin liquid）由以下两条性质定义。其一，它们具有由自旋-1/2 费米子描述的无能隙激发；其二，这些无能隙激发之间只有短程相互作用。$Z_2$ 线性自旋液体、$Z_2$ 二次方自旋液体和 $Z_2$ 无能隙自旋液体是费米自旋液体的例子。

在 $Z_2$ 线性自旋液体中，自旋子具有无质量的狄拉克色散。这些具有短程相互作用的自旋子在连续极限下由如下有效理论描述（见问题 9.1.4）：
\begin{equation*}
S=\int\mathrm{d}t\,\mathrm{d}^{2}x\,\left[\bar{\psi}\partial_{\mu}\gamma^{\mu}\psi+\left(\bar{\psi}M\psi\right)^{2}\right]
\end{equation*}
在 $2+1$ 维中，$\psi$ 的标度量纲为 $[\psi]=1$，因此作用量是无量纲的。于是，相互作用项 $(\psi)^{4}$ 的标度量纲为 $[\psi^{4}]=4$，比 3 要大。我们看到，无质量狄拉克费米子之间的短程相互作用在 $1+2$ 维中是不相关的。因此，$Z_2$ 线性自旋液体是稳定的态。

现在让我们考虑式 (9.6.10) 中 $Z_2$ 二次方自旋液体 Z2Bx2(12)n 的稳定性。在 $Z_2$ 二次方自旋液体中，自旋子具有无能隙的二次方色散 $\omega\propto k^{2}$。在这种情形下，空间和时间具有不同的标度量纲，分别为 $[x^{-1}]\equiv1$ 与 $[t^{-1}]=2$。从无量纲的连续作用量
\begin{equation*}
S=\int\mathrm{d}t\,\mathrm{d}^{2}x\,\left[\psi^{\dagger}\mathrm{i}\partial_{t}\psi-c\psi^{\dagger}(\partial_{x})^{2}\psi\right]
\end{equation*}
我们看到，自旋子场 $\psi$ 的标度量纲为 $[\psi]=1$，而四费米子相互作用项的量纲为 $[\psi^{4}]=4$。于是，相互作用作用量 $S_{\mathrm{int}}=\int\mathrm{d}t\,\mathrm{d}^{2}x\,g(\psi^{\dagger}\psi)^{2}$ 中的耦合常数 $g$ 的标度量纲为 $[g]=0$，因为 $S_{\mathrm{int}}$ 的标度量纲总是零。因此，与 $Z_2$ 线性自旋液体不同，$Z_2$ 二次方态中无能隙自旋子之间的短程相互作用在 $1+2$ 维中是边缘的。相互作用的高阶效应会使耦合变得相关还是不相关，仍需进一步的研究才能确定。这将决定 $Z_2$ 二次方自旋液体在平均场层次之外的动力学稳定性。

在 $1+2$ 维中，$Z_2$ 无能隙自旋液体与费米液体一样稳定。如果我们假定费米液体在 $1+2$ 维中是稳定的，那么 $Z_2$ 无能隙自旋液体也是稳定的。

\subsection{代数自旋液体}

代数自旋液体（algebraic spin liquid）是具有无能隙激发的态，但其中没有任何一种无能隙激发是由自由的玻色型或费米型准粒子描述的。$U(1)$ 线性自旋液体是代数自旋液体的例子。它们的低能激发由与 $U(1)$ 规范场耦合的无质量狄拉克费米子描述。由于费米子--规范场耦合在微扰论的所有阶都是严格边缘的（Appelquist and Nash, 1990），无能隙激发一直低到零能都具有有限的相互作用。结果，低能下既没有自由的玻色型准粒子（如规范玻色子），也没有自由的费米型准粒子（如自旋子）。这使得关于这些态的稳定性的讨论困难得多。关于代数自旋液体的讨论，感兴趣的读者可参阅 Rantner and Wen (2002) 与 Wen (2002c)。

已经证明，在 $2+1$ 维中，$U(1)$ 线性态在自旋模型的某些平均场涨落微弱的大 $N$ 极限下是稳定的。$U(1)$ 线性态的 PSG 阻止微扰涨落生成破坏稳定性的抵消项（counter term），例如费米子质量项。这保证了 $U(1)$ 线性态的稳定性。因此，至少在大 $N$ 极限下，相应的代数自旋液体作为物理自旋系统的相而存在。

代数自旋液体的存在是一个非常惊人的现象。按照通常的看法，如果玻色子/费米子在低能下相互作用，那么相互作用将为这些低能激发打开一个能隙。这蕴含着：一个系统要么在低能下具有自由的玻色型/费米型激发，要么根本没有低能激发。代数自旋液体的存在表明，这种通常的看法是不正确的。它提出了一个重要问题：是什么保护着无能隙激发（特别是当它们在所有能标上都相互作用时）？无能隙激发的存在应当有一个“理由”或“原理”。在下一节中我们将证明，PSG 阻止规范玻色子和费米子获得质量项。\emph{所以，正是量子序及其相应的 PSG 保护着无能隙激发。}

\section{量子序与无能隙规范玻色子和费米子}

\begin{keypoints}
\item 量子序能够产生并保护无能隙规范玻色子与无能隙费米子，正如对称性破缺能够产生并保护无能隙 Nambu--Goldstone 玻色子一样。
\end{keypoints}

无能隙激发在自然界和凝聚态系统中都非常罕见。因此，如果我们看到一个无能隙激发，我们就会想问它为什么存在。无能隙激发的一种来源是自发对称性破缺，它给出 Nambu--Goldstone 玻色子（Nambu, 1960; Goldstone, 1961）。无能隙激发与自发对称性破缺之间的这种关系非常重要。正是由于这种关系，我们可以在不了解系统细节的情况下，仅从对称性出发得到一个复杂态的低能物理。这一思路使得 Landau 关于相与相变的对称性破缺理论（Ginzburg and Landau, 1950; Landau and Lifschitz, 1958）成为一个研究相的低能性质的极为强有力的理论。然而，自发对称性破缺并不是无能隙激发的唯一来源。量子序及其相应的 PSG 也能产生并保护无能隙激发。令人惊奇的是，量子序产生并保护的是无能隙规范玻色子和无能隙费米子。无能隙费米子甚至可以从纯玻色模型中产生。在本节中，我们想较为详细地讨论量子序（及其 PSG）与无能隙规范/费米子激发之间的关系（Wen, 2002a,c; Wen and Zee, 2002）。

\subsection{投影对称群与无能隙规范玻色子}

\begin{keypoints}
\item （在平均场层次上）无能隙规范玻色子的规范群就是 PSG 的 IGG。
\end{keypoints}

无能隙规范涨落与量子序之间的关系简单而直接。量子有序态中无能隙规范涨落的规范群，就是刻画该量子序的 PSG 中的 IGG。

为了看清量子序及其 PSG 如何产生并保护无能隙规范玻色子，作为一个例子，让我们假定一个量子有序态的 IGG $\mathcal{G}$ 包含一个 $U(1)$ 子群，它由下列常值规范变换构成：
\begin{equation*}
\{W_{\pmb{i}}=\mathrm{e}^{\mathrm{i}\theta\tau^{3}}\,|\,\theta\in[0,2\pi)\}\subset\mathcal{G}
\end{equation*}
接着我们考虑围绕平均场解 $\bar{u}_{ij}$ 的一类涨落，即 $u_{ij}=\bar{u}_{ij}\mathrm{e}^{\mathrm{i}a_{ij}^{3}\tau^{3}}$。由于 $\bar{u}_{ij}$ 在常值规范变换 $\mathrm{e}^{\mathrm{i}\theta\tau^{3}}$ 下不变，一个依赖空间的规范变换 $\mathrm{e}^{\mathrm{i}\theta_{\pmb{i}}\tau^{3}}$ 将把涨落 $a_{ij}^{3}$ 变为 $\tilde{a}_{ij}^{3}=a_{ij}^{3}+\theta_{\pmb{i}}-\theta_{\pmb{j}}$。这意味着 $a_{ij}^{3}$ 与 $\tilde{a}_{ij}^{3}$ 标记同一个物理态，而 $a_{ij}^{3}$ 对应于规范涨落。涨落的能量具有规范不变性，即 $E(\{a_{ij}^{3}\})=E(\{\tilde{a}_{ij}^{3}\})$。由于能量的这种规范不变性，我们看到规范场的质量项 $(a_{ij}^{3})^{2}$ 是不被允许的。因此，由 $a_{ij}^{3}$ 描述的 $U(1)$ 规范涨落将在低能下出现。

如果 $\mathcal{G}$ 的 $U(1)$ 子群由如下依赖空间的规范变换构成：
\begin{equation*}
\{W_{\pmb{i}}=\mathrm{e}^{\mathrm{i}\theta\,\pmb{n}_{\pmb{i}}\cdot\pmb{\tau}}\,|\,\theta\in[0,2\pi),\ |\pmb{n}_{\pmb{i}}|=1\}\subset\mathcal{G},
\end{equation*}
那么我们总可以用一个 $SU(2)$ 规范变换把每个格点上的 $\pmb{n}_{\pmb{i}}$ 转到 $z$ 方向，从而把问题约化为上面讨论过的情形。因此，无论 IGG 中的规范变换是否依赖空间，低能规范涨落的规范群总是由 $\mathcal{G}$ 给出。由于 IGG 的每一个 $U(1)$ 子群都对应于一个无能隙 $U(1)$ 规范玻色子，低能规范玻色子的规范群由 IGG 给出，即使 IGG 是非阿贝尔的。

我们想说明的是，有时低能规范涨落不仅出现在 $\pmb{k}=0$ 附近，还出现在其他一些 $\pmb{k}$ 点附近。在这种情形下，我们将有几个低能规范场，每个 $\pmb{k}$ 点对应一个。9.8 节中讨论的 $SU(2)$ 从属玻色子理论的某些拟设给出了这种现象的例子，它们在低能下具有 $SU(2)\times SU(2)$ 规范结构。我们看到，低能规范结构 $SU(2)\times SU(2)$ 甚至可以比高能规范结构 $SU(2)$ 还要大。即使对于低能规范涨落出现在不同 $\pmb{k}$ 点附近这种复杂情形，IGG 仍然正确地描述相应拟设的低能规范结构。如果 IGG 包含不依赖空间坐标的规范变换，那么这类变换对应于 $\pmb{k}=0$ 附近无能隙规范涨落的规范群。如果 IGG 包含依赖空间坐标的规范变换，那么这些变换对应于非零 $\pmb{k}$ 附近无能隙规范涨落的规范群。例如，若 $\mathrm{IGG}=\{\mathrm{e}^{\mathrm{i}(-)^{\pmb{i}}\theta\tau^{3}}\}$，则在 $\pmb{k}=(\pi,\pi)$ 附近就会出现一个无能隙 $U(1)$ 规范玻色子，由规范场 $(-)^{\pmb{i}}a^{3}(\pmb{x})$ 描述。这样，IGG 为我们提供了对所有低能规范涨落的统一处理，无论它们的晶体动量如何。

在本章中，我们在许多地方使用了 $Z_2$ 自旋液体、$U(1)$ 自旋液体、$SU(2)$ 自旋液体和 $SU(2)\times SU(2)$ 自旋液体这些术语。现在我们可以对这些低能 $Z_2$、$U(1)$、$SU(2)$ 和 $SU(2)\times SU(2)$ 规范群给出精确定义。这些低能规范群就是相应拟设的 IGG。它们与出现在 $SU(2)$、$U(1)$ 或 $Z_2$ 从属玻色子方法中的高能规范群毫无关系。我们还使用过一个平均场态的 $Z_2$ 规范结构、$U(1)$ 规范结构和 $SU(2)$ 规范结构这些术语。它们精确的数学含义同样是相应拟设的 IGG。当我们说一个 $U(1)$ 规范结构被破缺到一个 $Z_2$ 规范结构时，我们的意思是：拟设被改变了，使得它的 IGG 从 $U(1)$ 群变成了 $Z_2$ 群。

\subsection{投影对称群与无能隙费米子}

\begin{keypoints}
\item 有时，一个 PSG 保证无能隙费米子的存在。在这种情形下，PSG 的普适性意味着无能隙费米子的普适性（或稳定性）。
\end{keypoints}

为了演示 PSG 与无能隙费米子之间的直接联系，本节我们将研究一个特定的自旋液体，其量子序由 Z2Azz13 刻画（Wen and Zee, 2002）。Z2Azz13 自旋液体中的自旋子谱由自旋子跳跃哈密顿量
\begin{equation*}
H=\sum_{ij}\psi_{\pmb{i}}^{\dagger}u_{ij}\psi_{\pmb{j}}
\end{equation*}
给出。Z2Azz13 拟设的一个例子由式 (9.6.2) 给出。我们注意到，拟设 $u_{ij}$ 可以看作一个把费米子波函数 $\psi(\pmb{i})$ 映为 $\psi'(\pmb{i})=\sum_{\pmb{j}}u_{ij}\psi(\pmb{j})$ 的算符。这个算符记作 $\hat{H}$，将被称为哈密顿量。$\hat{H}$ 的本征值决定费米子谱。无能隙费米子对应于哈密顿量的零本征值。

Z2Azz13 的 PSG 由
\begin{align*}
G_{x}(\pmb{i}) &= \tau^{0}, & G_{y}(\pmb{i}) &= \tau^{0}\\
G_{P_{x}}(\pmb{i}) &= (-)^{\pmb{i}}\mathrm{i}\tau^{3} & G_{P_{y}}(\pmb{i}) &= (-)^{\pmb{i}}\mathrm{i}\tau^{3}\\
G_{P_{xy}}(\pmb{i}) &= \mathrm{i}\tau^{1} & G_{T}(\pmb{i}) &= \mathrm{i}\tau^{3}\\
G_{0}(\pmb{i}) &= -\tau^{0},\tag{9.10.1}
\end{align*}
生成。复合变换（如 $G_{P_{x}}P_{x}$）也可以看作作用在 $\psi_{\pmb{i}}$ 上的幺正算符。哈密顿量的投影对称性意味着，例如 $G_{P_{x}}P_{x}\hat{H}(G_{P_{x}}P_{x})^{\dagger}=\hat{H}$。

由于 $G_{x}=G_{y}=\tau^{0}$，Z2Azz13 拟设是平移不变的。在动量空间中，哈密顿量具有形式
\begin{equation*}
H(\pmb{k})=\epsilon^{\mu}(\pmb{k})\tau^{\mu}
\end{equation*}
拟设 $u_{ij}$ 在 $G_{T}T$ 下的不变性蕴含 $G_{T}(\pmb{i})u_{ij}G_{T}^{\dagger}(\pmb{j})=-u_{ij}$。用 $\pmb{k}$ 空间的语言表达，我们有
\begin{equation*}
U_{T}H(k_{x},k_{y})U_{T}^{\dagger}=-H(k_{x},k_{y}),\qquad U_{T}=\mathrm{i}\tau^{3}.\tag{9.10.2}
\end{equation*}
其中非平凡的 $G_{T}$ 给出非平凡的 $U_{T}$。式 (9.10.2) 蕴含 $\epsilon^{0,3}(\pmb{k})=0$，以及
\begin{equation*}
H(\pmb{k})=\epsilon^{1}(\pmb{k})\tau^{1}+\epsilon^{2}(\pmb{k})\tau^{2}\tag{9.10.3}
\end{equation*}
拟设在 $G_{P_{xy}}P_{xy}$ 下的不变性给出
\begin{equation*}
U_{P_{xy}}H(k_{x},k_{y})U_{P_{xy}}^{\dagger}=H(k_{y},k_{x}),\qquad U_{P_{xy}}=\mathrm{i}\tau^{1},\tag{9.10.4}
\end{equation*}
其中 $k_{x}$ 与 $k_{y}$ 的交换由 $P_{xy}$ 生成。拟设在 $G_{P_{x}}P_{x}$ 下的不变性导致
\begin{equation*}
U_{P_{x}}H(k_{x},k_{y})U_{P_{x}}^{\dagger}=H(-k_{x}+\pi,\,k_{y}+\pi),\qquad U_{P_{x}}=\mathrm{i}\tau^{3},\tag{9.10.5}
\end{equation*}
\begin{figure}[H]
\centering
\includegraphics[width=0.72\textwidth]{fig_9.16.pdf}
\caption*{图 9.16\quad $H(\pmb{k})$ 零点的两种图案。$\epsilon^{1}=0$ 线与 $\epsilon^{2}=0$ 线分别用实线和虚线表示。实心圆是缠绕数为 1 的零点，空心圆是缠绕数为 $-1$ 的零点。阴影区域是布里渊区的四分之一。当 $\epsilon^{2}$ 的零线移动而穿过 $\epsilon^{1}$ 的零线时，一个缠绕数为 1 的零点变成一个缠绕数为 $-1$ 的零点加上两个缠绕数为 1 的零点。}
\end{figure}
而拟设在 $G_{P_{y}}P_{y}$ 下的不变性导致
\begin{equation*}
U_{P_{y}}H(k_{x},k_{y})U_{P_{y}}^{\dagger}=H(k_{x}+\pi,\,-k_{y}+\pi),\qquad U_{P_{y}}=\mathrm{i}\tau^{3}.\tag{9.10.6}
\end{equation*}
（译注：原书引导句误印作 $G_{P_{y}}T_{P_{y}}$，式 (9.10.6) 右端误印作 $U_{P_{x}}$，均按上下文订正为 $G_{P_{y}}P_{y}$ 与 $U_{P_{y}}$。）动量平移 $(\pi,\pi)$ 来自 $G_{P_{x}}$ 与 $G_{P_{y}}$ 中的 $(-)^{\pmb{i}}$ 因子。于是，非零的 $\epsilon^{1,2}(\pmb{k})$ 具有下列对称性：
\begin{align*}
\epsilon^{1}(k_{x},k_{y})&=-\epsilon^{1}(\pi-k_{x},\,\pi+k_{y})=-\epsilon^{1}(\pi+k_{x},\,\pi-k_{y})=\epsilon^{1}(k_{y},k_{x})\\
\epsilon^{2}(k_{x},k_{y})&=-\epsilon^{2}(\pi-k_{x},\,\pi+k_{y})=-\epsilon^{2}(\pi+k_{x},\,\pi-k_{y})=-\epsilon^{2}(k_{y},k_{x})\tag{9.10.7}
\end{align*}

事实上，式 (9.10.3) 与式 (9.10.7) 定义了 $\pmb{k}$ 空间中最一般的 Z2Azz13 拟设。

式 (9.10.7) 使我们能够由 $\epsilon^{1,2}(\pmb{k})$ 在布里渊区四分之一区域内的取值确定它们（见图 9.16）。在该四分之一布里渊区内，$\epsilon^{2}(\pmb{k})$ 在 $k_{x}$ 与 $k_{y}$ 互换下是反对称的。因此，当 $k_{x}=k_{y}$ 时 $\epsilon^{2}(\pmb{k})=0$。式 (9.10.7) 还蕴含 $\epsilon^{1}(\pmb{k})$ 在绕 $(0,\pi)$ 或 $(\pi,0)$ 的 $90^{\circ}$ 旋转下变号，即 $\epsilon^{1}(k_{x},k_{y}+\pi)=-\epsilon^{1}(k_{y},-k_{x}+\pi)$。因此，$\epsilon^{1}(\pmb{k})$ 必在 $(0,\pi)$ 与 $(\pi,0)$ 处为零。它在连接 $(0,\pi)$ 与 $(\pi,0)$ 的一条线上也必为零（见图 9.16）。结果，$\epsilon^{1}=0$ 线与 $\epsilon^{2}=0$ 线在四分之一布里渊区内至少相交一次。交点就是哈密顿量 $H(\pmb{k})$ 的零点。我们看到，无能隙自旋子至少出现在布里渊区中四个孤立的 $\pmb{k}$ 点上（见图 9.16）。拟设的 Z2Azz13 对称性直接导致无能隙自旋子。

$H(\pmb{k})$ 的零点的两类典型分布看上去可能就像图 9.16 中那样。我们看到，这些零点具有特殊的图案。为了理解图案中哪些特征是普适的，我们想研究当我们使拟设变形时零点的运动。在这样做之前，我们想指出，$H(\pmb{k})$ 的零点具有一种可以用缠绕数刻画的内部结构。当 $\pmb{k}$ 绕一个零点一周时，二维矢量 $(\epsilon^{1}(\pmb{k}),\epsilon^{2}(\pmb{k}))$ 就围绕 $(0,0)$ 画出一个回路。零点的缠绕数由该回路围绕 $(0,0)$ 缠绕的圈数给出。如果回路围绕 $(0,0)$ 逆时针缠绕，缠绕数为正；如果顺时针缠绕，则为负。典型的零点具有缠绕数 $\pm1$。比如说，一个缠绕数为 2 的零点，在我们微扰哈密顿量之后可以分裂成两个缠绕数为 1 的零点。

对于 Z2Azz13 拟设，$H(\pmb{k})$ 的零点由 $\epsilon^{1}(\pmb{k})$ 与 $\epsilon^{2}(\pmb{k})$ 的零线的交点给出。由于 $\epsilon^{0}(\pmb{k})=\epsilon^{3}(\pmb{k})=0$，所以当我们使拟设变形时，零点不能随意地出现或消失。零点只能成组地产生或湮没，并保持总缠绕数守恒（见图 9.16）。再结合拟设的对称性 (9.10.7)，我们发现下列性质对哈密顿量的任何不改变 PSG 的微小扰动都是稳健的：沿 $((0,0),(\pi,\pi))$ 连线的 $+1$ 与 $-1$ 零点的图案；以及 $(N_{+},N_{-})$，即三角形 $((0,0),(\pi,\pi),(0,\pi))$ \emph{内部} 的 $(+1,-1)$ 零点的数目（不包括边和角上的零点）。我们把这些性质定义为 Z2Azz13 自旋液体的零点图案（pattern of the zero's，POZ）。正如费米面拓扑的改变一样，POZ 的改变将导致基态能量出现奇异性，并标志一个相变。因此，POZ 也是刻画量子序的一个量子数。我们看到，自旋液体中的量子序并不能完全由 PSG 刻画。POZ 提供了一种额外的刻画。PSG 与 POZ 的结合为量子序提供了更完全的刻画。

\subsection*{问题 9.10.1}

考虑式 (9.6.1) 中 Z2A0013 态的一个拟设。你可以假定 $\gamma_{1,2}=\lambda=0$。

(a) 证明：改变 $a_{0}^{1}$ 将引起 POZ 的改变。证明 POZ 的改变会在平均场基态能量中引起奇异性。

(b) 求转变点处低能自旋子的谱。

\subsection*{问题 9.10.2}

(a) 求 Z2A003n PSG 的 $G_{x,y}$、$G_{P_{x},P_{y},P_{xy}}$ 与 $G_{T}$。

(b) 求 Z2A003n 态最一般的平均场哈密顿量 $H(\pmb{k})$。

(c) 证明：Z2A003n 态在 $\pmb{k}=(\pm\frac{\pi}{2},\pm\frac{\pi}{2})$ 处总是具有无能隙自旋子。

% ---------------- 新增术语（ch9_g） ----------------
% | rigid spin liquid | 刚性自旋液体 |
% | Bose spin liquid | 玻色自旋液体 |
% | Fermi spin liquid | 费米自旋液体 |
% | infra-red divergence | 红外发散 |
% | counter term | 抵消项 |
% | lattice scale | 格点尺度 |
% | spinon hopping Hamiltonian | 自旋子跳跃哈密顿量（hopping 系列词沿“跳跃”：跳跃问题/跳跃幅/跳跃算符） |
% | zero-line | 零线 |
% | pattern of the zero's (POZ) | 零点图案（POZ） |
% | winding number of a zero | 零点的缠绕数（缠绕数术语表已收） |
% | momentum shift | 动量平移 |
% | dynamical stability | 动力学稳定性 |
% | spatially-dependent (gauge transformation) | 依赖空间的（规范变换） |
% | constant gauge transformation | 常值规范变换 |
